\documentclass[10pt,a4paper]{article}
\usepackage[utf8]{inputenc}
\usepackage[T1]{fontenc}
\usepackage{lmodern}
\usepackage[english]{babel}
\usepackage[margin=2.0cm]{geometry}
\usepackage{amsmath,amssymb,mathtools,bm}
\usepackage{booktabs,array,tabularx}
\newcolumntype{L}[1]{>{\raggedright\arraybackslash}p{#1}}
\usepackage{graphicx}
\usepackage{float}
\usepackage[font=small]{caption}
\usepackage{xcolor}
\usepackage[hidelinks]{hyperref}
\usepackage{url}
\makeatletter\g@addto@macro{\UrlBreaks}{\do\_\do/}\makeatother
\numberwithin{equation}{section}
\setlength{\emergencystretch}{3em}
\newcommand{\note}[1]{\begin{center}\fcolorbox{red!70!black}{red!5}{\parbox{0.94\linewidth}{\small\textbf{Note.} #1}}\end{center}}
\newcommand{\koA}{1268}
\newcommand{\faA}{0.8}
\newcommand{\preA}{5}
\newcommand{\revA}{2}
\newcommand{\vmaxA}{0.0222}
\newcommand{\resA}{1.00}
\newcommand{\cmA}{244}
\newcommand{\gapMA}{2.7}
\newcommand{\gapSA}{12}
\newcommand{\tqA}{43}
\newcommand{\tcA}{0.23}
\newcommand{\kzfreeA}{0}
\newcommand{\pkMA}{1.4}
\newcommand{\pkVA}{4.8}
\newcommand{\zMA}{0.023}
\newcommand{\zVA}{0.073}
\newcommand{\zNA}{0.023}
\newcommand{\lfMA}{1.2}
\newcommand{\lfNA}{7.6}
\newcommand{\koP}{11824}
\newcommand{\faP}{55.5}
\newcommand{\preP}{14}
\newcommand{\revP}{2}
\newcommand{\vmaxP}{0.0034}
\newcommand{\resP}{0.99}
\newcommand{\cmP}{34}
\newcommand{\gapMP}{8.6}
\newcommand{\gapSP}{8}
\newcommand{\tqP}{6}
\newcommand{\tcP}{0.41}
\newcommand{\kzfreeP}{87}
\newcommand{\pkMP}{1.7}
\newcommand{\pkVP}{0.5}
\newcommand{\zMP}{0.027}
\newcommand{\zVP}{0.034}
\newcommand{\zNP}{0.014}
\newcommand{\lfMP}{4.0}
\newcommand{\lfNP}{9.7}
\newcommand{\measA}{9.9}
\newcommand{\measP}{5.5}
\newcommand{\peakMA}{183}
\newcommand{\peakMP}{24}
\newcommand{\peakDGA}{206}
\newcommand{\peakDGP}{27}
\newcommand{\olDA}{42}
\newcommand{\olMA}{$-$86}
\newcommand{\clEA}{91}
\newcommand{\clCA}{$-$973}
\newcommand{\clGA}{31}
\newcommand{\clGCA}{$-$726}
\newcommand{\clFA}{56}
\newcommand{\clFDA}{47}
\newcommand{\kivMA}{12.8}
\newcommand{\kivDA}{0.37}
\newcommand{\kpvMA}{147.6}
\newcommand{\kpvDA}{187.4}
\newcommand{\olDP}{13}
\newcommand{\olMP}{$-$52}
\newcommand{\clEP}{98}
\newcommand{\clCP}{$-$315}
\newcommand{\clGP}{65}
\newcommand{\clGCP}{$-$313}
\newcommand{\clFP}{45}
\newcommand{\clFDP}{37}
\newcommand{\kivMP}{60.0}
\newcommand{\kivDP}{1.28}
\newcommand{\kpvMP}{6.1}
\newcommand{\kpvDP}{53.5}
\newcommand{\fgKpA}{1.35}
\newcommand{\fgPkA}{$-$59.5}
\newcommand{\fgZA}{0.0231}
\newcommand{\fgBwA}{0.12}
\newcommand{\vibA}{3.0}
\newcommand{\kpDropA}{0.5}
\newcommand{\hiPkA}{$-$35.5}
\newcommand{\hiKpA}{22}
\newcommand{\biasA}{-0.03}
\newcommand{\zFourA}{0.037}
\newcommand{\pkFourA}{$-$63.5}
\newcommand{\fgKpP}{1.05}
\newcommand{\fgPkP}{$-$62.7}
\newcommand{\fgZP}{0.0272}
\newcommand{\fgBwP}{0.20}
\newcommand{\vibP}{1.6}
\newcommand{\kpDropP}{0.5}
\newcommand{\hiPkP}{$-$36.5}
\newcommand{\hiKpP}{21}
\newcommand{\biasP}{4.15}
\newcommand{\zFourP}{0.043}
\newcommand{\pkFourP}{$-$66.5}
\newcommand{\finKpA}{8.147}
\newcommand{\finKiA}{6.46}
\newcommand{\finBwA}{0.384}
\newcommand{\finTsA}{3.72}
\newcommand{\finOvA}{19}
\newcommand{\finModA}{0.87}
\newcommand{\finIA}{17.7}
\newcommand{\finErrA}{585}
\newcommand{\finOvDA}{2284}
\newcommand{\casBwA}{0.37}
\newcommand{\casTsA}{4.68}
\newcommand{\casErrA}{372}
\newcommand{\casOvDA}{43}
\newcommand{\casIA}{11.7}
\newcommand{\mModA}{0.94}
\newcommand{\mHzA}{2}
\newcommand{\mLA}{$-$17.0}
\newcommand{\stepVibA}{2}
\newcommand{\stepOvA}{68}
\newcommand{\stepIA}{32}
\newcommand{\sTwoA}{86}
\newcommand{\sThreeA}{0.294}
\newcommand{\sRaisedA}{19.4}
\newcommand{\biasDA}{0.00}
\newcommand{\finKpP}{1425}
\newcommand{\finKiP}{223}
\newcommand{\finBwP}{32.8}
\newcommand{\finTsP}{0.068}
\newcommand{\finOvP}{6}
\newcommand{\finModP}{0.59}
\newcommand{\finIP}{11.5}
\newcommand{\finErrP}{20}
\newcommand{\finOvDP}{14}
\newcommand{\casBwP}{1.04}
\newcommand{\casTsP}{1.62}
\newcommand{\casErrP}{355}
\newcommand{\casOvDP}{11}
\newcommand{\casIP}{5.2}
\newcommand{\mModP}{0.46}
\newcommand{\mHzP}{29}
\newcommand{\mLP}{17.3}
\newcommand{\stepVibP}{213}
\newcommand{\stepOvP}{39}
\newcommand{\stepIP}{1215}
\newcommand{\sTwoP}{$1.3\cdot10^{6}$}
\newcommand{\sThreeP}{$6.3\cdot10^{4}$}
\newcommand{\sRaisedP}{$2.5\cdot10^{6}$}
\newcommand{\biasDP}{0.01}
\newcommand{\nbKpA}{25.21}
\newcommand{\nbKiA}{0.491}
\newcommand{\nbQpA}{768}
\newcommand{\nbQiA}{0.246}
\newcommand{\nbQrA}{646}
\newcommand{\nbMmA}{0.78}
\newcommand{\nbKpP}{798.8}
\newcommand{\nbKiP}{29.4}
\newcommand{\nbQpP}{$6.4\cdot10^{5}$}
\newcommand{\nbQiP}{$1\cdot10^{3}$}
\newcommand{\nbQrP}{$7.6\cdot10^{5}$}
\newcommand{\nbMmP}{0.56}


\title{Tuning of the LQG-PI position controller of the LUCY mount\\ in the PL configuration of Gawronski}
\author{Generated by \texttt{build\_lqg\_tuning\_report.py}\\ \small\texttt{sim/gimbal\_multidof/report\_lqg\_tuning}}
\date{\today}
\begin{document}
\maketitle

\note{Working document, updated after each step. Status: the tuning in force is the one of Section~\ref{sec:nb}, with the no-bearing internal model and the Dahl plant as simulator, tested on the circle with disturbances and noise in Section~\ref{sec:circle}. Sections~\ref{sec:step0}--\ref{sec:compare} record the path, including the tuning with the Maxwell element, now superseded. All results are simulations. The H7 configuration still contains placeholder data, and the figures carry the label ``result from placeholder parameters''. This report replaces the tuning in \path{report_lqg_gawronski}, which did not follow the procedure of the book.}

\section{Purpose and roadmap}
The model of the mount is used as a test bench to evaluate a controller in the style of the Deep Space Network: LQG in the position loop and PI in the velocity loop (configuration PL in W.~Gawronski, \emph{Modeling and Control of Antennas and Telescopes}, Springer, 2008, Chapter~8). To keep academic consistency, the weights are tuned with the procedure of the book. The phases are:
\begin{enumerate}
\item tuning of the weights (this report);
\item tracking of the Attorrese circle with process noise and disturbances, compared with the present cascade;
\item tracking of real drone trajectories: reference from the estimators (EKF2, OpenVINS), pointing error computed on the TagSLAM ground truth (Section~\ref{sec:thesis}).
\end{enumerate}

\section{Decision log}\label{sec:decisions}
\begin{table}[H]\centering\footnotesize
\setlength{\tabcolsep}{4pt}\begin{tabular}{lL{6.3cm}L{6.3cm}L{2.5cm}}\toprule
\# & decision & reason & review \\ \midrule
D1 & Configuration PL: velocity and current loops stay in the drive; LQG in the master at 1\,kHz. & Requested architecture, Section~8.5.2 of the book. & -- \\
D2 & State $[\hat e_i,\hat e,\hat x_f]$, with $\hat e=r-\hat y$. & Requested; book (8.12). Figures 8.4 and 8.22c of the book use the measured $y$ instead. & -- \\
D3 & \texttt{control.dlqr}, \texttt{control.dlqe} of python-control 0.10.2 (no Slycot). & Requested; same definitions as MATLAB. & -- \\
D4 & \texttt{dlqe} on $[y,x_f]$; $\hat e_i$ computed in the controller. & $y_i$ is not observable from $y$; \texttt{dlqe} fails on the augmented model. & -- \\
D5 & Bearings in the synthesis model by harmonic linearisation, for both \texttt{dlqr} and \texttt{dlqe}. & A real-time LQG can only hold a linear model; the model without bearings misses the low frequency (Section~\ref{sec:step0}). & -- \\
D6 & Operating point: bearing amplitude of the frequency-response tests. & The only point validated by a measurement; the book identifies its plant from a test. & Phase 2: compare with the no-bearing design \\
D8 & \textbf{Maxwell element} $K_b(s)=k_1s/(s+a)$, that is, a spring-damper system: spring $k_1$ in series with damper $c=k_1/a$. One more state. No static spring ($k_0=0$). Azimuth $k_1=\koA$\,N\,m/rad, $c=\cmA$\,N\,m\,s/rad, $f_a=\faA$\,Hz; elevation $k_1=\koP$\,N\,m/rad, $c=\cmP$\,N\,m\,s/rad, $f_a=\faP$\,Hz. & A viscous $c_{eq}$ overestimates the damping of the modes. A static spring is not physical for this bearing, and in azimuth it was a local minimum of the fit. & Step~1: the limits of Section~\ref{sec:limits}; compare with the no-bearing internal model \\
D10 & \textbf{Priority to the resonances}: weight 3 on the error in the resonance bands of the fit, 1 elsewhere. & Requested: keep the knowledge of the resonances rather than the low frequency. Above weight 3 the low-frequency error grows and the peak moves away from the measured frequency, with no gain on the modes (Table~\ref{tab:weights}). & -- \\
D7 & No feedforward and no command preprocessor in the tuning. & Style of the DSN: Chapter~8 does not name them in the LQG of the DSN (Figures~8.1--8.4, Sections~8.2--8.4). They appear with LQG only in the LMT analysis (Section~8.5.4), and with PI in Section~7.5. Consistency with the theoretical formulation of Chapter~8. & -- \\
D9 & Circle of radius 10\,m, height 5\,m, centre at 100\,m, 0.2\,rad/s. & Same as the study of the Hilkert observer, so results compare. & -- \\
D11 & First guess $q_i=0.1$, $q_p=0.3$, $q_f=0$, $Q=V$. & The values of the book (Section~8.4.2, point~2); Section~8.2 only says ``small''. & -- \\
D12 & \textbf{Section~8.2 followed to the letter.} Closed-loop transfer functions and step response on the reduced Maxwell model only, $Q=V$ (8.26). $V$ becomes a parameter in the fine tuning. & User's decision: the reduced Maxwell model is the analogue of the reduced linear model of the book. & Fine tuning \\
D13 & \textbf{End of Step 1: all flexible weights small, $q_f=0.1$ on every flexible state,} with $q_i=0.1$, $q_p=0.3$. No mode raised. Full vector in Table~\ref{tab:s1vector}. & The peaks in $T_{ry}$ are not too high: near $-$60\,dB. The book raises a peak near 0\,dB and leaves one near $-$6\,dB (Figure~8.5). User's indication: keep them low and raise only the corresponding weight of a peak that is too high. & Step 4 of the book \\
D14 & ``Too high'' is judged against the peaks of Figure~8.5 of the book. & The book gives no threshold, and this is the only quantitative reference in it. Reading from a figure, so approximate. & When a pointing requirement exists \\
D15 & $V$ is the intensity of the continuous equation (8.8): \texttt{dlqe} receives $VT^2$ with $W=1$. & Correction of an error: with $V$ per sample the estimator is about 760 times faster than (8.8). See Section~\ref{sec:v}. & -- \\
D16 & The gains stop at the current limit of the doublets, on both axes. This closes A1. & Decision of the user: the doublets are the application case; on LUCY the disturbances are generic, not wind. Section~\ref{sec:wind} remains as the sensitivity to this choice. & -- \\ \bottomrule
\end{tabular}\end{table}

\note{D5 and D8 (Maxwell element in the internal model) are superseded by the decision of the user of Section~\ref{sec:nb}: the state is estimated with a model without bearings, and the plant simulated stays Dahl. They stay in this log as the path followed.}

\begin{table}[H]\centering\footnotesize
\setlength{\tabcolsep}{4pt}\caption{Arbitrary choices of Steps 2--4 and of the test of $V$. The book does not give these values.}\label{tab:arbitrary}
\begin{tabular}{lL{7.2cm}L{8.4cm}}\toprule
\# & choice & reason, and what depends on it \\ \midrule
A1 & Acceleration limit read as the current limit of 18\,A in the V2 doublets, one per amplitude (0.1--8$^\circ$), on the Dahl plant, no current limit in the simulation, peak of the current command. & Sections~8.1.7 and~7.2 speak of the limit at ``typical commands''. The reduced model has no current output. It sets $q_p$ in azimuth. \\
A2 & No derating of the limit. & Section~7.2 sets a PI at 75\,\% of the limit, Section~8.2 does not. Literal reading of Section~8.2. \\
A3 & Further stops on $q_p$ and $q_i$: $\min|1+L|\geq0.5$ on the reduced model, and stability on the full Maxwell model and on the full plant with bearings sliding and stuck. & Without them the rule of the fundamental gives $k_p\approx7400$\,1/s in elevation at 1\,kHz. The stuck plant sets $q_p$ and $q_i$ in elevation: it goes beyond the book, which evaluates on its model. \\
A4 & Scales of half decades, bisection on 5 steps. & Resolution of the tuning. \\
A5 & Bandwidth: first $-3$\,dB crossing of $T_{ry}$. Settling: 2\,\% of a 1$^\circ$ step on the reduced Maxwell model. & Definitions not given in Section~8.2. In azimuth $T_{ry}$ has a plateau near $-10$\,dB, so the bandwidth is sensitive to the definition. \\
A6 & ``Oscillations'': at least two extrema outside 2\,\% after the first crossing, with the step low-pass filtered at 20\,Hz. & Section~8.1.7 places them below the fundamental frequency. The filter keeps out the vibration of the modes, which belongs to Step~4. \\
A7 & $q_i$ also limited by current and stability, not only by oscillations and (7.9). & $q_i$ raises $k_p$ too. In both axes one of these stops it before oscillations or (7.9). \\
A8 & Step~4: threshold $-6$\,dB (D14), scale from $10^3$, stop at the first value below it. & In elevation $10^4$ gives $-5.8$\,dB, so $3\cdot10^4$ is selected. \\
A9 & Fine tuning: index (8.27) reduced to $f=t_s/t_{s0}-f_o/f_{o0}$, equal weights, relative to the result of Step~4; grid of 6 values of $q_p$ and up to 7 of $q_i$, with $q_i$ not below Step~3; adopted if $f<-0.05$. & Objective requested (settling, bandwidth), with disturbance terms left for later. The overshoot is not in the index, so it can grow. Without the lower limit on $q_i$, the index pushes $q_i$ to zero. \\
A10 & $q_p$ raised again after the fine tuning, with a coarse scale. & The rule of Step~2 after $q_i$ and $q_f$ have changed. \\
A11 & $V=\rho Q$ with a scalar $\rho$ from $10^{-4}$ to $10^3$. Adopted only if $f<-0.05$ and neither the median settling nor the median rms error on the doublets on the Dahl plant get worse. & The book gives no procedure for $V$. On the model $V$ does not change $T_{ry}$, so it can only be judged with $q_p$ raised again and on the Dahl plant. A separate $V_f$ was not tested. \\
A13 & Disturbances of Section~\ref{sec:circle}: torque at the bearing port; step 0.5\,N\,m at 15\,s; Gaussian torque in the band 0.1--50\,Hz, 0.1\,N\,m rms, one seed per axis, with zero mean and with a mean of 0.2\,N\,m as in A15; sines 0.2\,N\,m at 1, 10 and 56\,Hz. Noise on the position measurement only. & No disturbance spectrum of LUCY is available; the levels are of the order of the Coulomb torque. The results scale with them, and the ranking is what matters. Noise on the velocity feedback is the same for all controllers and is left out. \\
A14 & Wind gusts of Section~\ref{sec:wind}: Davenport spectrum (5.14) with mean speed 5\,m/s; torque at the bearing port, as model 2 of the book, with assumed rms of 0.1 and 0.5\,N\,m; no mean torque; 120\,s, two seeds. & These levels are not derived from LUCY, whose wind loads are not known; the ratios between controllers changed little between the two levels. \\
A15 & Choices of Section~\ref{sec:tp_setup}: position of LUCY, command held between estimates, alignment on the whole window, disturbance of 0.1\,N\,m rms; for the mean of 0.2\,N\,m requested by the user, the same value and sign on both axes, already held by the loop when the run starts. & They are not fixed by the request; the sensitivity to alignment and position is in Section~\ref{sec:tp_limits}. A static torque is already present before tracking starts, and a sudden one is the step of Section~\ref{sec:circle}. \\
A12 & Noise on the measurement assumed white, 1 and 10\,$\mu$rad, in Section~\ref{sec:compare}. & The measured position is constant at rest, so the data do not give a noise level. \\ \bottomrule
\end{tabular}\end{table}

\section{Test bench}
The model is the one of \path{report_dahl_resampled} (uniform grid in $\ln f$): modal structure with viscous modal damping, pure Dahl bearings, AKD2G drive and TBM(S)-7646-A motor, in closed loop on the V2 doublets. Table~\ref{tab:model} gives the main values. $K_{pp}$, $K_{pv}$, $K_{iv}$ are the effective gains estimated on the 1$^\circ$ doublet, not the drive file values. The present cascade uses all three; the LQG plant uses $K_{pv}$ and $K_{iv}$.

\begin{table}[H]\centering\footnotesize
\caption{Main parameters of the test bench. $K_{pp}$ [1/s], $K_{pv}$ [A\,s/rad], $K_{iv}$ [Hz].}\label{tab:model}
\begin{tabular}{lrrrrrr}\toprule
axis & $\sigma_0$ [N\,m/rad] & $T_c$ [N\,m] & $g_{K_t}$ & $K_{pp}$ & $K_{pv}$ & $K_{iv}$ \\ \midrule
azimuth & 1571 & 0.232 & 0.900 & 7.92 & 187.4 & 0.37 \\
elevation & 17837 & 0.407 & 1.051 & 8.02 & 53.5 & 1.28
\\ \bottomrule
\end{tabular}\end{table}

\section{Step 0: linearisation of the bearings}\label{sec:step0}
\subsection{Two ways of linearising the Dahl element}
With $\dot z=\dot\delta-\sigma_0|\dot\delta|z/T_c$ and $T_b=\sigma_0z$, a first-order Taylor expansion around a state $(\dot\delta_0,z_0)$ gives
\begin{equation}
\Delta\dot z=\Bigl(1-\frac{\sigma_0z_0}{T_c}\,\mathrm{sign}\,\dot\delta_0\Bigr)\Delta\dot\delta-\frac{\sigma_0|\dot\delta_0|}{T_c}\Delta z .
\end{equation}
It has two natural points:
\begin{itemize}
\item at rest, $\dot\delta_0=z_0=0$: the bearing is a spring $\sigma_0$ (``stuck'');
\item in steady sliding, $\sigma_0z_0=T_c\,\mathrm{sign}\,\dot\delta_0$: the coefficient of $\Delta\dot\delta$ vanishes, and the friction is a constant torque with no linear dynamics (``no bearings'', with $T_c$ as a disturbance).
\end{itemize}
The small oscillations of the frequency-response tests lie between these two. They are described by the first harmonic, $K_b(A)$, which depends on the amplitude $A$ and not on the frequency. Its real part is an equivalent stiffness and its imaginary part an equivalent dissipation, constant with frequency. This is not a state-space element, and it has to be approximated.

\paragraph{Viscous approximation.} $K_b\simeq k_{eq}+j\omega c_{eq}$. The dissipation grows with frequency without limit, while in the describing function it does not. A value of $c_{eq}$ chosen at 1--10\,Hz therefore overestimates the damping of the modes at 56 and 98\,Hz.

\paragraph{Maxwell approximation (selected).} \textbf{The Maxwell element is nothing more than a spring-damper system}: a spring $k_1$ in series with a linear damper $c$. It adds no friction law, only two ordinary mechanical elements. In series both carry the same torque, and their deflections add up. With $z$ the deflection of the spring,
\begin{equation}
T_b=k_1z=c\,(\dot\delta-\dot z)\quad\Longrightarrow\quad\dot z=\dot\delta-az,\qquad a=\frac{k_1}{c},\qquad K_b(s)=\frac{T_b}{\delta}=\frac{k_1s}{s+a},
\end{equation}
with one more state, $z$. The two elements take turns with frequency:
\begin{itemize}
\item well above $f_a=a/2\pi$ the damper is stiff, the spring deflects, and $K_b\to k_1$: a \emph{spring}, with little dissipation at the modes;
\item well below $f_a$ the spring is stiff, the damper moves, and $K_b\to cs$: a \emph{viscous damper}. There is no static restoring torque, as in the Dahl element.
\end{itemize}
The dissipation $\mathrm{Im}\,K_b=k_1a\omega/(\omega^2+a^2)$ has its largest value, $k_1/2$, at $f_a$, and then falls, which the viscous element cannot do. In the Dahl element the dissipation does not fall to zero at low frequency, and a viscous term does not saturate at $T_c$: these differences are the price of a linear model (Section~\ref{sec:limits}).

\paragraph{Fit with priority to the resonances.} The two parameters are chosen so that the plant $y/u$ follows the plant with the describing function at the amplitude of the test, over 1--150\,Hz, with error on $\ln(y/u)$. The reference is the model, not the measurement, so the measurement remains an independent check. The error has weight 3 in the bands of the resonances (azimuth 48--62 and 62--72\,Hz; elevation 48--64 and 90--110\,Hz) and 1 elsewhere. Table~\ref{tab:weights} shows the effect of the weight. In azimuth it is almost nil: $f_a$ is far below the modes, where the element is only a spring, and the 56\,Hz damping stays that of the structure. In elevation $f_a$ falls next to the mode, so the weight matters: from 1 to 3 the peak moves 1.2\,dB towards the measurement, at the same low-frequency error; above 3 the peak gains at most 0.4\,dB, then returns, its frequency moves further from the measured one, and the low-frequency error grows. Weight 3 is selected (D10). The priority is therefore moderate: a stronger variant, with weight 0.03 outside the bands and $k_0$ free, returned an azimuth stiffness of 1077\,N\,m/rad, which is the static spring that this choice removes.

\begin{table}[H]\centering\footnotesize\setlength{\tabcolsep}{4pt}
\caption{Effect of the weight on the resonance bands, with $k_0=0$. Error against the plant from the measured frequency response. Peak of $|y/u|$ and damping of the mode near 56\,Hz.}\label{tab:weights}
\begin{tabular}{lrrrrrrrrr}\toprule
axis & weight & $k_1$ [N\,m/rad] & $f_a$ [Hz] & $c$ [N\,m\,s/rad] & \multicolumn{2}{c}{1--5\,Hz [dB, deg]} & peak [Hz] & peak [dB\,s] & $\zeta$ \\ \midrule
azimuth & 1 & 1264 & 0.82 & 244 & 1.2 & 8 & 55.7 & $-$62.4 & 0.0231 \\
 & \textbf{3} & \textbf{1268} & \textbf{0.83} & \textbf{244} & \textbf{1.2} & \textbf{8} & \textbf{55.7} & \textbf{$-$62.4} & \textbf{0.0231} \\
 & 10 & 1301 & 0.87 & 239 & 1.1 & 9 & 55.7 & $-$62.4 & 0.0231 \\
 & 30 & 1355 & 1.04 & 207 & 1.0 & 11 & 55.7 & $-$62.4 & 0.0231 \\
\emph{measured} &  &  &  &  &  & 55.7 & $-$63.7 &  \\
elevation & 1 & 9000 & 45.79 & 31 & 3.9 & 14 & 56.8 & $-$61.9 & 0.0237 \\
 & \textbf{3} & \textbf{11824} & \textbf{55.51} & \textbf{34} & \textbf{4.0} & \textbf{15} & \textbf{56.8} & \textbf{$-$63.1} & \textbf{0.0272} \\
 & 10 & 13129 & 45.43 & 46 & 4.6 & 17 & 57.4 & $-$63.5 & 0.0284 \\
 & 30 & 13396 & 38.25 & 56 & 5.3 & 19 & 57.4 & $-$63.2 & 0.0279 \\
\emph{measured} &  &  &  &  &  & 56.3 & $-$64.8 & 
\\ \bottomrule
\end{tabular}\end{table}

The static spring $k_0$ of the earlier version of this report (azimuth 507, elevation 89\,N\,m/rad) is removed. With $k_0$ free the fit now gives \kzfreeA\,N\,m/rad in azimuth, so the earlier value was a local minimum, and \kzfreeP\,N\,m/rad in elevation, with a gain of only 0.3\,dB in 1--5\,Hz (Table~\ref{tab:cand}). The result is a linear model that a real-time controller can hold. The operating point remains an amplitude, not a state.

\subsection{Candidates against the measurement}
Table~\ref{tab:cand} and Figure~\ref{fig:plant} compare each candidate with the plant from velocity command to encoder obtained by closing the velocity loop of the model on the measured frequency response.

\begin{table}[H]\centering\footnotesize\setlength{\tabcolsep}{3.5pt}
\caption{Rms error against the plant from the measured frequency response.}\label{tab:cand}
\begin{tabular}{L{6.3cm}rrrrrrrr}\toprule
& \multicolumn{4}{c}{azimuth} & \multicolumn{4}{c}{elevation} \\
\cmidrule(lr){2-5}\cmidrule(lr){6-9}
& \multicolumn{2}{c}{1--5\,Hz} & \multicolumn{2}{c}{5--30\,Hz} & \multicolumn{2}{c}{1--5\,Hz} & \multicolumn{2}{c}{5--30\,Hz} \\
bearings in the model & dB & deg & dB & deg & dB & deg & dB & deg \\ \midrule
no bearings (Taylor, steady sliding) & 7.6 & 54 & 4.1 & 30 & 9.7 & 14 & 3.0 & 28 \\
stuck, $\sigma_0$ (Taylor, at rest) & 1.8 & 28 & 2.0 & 23 & 28.8 & 94 & 12.6 & 116 \\
describing function, test amplitude (not state space) & 0.5 & 3 & 1.7 & 22 & 0.4 & 8 & 1.5 & 9 \\
viscous, $k_{eq}+c_{eq}$ & 1.5 & 7 & 2.5 & 22 & 3.8 & 10 & 2.6 & 13 \\
\textbf{Maxwell, $k_1s/(s+a)$} & 1.2 & 8 & 2.1 & 22 & 4.0 & 15 & 2.4 & 12 \\
Maxwell with free $k_0$ (not adopted) & 1.2 & 8 & 2.1 & 22 & 3.7 & 13 & 2.3 & 12 \\
two Maxwell branches (not adopted) & 1.3 & 6 & 1.7 & 20 & 3.2 & 14 & 2.3 & 13 \\
\textbf{reduced synthesis model} & 1.2 & 8 & 2.1 & 22 & 4.0 & 15 & 2.4 & 12 \\
describing function, $A=10\,\mu$rad & 1.6 & 27 & 2.0 & 22 & 28.4 & 84 & 11.6 & 102 \\
describing function, $A=100\,\mu$rad & 0.6 & 13 & 1.9 & 20 & 17.9 & 41 & 4.7 & 35 \\
describing function, $A=1000\,\mu$rad & 4.6 & 39 & 3.7 & 27 & 1.8 & 12 & 2.6 & 22 \\
describing function, $A=10000\,\mu$rad & 7.1 & 53 & 4.1 & 29 & 7.8 & 12 & 2.9 & 28
\\ \bottomrule
\end{tabular}\end{table}

\begin{figure}[H]\centering
\includegraphics[width=0.95\linewidth]{figures/plant_operating_points.pdf}
\caption{Velocity-loop plant with the bearings at different operating points.}\label{fig:plant}
\end{figure}

\begin{table}[H]\centering\footnotesize
\caption{Dissipation at the modes. Height of the peak of $|y/u|$ near the structural modes [dB\,s].}\label{tab:peaks}
\begin{tabular}{lrrrr}\toprule
& \multicolumn{2}{c}{azimuth} & \multicolumn{2}{c}{elevation} \\
\cmidrule(lr){2-3}\cmidrule(lr){4-5}
& 50--62\,Hz & 90--110\,Hz & 50--62\,Hz & 90--110\,Hz \\ \midrule
no bearings & -62.2 & -84.7 & -57.8 & -79.8 \\
stuck, $\sigma_0$ & -62.5 & -84.7 & -56.4 & -77.8 \\
describing function, test amplitude & -63.0 & -84.7 & -63.2 & -79.0 \\
viscous, $k_{eq}+c_{eq}$ & -68.5 & -84.7 & -65.3 & -82.1 \\
\textbf{Maxwell} & -62.4 & -84.7 & -63.1 & -79.7 \\
measured & -63.7 & -86.4 & -64.8 & -78.9
\\ \bottomrule
\end{tabular}\end{table}

\begin{table}[H]\centering\footnotesize
\caption{Rms error in the resonance bands of the fit against the plant from the measured frequency response.}\label{tab:resbands}
\begin{tabular}{lrrrr}\toprule
& \multicolumn{2}{c}{azimuth} & \multicolumn{2}{c}{elevation} \\
\cmidrule(lr){2-3}\cmidrule(lr){4-5}
bearings in the model & dB & deg & dB & deg \\ \midrule
no bearings & 3.3 & 33 & 4.8 & 17 \\
viscous, $k_{eq}+c_{eq}$ & 4.1 & 44 & 4.2 & 17 \\
describing function (not state space) & 2.7 & 31 & 2.7 & 16 \\
\textbf{Maxwell} & 2.8 & 31 & 3.4 & 14 \\
\textbf{reduced synthesis model} & 2.8 & 31 & 3.4 & 14
\\ \bottomrule
\end{tabular}\end{table}

\begin{table}[H]\centering\footnotesize
\caption{Damping of the lightly damped poles of the velocity loop near the modes. In azimuth the 65\,Hz pole belongs to the resonator of the drive.}\label{tab:damp}
\begin{tabular}{lrrrr}\toprule
axis & $f$ [Hz] & no bearings & viscous & Maxwell \\ \midrule
azimuth & 56.0 & 0.023 & 0.073 & \textbf{0.023} \\
azimuth & 65.1 & 0.028 & 0.028 & \textbf{0.028} \\
elevation & 56.1 & 0.014 & 0.034 & \textbf{0.027} \\
elevation & 98.3 & 0.031 & 0.043 & \textbf{0.036}
\\ \bottomrule
\end{tabular}\end{table}

\begin{table}[H]\centering\footnotesize\setlength{\tabcolsep}{4pt}
\caption{Selected model. Amplitudes: range (median) over 1--10\,Hz in the test. $c=k_1/a$ and $c_{eq}$ in N\,m\,s/rad: damper of the Maxwell element, which acts only well below $f_a$, and viscous coefficient. Residue: of the pole at the origin, 1 without stiffness. Order: of the reduced synthesis model.}\label{tab:point}
\begin{tabular}{lrlrrrrrr}\toprule
axis & $T_c/\sigma_0$ [$\mu$rad] & amplitude [$\mu$rad] & $k_1$ [N\,m/rad] & $f_a$ [Hz] & $c$ & $c_{eq}$ & residue & order \\ \midrule
azimuth & 147 & 69--443 (185) & 1268 & 0.8 & 243.8 & 26.5 & 1.00 & 10 \\
elevation & 23 & 143--3282 (1155) & 11824 & 55.5 & 33.9 & 25.8 & 0.99 & 12
\\ \bottomrule
\end{tabular}\end{table}

\subsection{Where the bearings are during tracking}
Figure~\ref{fig:circle} shows the bearing state on the Attorrese circle (radius 10\,m, height 5\,m, centre at 100\,m, 0.2\,rad/s), with the cascade and without noise or disturbance. After the first 5\,s from rest, the bearings are in pre-sliding ($|\sigma_0z/T_c|<0.95$) for \preA\,\% of the orbit in azimuth, with \revA\ velocity reversals, and for \preP\,\% in elevation, with \revP\ reversals. Near each reversal the axis stops and restarts: in elevation it stays still for about 2\,s and then jumps, a stick-slip that no linear model describes. Largest velocity: \vmaxA\ and \vmaxP\,rad/s.

\begin{figure}[H]\centering
\includegraphics[width=0.95\linewidth]{figures/circle_bearing_state.pdf}
\caption{Bearing velocity and normalised state on the circle; dotted lines: limit of pre-sliding used above.}\label{fig:circle}
\end{figure}

\subsection{Choice and consequences}
The Maxwell element, a spring-damper system, is selected at the point of the frequency-response tests, with priority to the resonances (D6, D8, D10). It reaches about the same low-frequency error as the viscous element without overestimating the damping. The damping at 56\,Hz is \zMA\ in azimuth and \zMP\ in elevation, against \zVA\ and \zVP\ with the viscous element and \zNA\ and \zNP\ without bearings (Tables~\ref{tab:peaks} and~\ref{tab:damp}). A second branch was also fitted: it gains at most 0.8\,dB in 1--5\,Hz, not enough for one more state. Four facts must be kept in view.
\begin{itemize}
\item \textbf{The low frequency is improved, not solved.} At 1\,Hz the model is still \gapMA\,dB above the measurement in azimuth and \gapMP\,dB in elevation (Figure~\ref{fig:plant}), and the rms error in 1--5\,Hz is \lfMA\ and \lfMP\,dB, against \lfNA\ and \lfNP\,dB without bearings. This is the cost accepted in order to keep the resonances.
\item \textbf{The resonance peaks err on the side of less damping.} The peak at 56\,Hz is \pkMA\,dB above the measurement in azimuth and \pkMP\,dB in elevation, while the viscous element is \pkVA\,dB below in azimuth and \pkVP\,dB below in elevation. This is the side requested: damping is not overestimated. In the resonance bands the error stays near that of the describing function (Table~\ref{tab:resbands}): the residual of about 2.7\,dB and 31$^\circ$ in azimuth is also there with the describing function, so it does not come from the Maxwell approximation.
\item \textbf{It is not the most frequent point on the noise-free circle.} There the bearings slide for most of the orbit, which is the no-bearing point. The choice rests on the validation against a measurement and on the book, which identifies its plant from a test. The low-frequency disturbance regime, near velocity reversals and on slow real trajectories, is closer to this point. Phase~2 compares the design with the no-bearing design, with noise and disturbances, and this decision is reviewed there.
\item \textbf{The synthesis plant is not the plant in sliding.} Without the static spring the residue of the pole at the origin is \resA\ in azimuth and \resP\ in elevation, so the very low frequency agrees with the plant in sliding. The gap remains around the corner of the element: at 1\,Hz the plant in sliding has \gapSA\,dB more gain than the design model in azimuth and \gapSP\,dB in elevation. The margins must be checked with bearings sliding and stuck, as well as at the design point.
\end{itemize}

\paragraph{Synthesis model.} The model with the Maxwell element goes through the modal form, truncation, position as first state and integral augmentation of Section~8.1 without changes, and \texttt{dlqr} and \texttt{dlqe} solve on it with stable poles (test weights $q_p=100$, $q_i=1$). The reduced model reproduces the full one: the errors in Table~\ref{tab:cand} are nearly the same. Implemented in \path{closedLoopModel/lucy_cl/lqg.py} (\texttt{con\_cuscinetti}), with a unit test, and selected through the argument \texttt{cuscinetti} of \texttt{costruisci} in \texttt{analisi\_lqg\_gawronski.py}. The previous results in \path{results_lqg_gawronski} were not regenerated.

\begin{table}[H]\centering\footnotesize
\caption{Poles of the reduced synthesis model, with the rule $f\leq150$\,Hz, $\|G_i\|_\infty\geq10^{-5}$: frequency in Hz, damping in brackets if below 1. The pole at the origin is also kept.}\label{tab:poles}
\begin{tabular}{lL{13cm}}\toprule
axis & poles \\ \midrule
azimuth & 0.05; 4.20 (0.469); 13.32 (0.204); 56.36 (0.023); 65.10 (0.028) \\
elevation & 0.40; 10.10; 44.94; 49.37 (0.169); 56.90 (0.027); 96.44 (0.172); 98.95 (0.036)
\\ \bottomrule
\end{tabular}\end{table}

\subsection{Time-domain check of the Maxwell element}\label{sec:td}
To see whether the linear element keeps the correlation with the doublets, the V2 doublets are run with the Maxwell element replacing the Dahl element in the plant, in open loop (measured current in) and in closed loop (measured position reference in), with the gains of the drive files and with gains estimated on the 1$^\circ$ doublet as in the other cases. The model is linear and has no limits, so peaks above 18\,A are not physical. Dahl is the reference; its estimated-gain results reproduce those of \path{report_dahl_resampled} exactly.

\begin{table}[H]\centering\footnotesize\setlength{\tabcolsep}{4pt}
\caption{Median over the moves, without the first move of the 8$^\circ$ azimuth doublet. Position and current fit [\%] during the moves; largest model current [A]. Measured median peak current: \measA\,A in azimuth and \measP\,A in elevation.}\label{tab:td}
\begin{tabular}{llrrrrrrr}\toprule
& & open loop & \multicolumn{3}{c}{closed loop, estimated gains} & \multicolumn{3}{c}{closed loop, file gains} \\
\cmidrule(lr){4-6}\cmidrule(lr){7-9}
axis & bearings & position & position & current & $i_\mathrm{max}$ & position & current & $i_\mathrm{max}$ \\ \midrule
azimuth & Dahl & 42 & 96 & 20 & 11.6 & 47 & $-$287 & 42.4 \\
 & \textbf{Maxwell} & $-$86 & 91 & $-$973 & 75.0 & 56 & $-$1061 & 100.7 \\
elevation & Dahl & 13 & 96 & 50 & 5.1 & 37 & 24 & 4.1 \\
 & \textbf{Maxwell} & $-$52 & 98 & $-$315 & 18.8 & 45 & $-$232 & 8.7
\\ \bottomrule
\end{tabular}\end{table}

\begin{table}[H]\centering\footnotesize
\caption{Gains. $K_{pp}$ [1/s], $K_{pv}$ [A\,s/rad], $K_{iv}$ [Hz].}\label{tab:td_gains}
\begin{tabular}{llrrr}\toprule
axis & source & $K_{pp}$ & $K_{pv}$ & $K_{iv}$ \\ \midrule
azimuth & files (.akd, WorkBench) & 62.8 & 250.0 & 1.00 \\
 & estimated, Dahl & 7.9 & 187.4 & 0.37 \\
 & \textbf{estimated, Maxwell} & 8.6 & 147.6 & 12.76 \\
elevation & files (.akd, WorkBench) & 2.5 & 25.0 & 15.00 \\
 & estimated, Dahl & 8.0 & 53.5 & 1.28 \\
 & \textbf{estimated, Maxwell} & 7.4 & 6.1 & 60.00
\\ \bottomrule
\end{tabular}\end{table}

\begin{itemize}
\item \textbf{Correlation is lost, also without the static spring.} In open loop the median position fit falls from \olDA\ to \olMA\,\% in azimuth and from \olDP\ to \olMP\,\% in elevation. In closed loop with estimated gains the current fit is \clCA\,\% in azimuth and \clCP\,\% in elevation, with peaks up to \peakMA\,A and \peakMP\,A against the limit of 18\,A.
\item \textbf{The good closed-loop position fit (\clEA\ and \clEP\,\%) is misleading.} The controller imposes the position, and the current needed to do it is not limited in the linear model. The estimated gains also absorb the error and lose their meaning: $K_{iv}$ goes from \kivDA\ to \kivMA\,Hz in azimuth, and in elevation it reaches the upper bound of the estimate, 60\,Hz, while $K_{pv}$ falls from \kpvDP\ to \kpvMP\,A\,s/rad. Those gains must not be used. With the gains estimated for Dahl, without a new estimate, the median position fit of the Maxwell element falls to \clGA\,\% in azimuth and \clGP\,\% in elevation, and the current fit is \clGCA\ and \clGCP\,\%, with peaks of \peakDGA\ and \peakDGP\,A. With the file gains the position fit is \clFA\ and \clFP\,\%, against \clFDA\ and \clFDP\,\% with Dahl: similar, because both fit poorly.
\item \textbf{The cause is now the viscous term.} Removing $k_0$ removes the current that grew with the stroke, but during a move the element is the damper $c$. At 10$^\circ$/s it asks for about \tqA\,N\,m in azimuth and \tqP\,N\,m in elevation, against a true friction $T_c$ of \tcA\ and \tcP\,N\,m. The Dahl element saturates at $T_c$; a linear element cannot.
\end{itemize}

\paragraph{What this means.} The test confirms that a linear element is valid only near its operating point, as expected, and does not replace Dahl as a simulation model. In the phases that follow the plant remains Dahl, and the Maxwell element stays in the internal model of the controller, for small signals.

\subsection{Limits of the choice, to check in Step~1}\label{sec:limits}
These are risks identified from the model. The bias and the limit cycles have a first check in Section~\ref{sec:step1}; the others are not yet measured.
\begin{itemize}
\item \textbf{Prediction during large motion.} During tracking and moves the internal model sees a large viscous torque that the plant does not have, so the estimator has to correct the prediction continuously.
\item \textbf{Bias of the estimate.} The integrator acts on $\hat e$, so it brings $\hat y$, not $y$, to $r$. A persistent model error, with no disturbance state in the estimator, can leave a bias $y-\hat y$. The error must be measured on the true $y$. If the bias is significant, the alternatives are an error on the measured $y$ (Figures~8.4 and~8.22c of the book) or a disturbance state.
\item \textbf{The state $z$.} It has no measurement and no physical meaning in the plant, which holds the Dahl state. It enters $K_f$ through the modal coordinates, and the book gives no rule for the weight of such a state. Its weight is chosen and reported in Step~1.
\item \textbf{Margins.} To be checked with the plant in sliding (\gapSA\ and \gapSP\,dB of extra gain at 1\,Hz) and stuck, not only at the design point.
\item \textbf{Stick-slip and limit cycles.} The integrator together with the Coulomb level of the plant can produce hunting near reversals, which no linear analysis shows. To be checked in simulation on the Dahl plant.
\item \textbf{Elevation: $f_a$ next to the 56\,Hz mode.} There the element sets part of the damping of the mode (\zMP\ against \zNP\ without bearings). The true Dahl dissipation changes with amplitude, so the estimate at that mode may be sensitive to it.
\end{itemize}
The tuning uses the Maxwell internal model (D12). The comparison with the no-bearing design is left to Phase~2.

\begin{figure}[H]\centering\includegraphics[width=\linewidth]{figures/td_azimuth_open_loop.pdf}\caption{Azimuth, open loop.}\end{figure}
\begin{figure}[H]\centering\includegraphics[width=\linewidth]{figures/td_elevation_open_loop.pdf}\caption{Elevation, open loop.}\end{figure}
\begin{figure}[H]\centering\includegraphics[width=\linewidth]{figures/td_azimuth_closed_loop_estimated.pdf}\caption{Azimuth, closed loop, estimated gains.}\end{figure}
\begin{figure}[H]\centering\includegraphics[width=\linewidth]{figures/td_elevation_closed_loop_estimated.pdf}\caption{Elevation, closed loop, estimated gains.}\end{figure}
\begin{figure}[H]\centering\includegraphics[width=\linewidth]{figures/td_azimuth_closed_loop_akd.pdf}\caption{Azimuth, closed loop, file gains.}\end{figure}
\begin{figure}[H]\centering\includegraphics[width=\linewidth]{figures/td_elevation_closed_loop_akd.pdf}\caption{Elevation, closed loop, file gains.}\end{figure}

\section{The tuning parameters in the formulation of Gawronski}\label{sec:weights}
\subsection{Performance index and weights}
The plant is the velocity loop, $\dot x=Ax+Bu+v$, $y=Cx+w$, where $v$ and $w$ are process and measurement noise with covariances $V$ and $W=I$ (Section~8.1.1). The regulator minimises
\begin{equation}
J^2=E\!\left(\int_0^\infty\bigl(x^TQx+u^TRu\bigr)\,dt\right),\qquad R=1,\tag{8.3}
\end{equation}
with $u=-K_c\hat x$, $K_c=B^TS_c$ (8.4, 8.5), and $S_c$ solving the controller Riccati equation
\begin{equation}
A^TS_c+S_cA-S_cBB^TS_c+Q=0.\tag{8.6}
\end{equation}
The estimator $\dot{\hat x}=A\hat x+Bu+K_e(y-C\hat x)$ (8.2) has $K_e=S_eC^T$ (8.7), with $S_e$ solving the dual equation
\begin{equation}
AS_e+S_eA^T-S_eC^TCS_e+V=0.\tag{8.8}
\end{equation}
Since $A$, $B$, $C$ are fixed, $K_c$ depends only on $Q$ and $K_e$ only on $V$: the weights are the only design freedom. With the integral augmentation the tracking state and the gain are divided as
\begin{equation}
\hat x=\begin{Bmatrix}\hat e_i\\\hat e\\\hat x_f\end{Bmatrix},\qquad K_c=\begin{bmatrix}k_i&k_p&K_f\end{bmatrix},\qquad u=k_i\hat e_i+k_p\hat e-K_f\hat x_f.\tag{8.12, 8.14}
\end{equation}
\textbf{This law differs in sign from the one printed in the book.} Equation (8.15) reads
\begin{equation}
u=-k_ie_i-k_pe-K_f\hat x_f.\tag{8.15}
\end{equation}
It follows from (8.4), $u=-K_c\hat x$, with $\hat y_i$ and $\hat y$ replaced by $e_i$ and $e$. But the book defines $e=r-\hat y$, so for $r=0$ the replacement that keeps the feedback is $\hat y=-e$ and $\hat y_i=-e_i$, which gives $+k_i$ and $+k_p$, while $x_f$ is not affected and keeps $-K_f$. The law used here is the one implied by the closed-loop matrices of the book, (8.17) and (8.18), where $k_i$ and the command enter with the plus sign and $k_p$, $k_f$ multiply the output with the minus sign, and it matches Figure~8.22c. A rigid plant shows why: with $\dot y=u$ and $u=-k_p(r-y)$ the pole is at $+k_p$. On the reduced Maxwell model, at the first guess and at $q_p=100$, $q_i=25$, the signs of (8.15) with $e=r-\hat y$ give an unstable loop on both axes (largest $|z|$ between 1.0006 and 1.0043), and the law above a stable one. As printed, (8.15) holds only if $e$ is read as $\hat y-r$. As decided in D2, the error is the estimated one, $\hat e=r-\hat y$, as in Figure~8.3 of the book.

\subsection{The vector of parameters}
\textbf{The tuning looks for one set of parameters, collected in a column vector} that follows the state vector (8.23--8.25):
\begin{equation}
Q=\mathrm{diag}(q),\qquad
q=\begin{Bmatrix}q_i\\q_p\\q_f\end{Bmatrix},\qquad
q_f=\begin{Bmatrix}q_{f1m}\\q_{f2m}\\\vdots\end{Bmatrix},\qquad
q_{fim}=\begin{Bmatrix}q_{fi}\\\dot q_{fi}\end{Bmatrix}.
\end{equation}
Here $q_i$ weights the integral of the position, $q_p$ the position, and $q_{fim}$ the $i$th flexible mode. In $\dot q_{fi}$ the dot is only a name: it is the weight of the velocity state of the mode, not the derivative of a weight. The book then takes $V=Q=\mathrm{diag}(q)$ with two equal entries per mode (8.26):
\begin{equation}
q=\begin{bmatrix}q_{ei}&q_e&q_{f1}&q_{f1}&\cdots&q_{fn}&q_{fn}\end{bmatrix}.\tag{8.26}
\end{equation}
With $n$ modes $q$ has $2n+2$ entries but only $n+2$ independent values. The book's typical $n=4$ gives 10 entries and 6 values.

\subsection{The role of $V$}\label{sec:v}
$V=E(vv^T)$ is the intensity of the process noise, as the book defines it, and $W=E(ww^T)=I$ is that of the measurement noise (Section~8.1.1). It enters only the estimator, through (8.8): ``the covariance matrix specifies the noise in the system, and it impacts the estimator gains'' (p.~118). Its meaning is how much each state may be disturbed beyond what the model predicts, and so how much the estimator should trust the measurement rather than the model.
\begin{itemize}
\item \textbf{Size against $W$.} Larger $V/W$ means a faster estimator that follows the measurement, and more noise passes to the control. Smaller $V/W$ means a slower estimator that trusts the model. For the position state alone, the estimator pole is at $\sqrt{V/W}$.
\item \textbf{Distribution among the states.} $K_e$ divides the innovation $y-\hat y$ among the states in proportion to the $V$ of each, weighted by how much the state shows in $y$. A large $V$ on a mode tells the estimator that much of the innovation belongs to that mode.
\end{itemize}
In the book $V$ is not measured: with $Q=V$ it is used as a design parameter, an assumption that it calls ``opportunistic'' and that eases the search (pp.~118--119). What $Q=V$ gives is that (8.8) has the form of (8.6), so the estimator and the regulator have similar bandwidths: at the first guess $\sqrt{0.3}=0.55$\,rad/s against about 0.56\,rad/s for the rigid regulator poles.

On the model, $V$ does not change $T_{ry}$: the command enters plant and estimator alike and does not excite the estimation error, so $T_{ry}$ depends only on $Q$, and this was verified numerically on the reduced Maxwell model. $V$ acts on the disturbance response and on the sensitivity to model error.

\paragraph{Correction: $V$ in discrete time.} Equation (8.8) is continuous, while the implementation uses \texttt{dlqe}. Sampling at $T$ gives $V_d=VT$ and $W_d=W/T$, so at $W_d=1$ the weight is $VT^2$. Earlier versions of this report passed $V$ directly, as the box of p.~98 literally reads. For \texttt{dlqr} this makes no difference, since $Q$ and $R$ scale together. For \texttt{dlqe} it made the estimator about 760 times faster than (8.8): at the first guess 418\,rad/s against 0.55\,rad/s. With $VT^2$ the estimator poles coincide with those of (8.8), verified with \texttt{control.lqe} (and in a unit test). \textbf{This withdraws two conclusions of the earlier versions:} that $Q=V$ loses stability at low weights, and that the first guess ends in a limit cycle in elevation. Table~\ref{tab:s1v} shows the difference. The results in \path{report_lqg_gawronski} and the need for a small $\rho_e$ found there come from the same scale.

\begin{table}[H]\centering\footnotesize
\caption{The two discretisations of $V$, $Q=V$, weight on the 56\,Hz mode. $K_{e,y}$: estimator gain on the position [rad/s]. Largest pole on the full plant, bearings sliding and stuck; outside the procedure of the book.}\label{tab:s1v}
\begin{tabular}{lrlrll}\toprule
axis & $q_f$ & $V$ & $K_{e,y}$ & sliding & stuck \\ \midrule
azimuth & 0 & V per sample & 417.89 & stable & stable \\
 & 0 & V of (8.8) & 0.55 & stable & stable \\
 & $10^{4}$ & V per sample & 1838.11 & \textbf{unstable}, 0.7\,Hz & stable \\
 & $10^{4}$ & V of (8.8) & 10.34 & stable & stable \\
elevation & 0 & V per sample & 417.89 & stable & stable \\
 & 0 & V of (8.8) & 0.55 & stable & stable \\
 & $10^{4}$ & V per sample & 1890.76 & \textbf{unstable}, 1.3\,Hz & \textbf{unstable}, 21.2\,Hz \\
 & $10^{4}$ & V of (8.8) & 11.79 & stable & stable
\\ \bottomrule
\end{tabular}\end{table}

\subsection{What ``two-state mode'' means}
A structural mode is a second-order differential equation in the modal displacement $q_{mi}$ (2.20):
\begin{equation}
\ddot q_{mi}+2\zeta_i\omega_i\dot q_{mi}+\omega_i^2q_{mi}=b_{mi}u .
\end{equation}
A state-space model only holds first-order equations, so each mode needs two states, displacement and velocity. This is the ``two-state mode'' of the book, and it is the 2$\times$2 block of the modal state matrix (2.25, 2.27):
\begin{equation}
x_i=\begin{Bmatrix}\omega_iq_{mi}\\\dot q_{mi}\end{Bmatrix},\qquad
A_{mi}=\begin{bmatrix}0&\omega_i\\-\omega_i&-2\zeta_i\omega_i\end{bmatrix}.
\end{equation}
\textbf{A parameter appears also for the derivative because $Q$ has one entry per state}, and the velocity is a state like the displacement. The term is not redundant in meaning either. With equal weights the mode enters the index as
\begin{equation}
q_{fi}\,x_i^Tx_i=q_{fi}\bigl(\omega_i^2q_{mi}^2+\dot q_{mi}^2\bigr),
\end{equation}
which is twice the energy of the mode per unit modal mass, potential plus kinetic. During a vibration the energy passes from one state to the other every quarter of a period. With both weights the cost does not depend on the phase of the vibration, and for an undamped mode it stays constant in time; with the displacement weight alone it would vanish at each zero crossing. So the index weighs the vibration energy of the mode, and this is why the book gives the pair one value.

\paragraph{On LUCY.} The modal form in \path{closedLoopModel/lucy_cl/lqg.py} uses a different but equivalent block, $\bigl[\begin{smallmatrix}-\sigma&-\omega\\\omega&-\sigma\end{smallmatrix}\bigr]$, with the real and imaginary parts of the eigenvector as states and the scale $\|B_i\|=\|C_i\|$. It has the same poles. For small damping it is related to (2.27) by a rotation and a scale, and equal weights do not change under a rotation, so one value per mode keeps the same meaning. That scale is one reason why the values of $q_f$ are not comparable with those of the book. A real pole has a single state and a single weight. So $q$ has 11 entries in azimuth ($q_i$, $q_p$, the real pole at 0.05\,Hz, and the blocks at 4.2, 13.3, 56.4 and 65.1\,Hz), and 13 in elevation ($q_i$, $q_p$, the real poles at 0.4, 10.1 and 44.9\,Hz, and the blocks at 49.4, 56.9, 96.4 and 99.0\,Hz). Not all of these are structural modes: some are dynamics of the velocity loop, of the drive filters and of the Maxwell element.

\subsection{Why $q_i$, $q_p$ are tied to $k_i$, $k_p$}
The first two elements of $q$ are tied to the coefficients of a PI controller. The link has two origins.

\paragraph{The structure.} The first two components of $K_c$ multiply $\hat e_i$ and $\hat e$. The term $k_i\hat e_i+k_p\hat e$ of the control law is therefore a PI controller on the estimated error, and the book draws it as such in Figure~8.4. Each gain comes from the corresponding row of the Riccati solution.

\paragraph{The Riccati equation.} The book gives the dependence through empirical curves (Figures~8.9, 8.10). On a rigid plant it can be derived in closed form; the derivation is mine. Below the bandwidth of the velocity loop the plant is an integrator, $\dot y=u$, so with state $[y_i,\,y]$
\begin{equation}
A=\begin{bmatrix}0&1\\0&0\end{bmatrix},\quad B=\begin{Bmatrix}0\\1\end{Bmatrix},\quad Q=\begin{bmatrix}q_i&0\\0&q_p\end{bmatrix},\quad S_c=\begin{bmatrix}s_1&s_2\\s_2&s_3\end{bmatrix}.
\end{equation}
The elements of (8.6) give $q_i-s_2^2=0$, $s_1-s_2s_3=0$, $2s_2+q_p-s_3^2=0$. Since $K_c=B^TS_c=[s_2\ \ s_3]$,
\begin{equation}
k_i=\sqrt{q_i},\qquad k_p=\sqrt{q_p+2\sqrt{q_i}},\qquad\text{and conversely}\qquad q_i=k_i^2,\quad q_p=k_p^2-2k_i .
\end{equation}
The correspondence is one to one, so choosing $(q_i,q_p)$ is the same as choosing $(k_i,k_p)$. It also shows the coupling noted by the book: $q_p$ acts only on $k_p$, while $q_i$ acts on both. The book's statement that $q_p$ raises both gains refers to curves at constant $k_i/k_p$, where the two weights move together. The closed-loop poles are the roots of $s^2+k_ps+k_i$ (7.8), so the critical integral gain (7.9) becomes a condition on the weights:
\begin{equation}
k_i\leq0.25\,k_p^2\quad\Longleftrightarrow\quad q_p\geq2\sqrt{q_i}.
\end{equation}
The first guess of the book breaks it ($0.3<0.63$), which agrees with its overshoot of 18\,\%. The book's own values support the formula: $q_i=0.1$, $q_p=0.3$ give 0.316 and 0.966, against 0.304 and 0.9485 of its interface (Figure~8.17, with the first frequency weight at 5), and $q_i=3$, $q_p=10.05$ give 1.73 and 3.68 against 1.61 and 3.549.

\paragraph{Why it survives on the flexible plant.} For small $Q_f$ the tracking gains depend only weakly on the flexible part (Section~8.1.4), and the modes of LUCY lie far above the bandwidth. Table~\ref{tab:qgains} checks this numerically. $k_i=\sqrt{q_i}$ holds within 1\,\% on all models, and $k_p$ of the model without bearings is within about 1\,\% of the formula. With the Maxwell element $k_p$ is higher, by 39\,\% in azimuth and 8\,\% in elevation at the first guess, and less at high gains, because the slow pole of the element falls inside the low bandwidth. In this model $\dot y$ is not fed directly by $u$: $u$ acts through the velocity loop, so the integrator is a low-frequency description.

\begin{table}[H]\centering\footnotesize\setlength{\tabcolsep}{3.5pt}
\caption{$k_p$ [1/s] and $k_i$ [1/s$^2$] from the weights, $q_f=0$. Rigid: formula above and \texttt{dlqr} on $\dot y=u$ at 1\,kHz. Other columns: reduced synthesis models.}\label{tab:qgains}
\begin{tabular}{rrrrrrrrrrrr}\toprule
& & \multicolumn{2}{c}{formula} & \multicolumn{2}{c}{rigid, \texttt{dlqr}} & \multicolumn{2}{c}{azimuth, Maxwell} & \multicolumn{2}{c}{elevation, Maxwell} & \multicolumn{2}{c}{$k_p$, no bearings} \\
\cmidrule(lr){3-4}\cmidrule(lr){5-6}\cmidrule(lr){7-8}\cmidrule(lr){9-10}\cmidrule(lr){11-12}
$q_p$ & $q_i$ & $k_p$ & $k_i$ & $k_p$ & $k_i$ & $k_p$ & $k_i$ & $k_p$ & $k_i$ & az & el \\ \midrule
0.3 & 0.1 & 0.97 & 0.32 & 0.97 & 0.32 & 1.35 & 0.32 & 1.05 & 0.32 & 0.97 & 0.97 \\
25 & 4 & 5.39 & 2.00 & 5.37 & 1.99 & 6.20 & 2.00 & 5.69 & 2.00 & 5.45 & 5.42 \\
100 & 25 & 10.49 & 5.00 & 10.44 & 4.97 & 11.60 & 4.99 & 11.04 & 4.98 & 10.64 & 10.59 \\
400 & 100 & 20.49 & 10.00 & 20.29 & 9.90 & 21.55 & 9.96 & 21.22 & 9.95 & 20.72 & 20.66
\\ \bottomrule
\end{tabular}\end{table}

\subsection{Where $k_p$, $k_i$ are tuned in the procedure}
\textbf{The tuning of the first two elements of $q$ is ultimately the tuning of the coefficients $k_p$, $k_i$ of a PI controller.} It is also the first decision of the procedure: Section~8.1.4 says that the tuning starts with the ``initial choice of weights for the tracking subsystem'', and the first thing done in Step~1 is to fix small $q_i$, $q_p$, that is, a small PI. Table~\ref{tab:steps} places this in the steps of Section~8.2. There, Step~1 is named ``Tuning the flexible mode weights'', and $q_i$, $q_p$ are only set at the start of it, at small values, so that the resonances can be judged; they are tuned in Steps~2 and~3. These two steps are the tuning of a PI controller on a rigid antenna in Section~7.2, which also raises first the proportional and then the integral gain, and the resemblance is the point of Section~8.1.5.

\begin{table}[H]\centering\footnotesize
\caption{The elements of $q$ in the steps of Section~8.2.}\label{tab:steps}
\begin{tabular}{L{4.2cm}L{4.0cm}L{7.4cm}}\toprule
step of the book & elements of $q$ & effect \\ \midrule
Step 1, first action & $q_i$, $q_p$ fixed small & A small PI: on LUCY $k_p=\fgKpA$ and \fgKpP, $k_i=0.32$ \\
Step 1, flexible weights & $q_f$ & Resonances, raised only if too high; $k_p$, $k_i$ almost unchanged if the mode is above the bandwidth \\
Step 2, proportional weight & $q_p$ & $k_p$: bandwidth and disturbance rejection \\
Step 3, integral weight & $q_i$ & $k_i$, and part of $k_p$: rejection and overshoot, within (7.9) \\
Step 4, correction & $q_f$ & Resonances that resurface \\ \bottomrule
\end{tabular}\end{table}

\section{Audit of Section 8.2 of the book}\label{sec:audit}
Section~8.2 (pp.~108--110) gives four steps: (1) with small integral, position and flexible weights, check the closed-loop transfer function and raise the weight of a mode whose resonance is excessive, without unnecessarily large weights, because overdamped modes impair tracking; (2) raise the position weight until the bandwidth reaches the fundamental frequency; (3) raise the integral weight until oscillations appear, within the critical gain; (4) if resonances resurface, raise the corresponding flexible weights. Table~\ref{tab:audit} compares this text with its use on LUCY. The full audit is in \path{closedLoopModel/LQG_GAWRONSKI_PL_TUNING_AUDIT.md}, Section~8.

\begin{table}[H]\centering\footnotesize\setlength{\tabcolsep}{4pt}
\caption{Audit of Section~8.2, with the sections it depends on.}\label{tab:audit}
\begin{tabular}{L{4.4cm}L{5.8cm}L{5.8cm}}\toprule
point & book & LUCY \\ \midrule
Meaning of ``excessive'' & No threshold. The example of Section~8.4.2 speaks of ``sharp resonance peaks'' and ``visible flexible oscillations''. & Judged against the peaks of Figure~8.5 (D14). \\
First guess & ``Small'' in Section~8.2; $q_i=0.1$, $q_p=0.3$, frequency weights zero in Section~8.4.2. & Same $q_i$, $q_p$; $q_f=0.1$ on all flexible states, below the tracking weights, as (8.22) asks. \\
What to observe & $T_{ry}$ in Section~8.2; also the step and the disturbance response in Section~8.4.2. & $T_{ry}$ and 1$^\circ$ step on the reduced Maxwell model (D12). \\
Estimator weights & Absent from Section~8.2. $Q=V$ is the choice of (8.26); separate tuning ``may improve'' (p.~119). & $Q=V$ (D12), with $V$ the intensity of (8.8) (D15). Separate $V$ is left to the fine tuning. \\
Size of $q_f$ & Low authority, $\|Q_t\|\gg\|Q_f\|$ (8.22); $q_f=5$ against $q_p=10$ in the example. & Numbers not transferable: they depend on the scale of the modal states and on the decay rate of the mode. Here $q_f=0.1$ leaves the 56\,Hz mode untouched, and an effect starts near $10^3$ (Table~\ref{tab:s1mode}). \\
Independence of tracking and flexible parts & For $Q_f=0$, $S_{cf}=S_{ctf}=0$, so tracking gains do not depend on the flexible part (Section~8.1.4). & Only partly. $K_f\neq0$ at $q_f=0$, and $k_p$ is \fgKpA\ and \fgKpP\ against 0.97 for the rigid plant, because of the state of the Maxwell element: without bearings it is 0.97 on both axes (Table~\ref{tab:qgains}). The dependence of $k_p$ on the weight of the 56\,Hz mode is weak, as stated: \kpDropA\,\% and \kpDropP\,\% at $q_f=10^4$. \\
Stop rule of step 2 & Bandwidth equal to the fundamental frequency. & Does not bind, with modes at 56\,Hz. The limit is current (Section~8.1.7). \\
Condition of step 3 & Cites (8.15), which is the control law. & The critical gain is (7.9); the example of p.~104 ($k_p=16$, critical 64) confirms it. \\
Plant in the tuning & The same identified model for design and evaluation. & The same: reduced Maxwell model (D12). The Dahl plant is only a final check, outside the book. \\ \bottomrule
\end{tabular}\end{table}

\section{Step 1: flexible-mode weights}\label{sec:step1}
\subsection{Method}
Step~1 follows Section~8.2 to the letter (D12): small weights on the integral, the position and the flexible modes; check of the closed-loop transfer function; increase of the weight of the corresponding mode only if its resonance is too high; no unnecessarily large weights. All transfer functions and steps are computed on the reduced Maxwell model, with $Q=V$ and the 1-cycle delay.

\subsection{The vector of weights}
$q$ is a vector with one entry per state (Section~\ref{sec:weights}). The first guess is $q_i=0.1$, $q_p=0.3$ from Section~8.4.2, and the same small value $q_f=0.1$ on every flexible state, as in $Q_f=5I$ of Figure~8.5 but below the tracking weights, as (8.22) asks (Table~\ref{tab:s1vector}). A pair of equal entries belongs to a two-state mode.

\begin{table}[H]\centering\footnotesize
\caption{The selected vector $q$, in the order of the states. Frequencies: poles of the reduced model.}\label{tab:s1vector}
\begin{tabular}{lrrL{10.5cm}}\toprule
axis & $q_i$ & $q_p$ & $q_f$ \\ \midrule
azimuth & 0.1 & 0.3 & 0.05 Hz: 0.1; 4.20 Hz: $[0.1,\ 0.1]$; 13.32 Hz: $[0.1,\ 0.1]$; 56.36 Hz: $[0.1,\ 0.1]$; 65.10 Hz: $[0.1,\ 0.1]$ \\
elevation & 0.1 & 0.3 & 0.40 Hz: 0.1; 10.10 Hz: 0.1; 44.94 Hz: 0.1; 49.37 Hz: $[0.1,\ 0.1]$; 56.90 Hz: $[0.1,\ 0.1]$; 96.44 Hz: $[0.1,\ 0.1]$; 98.95 Hz: $[0.1,\ 0.1]$
\\ \bottomrule
\end{tabular}\end{table}

\subsection{Closed-loop transfer function and step}
Tables~\ref{tab:s1first} and~\ref{tab:s1modes} and Figure~\ref{fig:s1first} give the result. The book reports 8.52\,s, 18\,\% and 0.21\,Hz on the 34-m antenna at this stage (Section~8.4.2), so the low-frequency behaviour is comparable.

\begin{table}[H]\centering\footnotesize\setlength{\tabcolsep}{4pt}
\caption{First guess on the reduced Maxwell model. Step of 1$^\circ$: overshoot, settling to 2\,\%, and peak of $y$ in the band 30--150\,Hz over the whole response and from 0.1\,s after the step.}\label{tab:s1first}
\begin{tabular}{lrrrrrrr}\toprule
axis & $k_p$ [1/s] & $k_i$ [1/s$^2$] & band [Hz] & overshoot [\%] & $t_s$ [s] & vib.\ [$\mu$rad] & after 0.1\,s [$\mu$rad] \\ \midrule
azimuth & 1.37 & 0.32 & 0.12 & 28 & 12.1 & 3.0 & 0.35 \\
elevation & 1.05 & 0.32 & 0.20 & 21 & 9.6 & 1.6 & 0.27
\\ \bottomrule
\end{tabular}\end{table}

\begin{table}[H]\centering\footnotesize
\caption{Lightly damped modes at the first guess. Largest $T_{ry}$ in the band of the mode, and height of the peak above the higher of the minima on its sides.}\label{tab:s1modes}
\begin{tabular}{lrrrr}\toprule
axis & $f$ [Hz] & $\zeta$ & $T_{ry}$ [dB] & height [dB] \\ \midrule
azimuth & 56.4 & 0.023 & $-$59.5 & 5.9 \\
azimuth & 65.1 & 0.028 & $-$70.7 & 6.5 \\
elevation & 56.9 & 0.027 & $-$62.7 & 14.5 \\
elevation & 98.9 & 0.036 & $-$79.4 & 3.6
\\ \bottomrule
\end{tabular}\end{table}

\begin{figure}[H]\centering\includegraphics[width=\linewidth]{figures/s1_first_guess.pdf}\caption{First guess on the reduced Maxwell model: $T_{ry}$, with the levels of the peaks of Figure~8.5 of the book; zoom on the modes; 1$^\circ$ step.}\label{fig:s1first}\end{figure}

\subsection{Are the peaks too high?}
\textbf{No, against the only reference in the book.} In Figure~8.5 the peak near 2\,Hz, which the flexible weights remove, stands near 0\,dB, and the one near 5\,Hz remains near $-$6\,dB with $Q_f=5I$ (values read from the figure, so approximate). Here the largest peaks are \fgPkA\,dB in azimuth and \fgPkP\,dB in elevation, that is, more than 50\,dB lower. In the step they give \vibA\ and \vibP\,$\mu$rad on 17\,453\,$\mu$rad, invisible on the position plot, while the book sees oscillations in its step. The modes of LUCY are 280--450 times above the bandwidth; those of the book about 10--25 times. So no corresponding weight is raised at this step, and the vector stays as in Table~\ref{tab:s1vector} (D13).

The peaks are sharp in shape: 5.9\,dB in azimuth and 14.5\,dB in elevation above their surroundings. If a pointing requirement makes them relevant, the judgement changes.

\subsection{Process: when the system becomes overdamped}
The book warns that large flexible weights overdamp the antenna and reduce its bandwidth (Figure~8.8, Section~8.1.7). Table~\ref{tab:s1uniform} and Figure~\ref{fig:s1uniform} raise the same weight on all flexible states. The warning is confirmed: with $q_f\geq1$ in azimuth the bandwidth falls and the settling lengthens, and at $q_f=10$ the bandwidth is half. The cause is one state, not the modes at 56\,Hz, whose damping does not change. Weighting one state at a time with $q_f=10$ shows it: only the slow real pole added by the Maxwell element (0.05\,Hz in azimuth, 0.4\,Hz in elevation) reproduces the effect, with the same $k_p$ and bandwidth as the full vector, and every other state alone leaves them unchanged. This pole lies inside the bandwidth and is not a structural mode, so its weight must stay small. $q_f=0.1$ costs 2\,\% of bandwidth in azimuth and nothing in elevation.

\begin{table}[H]\centering\footnotesize\setlength{\tabcolsep}{4pt}
\caption{The same $q_f$ on all flexible states. $T_{ry}$ and $\zeta$ of the 56\,Hz mode. In bold, the selected value.}\label{tab:s1uniform}
\begin{tabular}{lrrrrrrrr}\toprule
axis & $q_f$ & $k_p$ & band [Hz] & overshoot [\%] & $t_s$ [s] & $T_{ry}$ [dB] & $\zeta$ & $\min|1+L|$ \\ \midrule
azimuth & 0 & 1.35 & 0.124 & 29 & 11.6 & $-$59.7 & 0.0231 & 0.67 \\
 & 0.01 & 1.35 & 0.124 & 29 & 11.7 & $-$59.7 & 0.0231 & 0.67 \\
 & \textbf{0.1} & \textbf{1.37} & \textbf{0.121} & \textbf{28} & \textbf{12.1} & \textbf{$-$59.5} & \textbf{0.0231} & \textbf{0.67} \\
 & 0.3 & 1.43 & 0.117 & 26 & 13.1 & $-$59.2 & 0.0231 & 0.69 \\
 & 1 & 1.57 & 0.104 & 22 & 15.5 & $-$58.3 & 0.0231 & 0.75 \\
 & 5 & 2.01 & 0.073 & 19 & $>$19 & $-$56.2 & 0.0231 & 0.91 \\
 & 10 & 2.33 & 0.060 & 19 & $>$19 & $-$54.9 & 0.0231 & 0.96 \\
 & 100 & 4.01 & 0.032 & 19 & $>$19 & $-$50.1 & 0.0233 & 0.94 \\
elevation & 0 & 1.05 & 0.200 & 21 & 9.5 & $-$62.7 & 0.0272 & 0.75 \\
 & 0.01 & 1.05 & 0.200 & 21 & 9.5 & $-$62.7 & 0.0272 & 0.75 \\
 & \textbf{0.1} & \textbf{1.05} & \textbf{0.199} & \textbf{21} & \textbf{9.6} & \textbf{$-$62.7} & \textbf{0.0272} & \textbf{0.74} \\
 & 0.3 & 1.05 & 0.199 & 21 & 9.6 & $-$62.7 & 0.0272 & 0.74 \\
 & 1 & 1.06 & 0.195 & 21 & 9.7 & $-$62.6 & 0.0272 & 0.74 \\
 & 5 & 1.10 & 0.179 & 20 & 10.4 & $-$62.2 & 0.0272 & 0.79 \\
 & 10 & 1.15 & 0.165 & 20 & 11.1 & $-$61.9 & 0.0273 & 0.83 \\
 & 100 & 1.58 & 0.101 & 19 & 16.5 & $-$59.1 & 0.0274 & 0.92
\\ \bottomrule
\end{tabular}\end{table}

\begin{figure}[H]\centering\includegraphics[width=\linewidth]{figures/s1_process_uniform.pdf}\caption{Process: the same weight on all flexible states. Large weights overdamp the response.}\label{fig:s1uniform}\end{figure}

\paragraph{What raising the corresponding weight would do.} Table~\ref{tab:s1mode} and Figures~\ref{fig:s1ladder} and~\ref{fig:s1trace} raise only the pair of the 56\,Hz mode, the others at 0.1. This is for the record, and for Step~4 of the book; no such weight is selected. Unlike the previous case the bandwidth and the step do not change, since this state lies far above the bandwidth, so the mode can be damped without overdamping the response: the damping goes from \fgZA\ to \zFourA\ in azimuth and from \fgZP\ to \zFourP\ in elevation at $q_f=10^4$, with the peak at \pkFourA\ and \pkFourP\,dB. Below about $10^3$ nothing happens. To first order the decay rate grows by $\Delta\sigma\simeq q_f\|b_i\|^2/(4\sigma)$; the mode at 56\,Hz has $\sigma\approx8$\,rad/s, and a large $\sigma$ needs a large $q_f$, so the book's 5 on a mode of a few hertz has no equivalent in the same number here. With $q_f=0.1$ the flexible controller practically does not touch the mode: this is a consequence of the small weights.

\begin{table}[H]\centering\footnotesize
\caption{Weight of the 56\,Hz pair only, other flexible states at 0.1; not selected.}\label{tab:s1mode}
\begin{tabular}{lrrrrrrr}\toprule
axis & $q_f$ & $\zeta$ & $T_{ry}$ [dB] & height [dB] & $\min|1+L|$ & band [Hz] & $k_p$ \\ \midrule
azimuth & 0.1 & 0.0231 & $-$59.5 & 5.9 & 0.67 & 0.121 & 1.374 \\
 & 1 & 0.0231 & $-$59.5 & 5.9 & 0.67 & 0.121 & 1.374 \\
 & 10 & 0.0231 & $-$59.5 & 5.9 & 0.67 & 0.121 & 1.374 \\
 & 100 & 0.0233 & $-$59.6 & 5.8 & 0.68 & 0.121 & 1.374 \\
 & $10^{3}$ & 0.0249 & $-$60.1 & 5.3 & 0.75 & 0.121 & 1.374 \\
 & $10^{4}$ & 0.0375 & $-$63.5 & 2.3 & 0.90 & 0.121 & 1.368 \\
 & $3\cdot10^{4}$ & 0.0560 & $-$66.3 & 0.1 & 0.58 & 0.121 & 1.360 \\
 & $10^{5}$ & 0.0957 & $-$67.6 & none & 0.18 & 0.121 & 1.343 \\
elevation & 0.1 & 0.0272 & $-$62.7 & 14.5 & 0.74 & 0.199 & 1.048 \\
 & 1 & 0.0272 & $-$62.7 & 14.5 & 0.74 & 0.199 & 1.048 \\
 & 10 & 0.0273 & $-$62.7 & 14.5 & 0.74 & 0.199 & 1.048 \\
 & 100 & 0.0274 & $-$62.8 & 14.5 & 0.73 & 0.199 & 1.048 \\
 & $10^{3}$ & 0.0292 & $-$63.3 & 13.9 & 0.78 & 0.199 & 1.047 \\
 & $10^{4}$ & 0.0433 & $-$66.5 & 10.6 & 0.97 & 0.199 & 1.042 \\
 & $3\cdot10^{4}$ & 0.0642 & $-$69.8 & 7.3 & 0.70 & 0.199 & 1.035 \\
 & $10^{5}$ & 0.1091 & $-$74.3 & 3.4 & 0.15 & 0.199 & 1.021
\\ \bottomrule
\end{tabular}\end{table}

\begin{figure}[H]\centering\includegraphics[width=\linewidth]{figures/s1_process_mode_ladder.pdf}\caption{What raising the weight of the 56\,Hz mode would do: $T_{ry}$ and step.}\label{fig:s1ladder}\end{figure}
\begin{figure}[H]\centering\includegraphics[width=\linewidth]{figures/s1_process_trace.pdf}\caption{Effect of the weight of the 56\,Hz mode on peak, damping, margin and bandwidth.}\label{fig:s1trace}\end{figure}

\subsection{Check on the Dahl plant}
This check is outside the procedure of the book and has not affected the selection (Table~\ref{tab:s1dahl}, Figure~\ref{fig:s1dahl}). The loop is stable on all plants. The step settles at 1000\,mdeg in azimuth. In elevation it remains about 4\,mdeg above the command after 120\,s, with $y-\hat y=$ \biasP\,mdeg: this is the bias of Section~\ref{sec:limits}, since the integrator brings $\hat y$, not $y$, to $r$. With $q_f=0$ it is 1.91\,mdeg. It is rechecked at the higher gains, and $V$ acts on it.

\begin{table}[H]\centering\footnotesize\setlength{\tabcolsep}{4pt}
\caption{Check of the first guess. Stability on the full Maxwell model and on the full plant with bearings sliding and stuck. 1$^\circ$ step on the Dahl plant, no current limit: overshoot [\%], range of $y$ [mdeg] from 100 to 120\,s, $y-\hat y$ [mdeg] at 120\,s, largest current [A].}\label{tab:s1dahl}
\begin{tabular}{llllllrrr}\toprule
axis & weights & full Maxwell & sliding & stuck & over. & $y$ & $y-\hat y$ & $i_\mathrm{max}$ \\ \midrule
azimuth & $q_f=0.1$ & stable & stable & stable & 11 & 1000.0--1000.0 & -0.03 & 5.5 \\
 & $q_f=0$ & stable & stable & stable & 16 & 999.9--1000.0 & -0.05 & 5.4 \\
elevation & $q_f=0.1$ & stable & stable & stable & 12 & 1004.1--1004.4 & 4.15 & 1.7 \\
 & $q_f=0$ & stable & stable & stable & 11 & 1001.9--1002.0 & 1.91 & 1.7
\\ \bottomrule
\end{tabular}\end{table}

\begin{figure}[H]\centering\includegraphics[width=\linewidth]{figures/s1_dahl_check.pdf}\caption{Check on the Dahl plant, 1$^\circ$ step.}\label{fig:s1dahl}\end{figure}

\subsection{Check at higher gains}
Since the peak in $T_{ry}$ scales with $k_p$, Table~\ref{tab:s1preview} anticipates what Step~2 will find on the model. This is not a tuning, and $q_p$, $q_i$ are not selected. At $k_p\approx$ \hiKpA\ the peaks reach \hiPkA\ and \hiPkP\,dB, still about 30\,dB below those of Figure~8.5. Step~4 of the book checks them at the final gains.

\begin{table}[H]\centering\footnotesize
\caption{Preview on the reduced Maxwell model, with $q_f=0.1$: largest $T_{ry}$ near 56\,Hz and $\min|1+L|$.}\label{tab:s1preview}
\begin{tabular}{lrrrrrr}\toprule
axis & $q_p$ & $q_i$ & $k_p$ & $k_i$ & $T_{ry}$ [dB] & $\min|1+L|$ \\ \midrule
azimuth & 25 & 4 & 6.2 & 2.0 & $-$46.4 & 0.92 \\
 & 100 & 25 & 11.6 & 5.0 & $-$40.9 & 0.87 \\
 & 400 & 100 & 21.6 & 10.0 & $-$35.5 & 0.76 \\
elevation & 25 & 4 & 5.7 & 2.0 & $-$48.0 & 0.76 \\
 & 100 & 25 & 11.0 & 5.0 & $-$42.2 & 0.78 \\
 & 400 & 100 & 21.2 & 10.0 & $-$36.5 & 0.80
\\ \bottomrule
\end{tabular}\end{table}

\section{Steps 2--4, iteration and $V$}\label{sec:steps24}
The objective, set for this phase, is low settling time and high bandwidth; the disturbance rejection is analysed afterwards. Transfer functions and steps are computed on the reduced Maxwell model with $Q=V$ (D12), starting from the vector of Step~1. The choices that the book leaves open are listed in Table~\ref{tab:arbitrary}. In the tables below, ``stop'' gives the limit violated by that point, and bold marks the selected value.

\begin{figure}[H]\centering\includegraphics[width=\linewidth]{figures/s24_step_2.pdf}\caption{Step~2: raising the proportional weight, $q_i=0.1$.}\label{fig:s2}\end{figure}

\subsection{Step 2: proportional weight}
Section~8.2 raises $q_p$ until the bandwidth reaches the fundamental frequency, 56\,Hz here, and Section~8.1.7 limits it through the acceleration, here the current (A1).
\begin{itemize}
\item \textbf{Azimuth: the current stops it} at $q_p=$ \sTwoA, on the 3$^\circ$ doublet. The bandwidth is 0.34\,Hz, very far from the fundamental. The current grows with $k_p$ because the loop tries to follow a command whose acceleration, followed exactly, would need about 48\,A: this is the limit described in Section~8.3.5.
\item \textbf{Elevation: the current never stops it,} with at most 16.7\,A, and the rule of the fundamental would be met only near $q_p\approx10^8$, $k_p\approx7400$\,1/s, at a 1\,kHz controller. Choice A3 stops it at $q_p=$ \sTwoP, where the plant with stuck bearings becomes unstable near 58\,Hz. This value does not come from the book.
\end{itemize}

\begin{table}[H]\centering\footnotesize\setlength{\tabcolsep}{3pt}
\caption{Step~2, with $q_i=0.1$. $T_{ry}$: largest value near the mode at 56\,Hz. $i$: largest current on the doublets, Dahl plant.}\label{tab:s2}
\begin{tabular}{lrrrrrrrrrl}\toprule
axis & $q_p$ & $k_p$ & $k_i$ & band [Hz] & $t_s$ [s] & over.\ [\%] & $i$ [A] & $\min|1+L|$ & $T_{ry}$ [dB] & stop \\ \midrule
azimuth & 0.3 & 1.374 & 0.316 & 0.121 & 12.1 & 27.9 & 3.6 & 0.67 & $-$59.5 & -- \\
 & 1 & 1.627 & 0.316 & 0.129 & 13 & 23.2 & 4.0 & 0.72 & $-$58.0 & -- \\
 & 3 & 2.179 & 0.316 & 0.145 & 15.6 & 16.6 & 4.8 & 0.79 & $-$55.5 & -- \\
 & 10 & 3.449 & 0.316 & 0.182 & 16.2 & 9.6 & 6.7 & 0.87 & $-$51.5 & -- \\
 & 30 & 5.656 & 0.316 & 0.242 & 9.73 & 5.5 & 10.5 & 0.95 & $-$47.2 & -- \\
 & \textbf{86} & \textbf{9.374} & \textbf{0.316} & \textbf{0.339} & \textbf{3.17} & \textbf{3.2} & \textbf{17.8} & \textbf{0.90} & \textbf{$-$42.8} & \textbf{--} \\
 & 89.3 & 9.547 & 0.316 & 0.344 & 3.1 & 3.2 & 18.2 & 0.90 & $-$42.6 & current \\
elevation & 0.3 & 1.048 & 0.316 & 0.199 & 9.56 & 20.9 & 2.2 & 0.74 & $-$62.7 & -- \\
 & 1 & 1.354 & 0.316 & 0.229 & 10.7 & 15.4 & 2.4 & 0.76 & $-$60.5 & -- \\
 & 3 & 1.973 & 0.316 & 0.288 & 11.7 & 9.5 & 2.7 & 0.77 & $-$57.2 & -- \\
 & $3\cdot10^{3}$ & 54.26 & 0.313 & 2.91 & 0.186 & 1.2 & 13.7 & 0.78 & $-$28.2 & -- \\
 & $1\cdot10^{4}$ & 98.49 & 0.311 & 4.82 & 0.117 & 0.8 & 16.7 & 0.75 & $-$22.9 & -- \\
 & $3\cdot10^{4}$ & 169.3 & 0.309 & 7.49 & 0.075 & 1.1 & 15.7 & 0.73 & $-$18.0 & -- \\
 & $1\cdot10^{5}$ & 305.2 & 0.305 & 11.5 & 0.072 & 2.5 & 14.9 & 0.69 & $-$12.7 & -- \\
 & $3\cdot10^{5}$ & 520.1 & 0.3 & 16.6 & 0.055 & 4.0 & 13.4 & 0.67 & $-$8.0 & -- \\
 & $1\cdot10^{6}$ & 925.9 & 0.293 & 24.3 & 0.086 & 5.6 & 11.9 & 0.63 & $-$3.5 & -- \\
 & \textbf{$1.3\cdot10^{6}$} & \textbf{1038} & \textbf{0.291} & \textbf{26.3} & \textbf{0.087} & \textbf{5.7} & \textbf{11.8} & \textbf{0.62} & \textbf{$-$2.7} & \textbf{--} \\
 & $1.3\cdot10^{6}$ & 1055 & 0.291 & 26.5 & 0.087 & 5.7 & 11.7 & 0.62 & $-$2.6 & unstable on a full plant
\\ \bottomrule
\end{tabular}\end{table}

\subsection{Step 3: integral weight}
Section~8.2 raises $q_i$ until oscillations appear, within (7.9). On neither axis do oscillations or the critical gain arrive first (A7, Figure~\ref{fig:s3}, Table~\ref{tab:s3}). In azimuth the current stops $q_i$ at \sThreeA, and $t_s$ gets worse, from 3.2 to 3.9\,s, as the overshoot grows; the book assigns to this weight the disturbance rejection, not the settling (Table~8.2). In elevation stability with stuck bearings stops it at \sThreeP, with $k_i\approx230$, 0.1\,\% of the critical gain.

\begin{table}[H]\centering\footnotesize\setlength{\tabcolsep}{3pt}
\caption{Step~3, at the $q_p$ of Step~2.}\label{tab:s3}
\begin{tabular}{lrrrrrrrrrl}\toprule
axis & $q_i$ & $k_p$ & $k_i$ & band [Hz] & $t_s$ [s] & over.\ [\%] & $i$ [A] & $\min|1+L|$ & $T_{ry}$ [dB] & stop \\ \midrule
azimuth & 0.1 & 9.374 & 0.316 & 0.339 & 3.17 & 3.2 & 17.8 & 0.90 & $-$42.8 & -- \\
 & \textbf{0.294} & \textbf{9.456} & \textbf{0.542} & \textbf{0.345} & \textbf{3.85} & \textbf{4.0} & \textbf{18.0} & \textbf{0.90} & \textbf{$-$42.7} & \textbf{--} \\
 & 0.305 & 9.459 & 0.551 & 0.345 & 3.89 & 4.0 & 18.0 & 0.90 & $-$42.7 & current \\
elevation & 0.1 & 1038 & 0.291 & 26.3 & 0.087 & 5.7 & 11.8 & 0.62 & $-$2.7 & -- \\
 & 0.316 & 1038 & 0.517 & 26.3 & 0.087 & 5.7 & 11.8 & 0.62 & $-$2.7 & -- \\
 & 1 & 1038 & 0.92 & 26.3 & 0.087 & 5.7 & 11.8 & 0.62 & $-$2.7 & -- \\
 & 100 & 1038 & 9.2 & 26.3 & 0.087 & 5.7 & 11.8 & 0.62 & $-$2.7 & -- \\
 & 316 & 1038 & 16.4 & 26.3 & 0.087 & 5.7 & 11.8 & 0.62 & $-$2.7 & -- \\
 & $1\cdot10^{3}$ & 1038 & 29.1 & 26.3 & 0.087 & 5.7 & 11.8 & 0.62 & $-$2.7 & -- \\
 & $3.2\cdot10^{3}$ & 1038 & 51.7 & 26.3 & 0.087 & 5.7 & 11.8 & 0.62 & $-$2.7 & -- \\
 & $1\cdot10^{4}$ & 1039 & 92 & 26.3 & 0.087 & 5.8 & 11.8 & 0.62 & $-$2.7 & -- \\
 & $3.2\cdot10^{4}$ & 1040 & 164 & 26.3 & 0.087 & 5.9 & 11.8 & 0.62 & $-$2.7 & -- \\
 & \textbf{$6.3\cdot10^{4}$} & \textbf{1040} & \textbf{230} & \textbf{26.3} & \textbf{0.087} & \textbf{6.0} & \textbf{11.8} & \textbf{0.62} & \textbf{$-$2.7} & \textbf{--} \\
 & $6.5\cdot10^{4}$ & 1041 & 234 & 26.3 & 0.087 & 6.0 & 11.8 & 0.62 & $-$2.7 & unstable on a full plant
\\ \bottomrule
\end{tabular}\end{table}

\begin{figure}[H]\centering\includegraphics[width=\linewidth]{figures/s24_step_3.pdf}\caption{Step~3: raising the integral weight.}\label{fig:s3}\end{figure}

\subsection{Step 4: correction of the flexible weights}
The peaks are judged as in D14. In azimuth they stay far below the threshold, at $-$43\,dB, so nothing changes. In elevation the peak at 56\,Hz has resurfaced, from $-$63 to $-$2.7\,dB, as the book foresees. Its corresponding weight is raised to $3\cdot10^4$, the others unchanged (Table~\ref{tab:s4}).

\begin{table}[H]\centering\footnotesize
\caption{Step~4. Weight of the pair of the 56\,Hz mode, the largest $T_{ry}$ near each mode, damping of the regulator pole.}\label{tab:s4}
\begin{tabular}{lrrrrrrr}\toprule
axis & $q_f$ & $T_{ry}$, 56\,Hz [dB] & $T_{ry}$, second mode [dB] & $\zeta$ & $\min|1+L|$ & band [Hz] & $k_p$ \\ \midrule
azimuth & 0.1 & $-$42.7 & $-$53.9 & 0.023 & 0.90 & 0.345 & 9.456 \\
elevation & 0.1 & $-$2.7 & $-$17.3 & 0.035 & 0.62 & 26.3 & 1040 \\
elevation & $1\cdot10^{3}$ & $-$3.1 & $-$17.3 & 0.037 & 0.62 & 26.3 & 1040 \\
elevation & $3\cdot10^{3}$ & $-$3.9 & $-$17.3 & 0.040 & 0.62 & 26.3 & 1039 \\
elevation & $1\cdot10^{4}$ & $-$5.8 & $-$17.2 & 0.049 & 0.62 & 26.3 & 1035 \\
elevation & $3\cdot10^{4}$ & $-$8.9 & $-$17.0 & 0.068 & 0.62 & 26.2 & 1028
\\ \bottomrule
\end{tabular}\end{table}

\subsection{Fine tuning: is another iteration worth it?}
Section~8.4.3 follows the manual tuning with an optimisation of (8.27); here it is reduced to a grid on the requested objective (A9, Table~\ref{tab:fine}).
\begin{itemize}
\item \textbf{Azimuth: yes.} The best point has lower $q_p$ and much higher $q_i$, at nearly the same $k_p$: the current depends on $k_p$, so $q_p$ can be exchanged for $q_i$, since $k_p^2\simeq q_p+2k_i$. The index improves by 13\,\%, with $k_i$ near $0.8k_p$, inside the range of the book. The price is the overshoot, from 4 to 21\,\%, which is outside the index.
\item \textbf{Elevation: no.} No point of the grid improves the index.
\end{itemize}
After this, $q_p$ is raised again by the rule of Step~2 (A10, Table~\ref{tab:raise}): to \sRaisedA in azimuth, stopped by the current, and to \sRaisedP in elevation, where the larger $q_f$ has moved the instability with stuck bearings further away. Peaks stay below the threshold.

\begin{table}[H]\centering\footnotesize\setlength{\tabcolsep}{4pt}
\caption{Fine tuning: the four best points that respect the limits. $f$: index. The reference, $f=0$, is the result of Step~4.}\label{tab:fine}
\begin{tabular}{lrrrrrrrrr}\toprule
axis & $f$ & $q_p$ & $q_i$ & $k_p$ & $k_i$ & band [Hz] & $t_s$ [s] & over.\ [\%] & $i$ [A] \\ \midrule
azimuth & $-$0.130 & 12.9 & 41.9 & 7.75 & 6.46 & 0.376 & 3.69 & 20.6 & 16.9 \\
 & 0.000 & 86 & 0.294 & 9.46 & 0.542 & 0.345 & 3.85 & 4.0 & 18.0 \\
 & 0.306 & 60.2 & 12.1 & 9.19 & 3.48 & 0.371 & 5.32 & 12.1 & 17.8 \\
 & 0.333 & 43 & 12.1 & 8.21 & 3.48 & 0.349 & 5.17 & 14.1 & 16.3 \\
elevation & 0.000 & $1.3\cdot10^{6}$ & $6.3\cdot10^{4}$ & 1.03e+03 & 228 & 26.2 & 0.051 & 5.7 & 11.9 \\
 & 0.127 & $8.9\cdot10^{5}$ & $6.3\cdot10^{4}$ & 869 & 230 & 23.4 & 0.052 & 5.4 & 12.1 \\
 & 0.197 & $1.8\cdot10^{6}$ & $6.3\cdot10^{4}$ & 1.21e+03 & 226 & 29.2 & 0.067 & 5.3 & 11.7 \\
 & 0.233 & $1.8\cdot10^{6}$ & $6.4\cdot10^{5}$ & 1.21e+03 & 720 & 29.3 & 0.069 & 5.8 & 11.7
\\ \bottomrule
\end{tabular}\end{table}

\begin{table}[H]\centering\footnotesize\setlength{\tabcolsep}{3pt}
\caption{$q_p$ raised after the fine tuning. Stops also include the resonance peak.}\label{tab:raise}
\begin{tabular}{lrrrrrrrrrl}\toprule
axis & $q_p$ & $k_p$ & $k_i$ & band [Hz] & $t_s$ [s] & over.\ [\%] & $i$ [A] & $\min|1+L|$ & $T_{ry}$ [dB] & stop \\ \midrule
azimuth & 12.9 & 7.747 & 6.46 & 0.376 & 3.69 & 20.6 & 16.9 & 0.84 & $-$44.4 & -- \\
 & 15.5 & 7.909 & 6.46 & 0.379 & 3.7 & 20.1 & 17.3 & 0.85 & $-$44.3 & -- \\
 & \textbf{19.4} & \textbf{8.147} & \textbf{6.46} & \textbf{0.384} & \textbf{3.72} & \textbf{19.4} & \textbf{17.7} & \textbf{0.87} & \textbf{$-$44.0} & \textbf{--} \\
 & 25.8 & 8.527 & 6.46 & 0.392 & 3.74 & 18.4 & 18.3 & 0.90 & $-$43.6 & current \\
elevation & $1.3\cdot10^{6}$ & 1028 & 228 & 26.2 & 0.051 & 5.7 & 11.9 & 0.62 & $-$8.9 & -- \\
 & $1.5\cdot10^{6}$ & 1121 & 227 & 27.7 & 0.059 & 5.5 & 11.8 & 0.61 & $-$8.2 & -- \\
 & $1.9\cdot10^{6}$ & 1245 & 225 & 29.8 & 0.067 & 5.3 & 11.6 & 0.60 & $-$7.3 & -- \\
 & \textbf{$2.5\cdot10^{6}$} & \textbf{1425} & \textbf{223} & \textbf{32.8} & \textbf{0.068} & \textbf{5.8} & \textbf{11.5} & \textbf{0.59} & \textbf{$-$6.2} & \textbf{--} \\
 & $3.8\cdot10^{6}$ & 1722 & 220 & 37.4 & 0.068 & 6.5 & 11.5 & 0.57 & $-$4.7 & unstable on a full plant, peak
\\ \bottomrule
\end{tabular}\end{table}

\subsection{Changing $V$}
\textbf{At fixed $Q$, $V$ cannot improve the objective on the model}: band and $t_s$ are the same in every row of Table~\ref{tab:v2} (Section~\ref{sec:v}). So $V$ can act in two ways only: on the plant that differs from the model, and on the margins, by leaving room to raise $q_p$. By rule A11, no value is adopted and $Q=V$ stays.
\begin{itemize}
\item \textbf{Small $V$ is bad on the Dahl plant.} The estimator follows the model and not the encoder: in azimuth at $\rho=10^{-4}$ the rms error is about 2700\,mdeg.
\item \textbf{Azimuth, large $V$: a real trade-off, not adopted.} With $\rho=100$--1000 the rms error falls from 585 to about 350\,mdeg, the overshoot on the doublets falls from 2284 to about 300\,mdeg, and the lower current allows $q_p$ to rise, with a 12\,\% larger bandwidth. However, the median settling on the doublets grows about threefold, so rule A11 rejects it. This option should be reconsidered in the analysis of the disturbance rejection.
\item \textbf{Elevation: nothing to gain.} $\rho$ between 0.01 and 0.3 changes little, smaller values worsen the error, $\rho\geq3$ is unstable with stuck bearings, and raising $q_p$ is stopped by the resonance peak.
\end{itemize}

\begin{table}[H]\centering\footnotesize\setlength{\tabcolsep}{3pt}
\caption{$V=\rho Q$ at the final $q$: band and $t_s$ on the Maxwell model, $\min|1+L|$; on the doublets, Dahl plant, largest current, median settling to 5\,mdeg, median rms error, largest overshoot. Last columns: largest $q_p$ within the limits, its band and index, referred to $\rho=1$.}\label{tab:v2}
\begin{tabular}{lrrrrrrrrrrr}\toprule
axis & $\rho$ & band [Hz] & $t_s$ [s] & $\min|1+L|$ & $i$ [A] & $t_\mathrm{set}$ [s] & err [mdeg] & over [mdeg] & $q_p$ & band [Hz] & $f$ \\ \midrule
azimuth & 0.0001 & 0.384 & 3.72 & 0.98 & 14.6 & 11.72 & 2674 & 22712 & 96.8 & 0.472 & $-$0.165 \\
 & 0.001 & 0.384 & 3.72 & 0.94 & 14.7 & 11.72 & 2482 & 18970 & 96.8 & 0.472 & $-$0.165 \\
 & 0.01 & 0.384 & 3.72 & 0.86 & 13.8 & 11.72 & 1992 & 11307 & 58.1 & 0.429 & $-$0.079 \\
 & 0.1 & 0.384 & 3.72 & 0.79 & 15.4 & 10.87 & 1217 & 4755 & 29 & 0.395 & $-$0.018 \\
 & 0.3 & 0.384 & 3.72 & 0.81 & 17.2 & 10.57 & 939 & 3535 & 23.2 & 0.388 & $-$0.007 \\
 & 1 & 0.384 & 3.72 & 0.87 & 17.7 & 3.06 & 585 & 2284 & 19.4 & 0.384 & -- \\
 & 3 & 0.384 & 3.72 & 0.92 & 16.6 & 8.28 & 404 & 1341 & 29 & 0.395 & $-$0.018 \\
 & 10 & 0.384 & 3.72 & 0.88 & 16.2 & 10.39 & 384 & 834 & 29 & 0.395 & $-$0.018 \\
 & 100 & 0.384 & 3.72 & 0.85 & 14.7 & 9.70 & 349 & 350 & 58.1 & 0.429 & $-$0.079 \\
 & 1000 & 0.384 & 3.72 & 0.86 & 13.5 & 9.66 & 348 & 298 & 58.1 & 0.429 & $-$0.079 \\
elevation & 0.0001 & 32.8 & 0.068 & 0.88 & 19.4 & 1.44 & 128 & 301 & -- & -- & -- \\
 & 0.001 & 32.8 & 0.068 & 0.75 & 16.1 & 0.99 & 41 & 106 & $2.5\cdot10^{6}$ & 32.8 & 0.000 \\
 & 0.01 & 32.8 & 0.068 & 0.63 & 12.9 & 0.78 & 19 & 29 & $2.5\cdot10^{6}$ & 32.8 & 0.000 \\
 & 0.1 & 32.8 & 0.068 & 0.59 & 11.7 & 0.68 & 19 & 16 & $2.5\cdot10^{6}$ & 32.8 & 0.000 \\
 & 0.3 & 32.8 & 0.068 & 0.59 & 11.6 & 0.66 & 19 & 14 & $2.5\cdot10^{6}$ & 32.8 & 0.000 \\
 & 1 & 32.8 & 0.068 & 0.59 & 11.5 & 0.65 & 20 & 14 & $2.5\cdot10^{6}$ & 32.8 & -- \\
 & 3 & 32.8 & 0.068 & 0.59 & 11.5 & 0.64 & 20 & 13 & -- & -- & -- \\
 & 10 & 32.8 & 0.068 & 0.59 & 11.5 & 0.64 & 20 & 13 & -- & -- & -- \\
 & 100 & 32.8 & 0.068 & 0.59 & 11.5 & 0.64 & 20 & 13 & -- & -- & -- \\
 & 1000 & 32.8 & 0.068 & 0.59 & 11.5 & 0.64 & 20 & 13 & -- & -- & --
\\ \bottomrule
\end{tabular}\end{table}

\begin{figure}[H]\centering\includegraphics[width=\linewidth]{figures/s24_effect_of_V.pdf}\caption{Effect of $V$.}\label{fig:v}\end{figure}

\subsection{Result}
\begin{table}[H]\centering\footnotesize
\caption{The final vector $q$, $Q=V$. Frequencies: poles of the reduced model; a pair belongs to a two-state mode.}\label{tab:finvector}
\begin{tabular}{lrrL{10.5cm}}\toprule
axis & $q_i$ & $q_p$ & $q_f$ \\ \midrule
azimuth & 41.9 & 19.4 & 0.05 Hz: 0.1; 4.20 Hz: $[0.1,\ 0.1]$; 13.32 Hz: $[0.1,\ 0.1]$; 56.36 Hz: $[0.1,\ 0.1]$; 65.10 Hz: $[0.1,\ 0.1]$ \\
elevation & $6.3\cdot10^{4}$ & $2.5\cdot10^{6}$ & 0.40 Hz: 0.1; 10.10 Hz: 0.1; 44.94 Hz: 0.1; 49.37 Hz: $[0.1,\ 0.1]$; 56.90 Hz: $[3\cdot10^{4},\ 3\cdot10^{4}]$; 96.44 Hz: $[0.1,\ 0.1]$; 98.95 Hz: $[0.1,\ 0.1]$
\\ \bottomrule
\end{tabular}\end{table}

\begin{table}[H]\centering\footnotesize
\caption{Final tuning against the present cascade, on the same plants: the cascade is also evaluated on the Maxwell model, so the band differs from that of a plant with no bearings.}\label{tab:final}
\begin{tabular}{lrrrr}\toprule
& \multicolumn{2}{c}{azimuth} & \multicolumn{2}{c}{elevation} \\
\cmidrule(lr){2-3}\cmidrule(lr){4-5}
& LQG-PI & cascade & LQG-PI & cascade \\ \midrule
$k_p$ [1/s]; $K_{pp}$ & 8.147 & 7.92 & 1425 & 8.02 \\
$k_i$ [1/s$^2$] & 6.463 & -- & 223.2 & -- \\
band, Maxwell [Hz] & 0.384 & 0.37 & 32.8 & 1.04 \\
$t_s$, 1$^\circ$, 2\,\%, Maxwell [s] & 3.72 & 4.68 & 0.068 & 1.62 \\
overshoot, Maxwell [\%] & 19.4 & 18.1 & 5.8 & 17.2 \\
largest current, doublets, Dahl [A] & 17.7 & 11.7 & 11.5 & 5.2 \\
median rms error, doublets, Dahl [mdeg] & 585 & 372 & 20 & 355 \\
median $t_\mathrm{set}$, doublets, Dahl [s] & 3.06 & 1.27 & 0.65 & 1.06 \\
largest overshoot, doublets, Dahl [mdeg] & 2284 & 43 & 14 & 11
\\ \bottomrule
\end{tabular}\end{table}

\begin{figure}[H]\centering\includegraphics[width=\linewidth]{figures/s24_final_tf_step.pdf}\caption{$T_{ry}$ and step on the reduced Maxwell model, along the tuning.}\label{fig:final}\end{figure}
\begin{figure}[H]\centering\includegraphics[width=\linewidth]{figures/s24_final_doublet.pdf}\caption{3$^\circ$ doublet on the Dahl plant, against the cascade.}\label{fig:doublet}\end{figure}

\paragraph{On the model, as requested, the objective improves on both axes,} by a great deal in elevation (band \finBwP\ against \casBwP\,Hz, $t_s$ \finTsP\ against \casTsP\,s) and little in azimuth (band \finBwA\ against \casBwA\,Hz, $t_s$ \finTsA\ against \casTsA\,s). The difference comes from the current: on the same command azimuth needs about four times the current for exact tracking, and it reaches the limit at $k_p\approx8$ while elevation does not reach it at all.

\paragraph{Four points must be kept in view.}
\begin{itemize}
\item \textbf{Azimuth is worse than the cascade on the Dahl plant.} The following error oscillates on the doublets (Figure~\ref{fig:doublet}): rms error \finErrA\ against \casErrA\,mdeg, overshoot \finOvDA\ against \casOvDA\,mdeg. The tuning on the Maxwell model does not see this. The cause, verified in Section~\ref{sec:compare}, is the Maxwell state in the estimator. The step of Step~3 with $q_p=86$, $q_i=0.29$ is also worse than the cascade, with 437\,mdeg and 991\,mdeg, so the result does not depend only on the fine tuning.
\item \textbf{The elevation gains are very high.} $k_p=\finKpP$\,1/s is about 1.4 per controller sample. On the doublets the bench gives a rms error of \finErrP\ against \casErrP\,mdeg, with a current within the limit, but these simulations have no encoder noise or quantisation. The settling on the 1$^\circ$ step is not achievable as such: on the Dahl plant, with no current limit, that step asks for \stepIP\,A, in azimuth \stepIA\,A, and in elevation it leaves \stepVibP\,$\mu$rad of vibration after 0.1\,s, with \stepOvP\,\% overshoot. Phase~2, with noise, may well reduce this result.
\item \textbf{Margin on the measured plant.} Outside the procedure, with the loop closed on the measured frequency response, $\min|1+L|$ is \mModP\ near \mHzP\,Hz in elevation, below the 0.5 of A3. In azimuth it is \mModA, but there $|L|$ stays at most \mLA\,dB in the band of the measurement, so it says nothing about the crossover.
\item \textbf{The bias does not appear at these gains:} $y-\hat y$ at 120\,s on the 1$^\circ$ step is \biasDA\ and \biasDP\,mdeg. The disturbance rejection has not been analysed.
\end{itemize}

\section{Comparison with the measured doublets, and investigation}\label{sec:compare}
This section answers two questions raised by the result of Section~\ref{sec:steps24}: why azimuth is worse than the cascade, and whether the elevation result is credible. It does not change the tuning. Everything here is evaluated on the Dahl plant, unlike the tuning.

\subsection{Comparison with the measurement}
The doublets were measured with the real cascade, so no measurement exists with the LQG. The comparison goes through the model of the cascade, which shows how much the bench can be trusted (Table~\ref{tab:cmd}, Figure~\ref{fig:cmd}). The bench reproduces the position well, with the median rms error 2\,\% below the measurement in azimuth and 10\,\% below in elevation, and the plateau at 5$^\circ$ is 12\,\% too low in elevation. The current is less good: the oscillation of the measured current after the start of the move, near 12\,Hz in azimuth, is not in the model. This is the limit of the bench that was already known from the fit of the current.

\begin{table}[H]\centering\footnotesize\setlength{\tabcolsep}{4pt}
\caption{One doublet per amplitude: M, measured; C, cascade on the model; L, final LQG-PI on the model.}\label{tab:cmd}
\begin{tabular}{lrrrrrrrrrr}\toprule
& & \multicolumn{3}{c}{rms error [mdeg]} & \multicolumn{3}{c}{largest error [mdeg]} & \multicolumn{3}{c}{largest current [A]} \\
\cmidrule(lr){3-5}\cmidrule(lr){6-8}\cmidrule(lr){9-11}
axis & stroke [$^\circ$] & M & C & L & M & C & L & M & C & L \\ \midrule
azimuth & 0.1 & 28 & 28 & 56 & 96 & 95 & 95 & 4.6 & 3.4 & 3.5 \\
 & 0.5 & 128 & 127 & 216 & 409 & 417 & 417 & 9.3 & 7.9 & 8.2 \\
 & 1 & 237 & 242 & 404 & 720 & 753 & 750 & 10.1 & 11.0 & 11.8 \\
 & 3 & 522 & 503 & 988 & 1292 & 1234 & 1838 & 10.3 & 11.7 & 14.4 \\
 & 5 & 669 & 628 & 1137 & 1294 & 1226 & 2284 & 10.5 & 11.7 & 17.7 \\
 & 8 & 976 & 942 & 767 & 1297 & 1283 & 1337 & 10.1 & 11.6 & 13.5 \\
elevation & 0.1 & 45 & 64 & 5 & 100 & 99 & 36 & 1.9 & 1.7 & 11.5 \\
 & 0.5 & 136 & 124 & 9 & 450 & 417 & 42 & 4.0 & 3.6 & 11.2 \\
 & 1 & 261 & 249 & 14 & 837 & 789 & 52 & 6.1 & 4.9 & 11.5 \\
 & 3 & 527 & 461 & 25 & 1285 & 1107 & 58 & 5.7 & 5.1 & 11.2 \\
 & 5 & 673 & 598 & 36 & 1297 & 1143 & 75 & 6.0 & 5.2 & 11.2 \\
 & 8 & 805 & 737 & 48 & 1306 & 1240 & 87 & 6.1 & 5.1 & 11.2
\\ \bottomrule
\end{tabular}\end{table}

\begin{figure}[H]\centering\includegraphics[width=\linewidth]{figures/cm_measured_doublets.pdf}\caption{Measured doublets, cascade on the model and final LQG-PI on the model.}\label{fig:cmd}\end{figure}

\subsection{Azimuth: where the performance is lost}
A chain of controllers changes one thing at a time, from a plain proportional controller with the gain of the cascade to the final LQG-PI (Table~\ref{tab:cmc}, Figure~\ref{fig:cma}).
\begin{itemize}
\item \textbf{The cascade can be reproduced.} A P on the encoder at the 1\,kHz of the master, with one cycle of delay, gives the same result. So does the P on the estimate with the no-bearing model.
\item \textbf{The loss comes in with the Maxwell state in the estimator:} the same P on the estimate goes from 44 to 1138\,mdeg of overshoot, and the dominant closed-loop pair on the plant in sliding falls from 2.0\,Hz with damping 0.63 to 1.1\,Hz with 0.30. The gains $k_i$ and $K_f$ are not the cause, since this row has neither. This is the risk ``prediction during large motion'' of Section~\ref{sec:limits}: during a move the model predicts a viscous torque of the Maxwell damper that the plant does not have, and with $Q=V$ the small $q_p$ gives a slow correction ($K_{e,y}\approx8$\,rad/s).
\item \textbf{This also explains the difference between the axes.} In elevation the same P on the estimate shows the same loss, 16 to 531\,mdeg. The final tuning does not, since $q_p$ is high and so $K_{e,y}\approx770$\,rad/s: the estimate follows the encoder and the model has little say.
\end{itemize}

\begin{table}[H]\centering\footnotesize\setlength{\tabcolsep}{3pt}
\caption{From the cascade to the final LQG-PI: doublets on the Dahl plant. Last column: complex closed-loop poles below 8\,Hz on the plant in sliding, frequency in Hz (damping).}\label{tab:cmc}
\begin{tabular}{lL{5.4cm}rrrrrl}\toprule
axis & controller & $i$ [A] & err [mdeg] & over [mdeg] & $t_\mathrm{set}$ [s] & $K_{e,y}$ [rad/s] & poles \\ \midrule
azimuth & Measured (real cascade) & 10.5 & 380 & 18 & 0.42 & -- & -- \\
 & Cascade & 11.7 & 372 & 43 & 1.27 & -- & -- \\
 & P on the encoder, $k=K_{pp}$ & 11.8 & 373 & 56 & 1.40 & -- & -- \\
 & P on the estimate, no-bearing model & 11.7 & 377 & 44 & 2.63 & 7.89 & 2.0 (0.63) \\
 & P on the estimate, Maxwell model & 14.2 & 430 & 1138 & 8.55 & 7.89 & 1.1 (0.30); 4.6 (0.40) \\
 & PI on the estimate, Maxwell model & 14.5 & 415 & 1183 & 7.52 & 7.89 & 1.1 (0.30); 4.6 (0.40) \\
 & LQG, Maxwell model, $k_p\approx K_{pp}$ & 14.8 & 468 & 901 & 10.87 & 7.89 & 1.2 (0.26); 4.4 (0.49) \\
 & LQG, no-bearing model, $k_p\approx K_{pp}$ & 9.6 & 447 & 96 & 10.87 & 7.89 & 2.0 (0.96) \\
 & Final LQG (Maxwell) & 17.7 & 585 & 2284 & 3.06 & 4.4 & 1.0 (0.24); 4.3 (0.48) \\
 & Final weights, no-bearing model & 8.5 & 509 & 1476 & 4.95 & 4.4 & -- \\
elevation & Measured (real cascade) & 6.1 & 394 & 5 & 1.12 & -- & -- \\
 & Cascade & 5.2 & 355 & 11 & 1.06 & -- & -- \\
 & P on the encoder, $k=K_{pp}$ & 5.2 & 355 & 16 & 1.02 & -- & -- \\
 & P on the estimate, no-bearing model & 5.2 & 363 & 2 & 0.99 & 7.99 & 2.3 (0.38); 1.9 (0.71) \\
 & P on the estimate, Maxwell model & 5.6 & 370 & 531 & 2.72 & 7.99 & 1.6 (0.26) \\
 & PI on the estimate, Maxwell model & 5.7 & 364 & 555 & 6.84 & 7.99 & 1.6 (0.26) \\
 & LQG, Maxwell model, $k_p\approx K_{pp}$ & 5.0 & 401 & 66 & 11.27 & 7.99 & 1.7 (0.33) \\
 & LQG, no-bearing model, $k_p\approx K_{pp}$ & 4.4 & 432 & 45 & 11.27 & 7.99 & 2.2 (0.68); 1.9 (0.71) \\
 & Final LQG (Maxwell) & 11.5 & 20 & 14 & 0.65 & 769 & -- \\
 & Final weights, no-bearing model & 9.0 & 40 & 16 & 0.67 & 769 & 1.9 (0.71)
\\ \bottomrule
\end{tabular}\end{table}

\begin{figure}[H]\centering\includegraphics[width=\linewidth]{figures/cm_azimuth_chain.pdf}\caption{Azimuth: the same P gain through different estimators, and the final LQG-PI.}\label{fig:cma}\end{figure}

\paragraph{Other configurations.} Table~\ref{tab:cmo} changes the internal model, $V$ and $q_i$, raising $q_p$ each time until the current limit of A1. These $q_p$ are selected on the Dahl plant, not on the Maxwell model, so they lie outside decision D12. With the no-bearing internal model and small $q_i$ the LQG-PI is better than the cascade at the same current limit: $k_p\approx25$, with an error 30\,\% lower, an overshoot of 28\,mdeg against 43, and a bandwidth on the plant in sliding of 3.1 against 2.1\,Hz, with a longer median settling (1.5--2.1 against 1.3\,s). It also beats a P on the encoder raised to the same current: at $k=12$ the P gives 16\,A, 284\,mdeg and 154\,mdeg of overshoot. With this model $V$ has almost no effect. With the Maxwell model a large $V$ recovers much of the loss, because the estimate then follows the encoder, but the overshoot and the settling remain worse. A large $q_i$ raises the overshoot on both models.

\begin{table}[H]\centering\footnotesize\setlength{\tabcolsep}{3pt}
\caption{Azimuth: configurations, with $q_p$ raised to the current limit and $q_f=0.1$. Doublets on the Dahl plant. Band: on the plant in sliding. $\min|1+L|$: on the measured plant.}\label{tab:cmo}
\begin{tabular}{lrrrrrrrrrrr}\toprule
internal model & $\rho$ & $q_i$ & $q_p$ & $k_p$ & $k_i$ & $i$ [A] & err [mdeg] & over [mdeg] & $t_\mathrm{set}$ [s] & band [Hz] & $\min|1+L|$ \\ \midrule
Maxwell & 1 & 0.1 & 69.1 & 8.4 & 0.32 & 15.4 & 463 & 918 & 10.87 & 2.11 & 0.91 \\
Maxwell & 100 & 0.1 & 94.2 & 9.8 & 0.32 & 15.9 & 336 & 120 & 10.87 & 2.76 & 0.85 \\
Maxwell & 100 & 41.88 & 53.4 & 10.0 & 6.46 & 17.6 & 312 & 326 & 10.30 & 2.81 & 0.83 \\
Maxwell & $1\cdot10^{4}$ & 0.1 & 126 & 11.3 & 0.32 & 16.9 & 309 & 132 & 10.87 & 3.01 & 0.85 \\
Maxwell & $1\cdot10^{4}$ & 41.88 & 69.1 & 10.7 & 6.46 & 17.1 & 298 & 237 & 10.31 & 2.98 & 0.85 \\
Maxwell & $1\cdot10^{6}$ & 0.1 & 126 & 11.3 & 0.32 & 16.8 & 309 & 129 & 10.87 & 3.01 & 0.85 \\
Maxwell & $1\cdot10^{6}$ & 41.88 & 69.1 & 10.7 & 6.46 & 17.0 & 298 & 238 & 10.30 & 2.99 & 0.85 \\
no bearings & 1 & 0.1 & 628 & 24.9 & 0.31 & 17.3 & 255 & 29 & 2.14 & 3.05 & 0.79 \\
no bearings & 1 & 41.88 & 628 & 25.3 & 6.41 & 17.6 & 250 & 271 & 8.11 & 3.13 & 0.78 \\
no bearings & 100 & 0.1 & 628 & 24.9 & 0.31 & 17.3 & 254 & 28 & 1.55 & 3.05 & 0.80 \\
no bearings & 100 & 41.88 & 628 & 25.3 & 6.41 & 17.6 & 249 & 271 & 7.64 & 3.13 & 0.79 \\
no bearings & $1\cdot10^{4}$ & 0.1 & 628 & 24.9 & 0.31 & 17.3 & 254 & 28 & 1.48 & 3.05 & 0.81 \\
no bearings & $1\cdot10^{4}$ & 41.88 & 628 & 25.3 & 6.41 & 17.6 & 249 & 271 & 7.59 & 3.13 & 0.81 \\
no bearings & $1\cdot10^{6}$ & 0.1 & 628 & 24.9 & 0.31 & 17.3 & 254 & 28 & 1.48 & 3.05 & 0.81 \\
no bearings & $1\cdot10^{6}$ & 41.88 & 628 & 25.3 & 6.41 & 17.6 & 249 & 271 & 7.59 & 3.13 & 0.81
\\ \bottomrule
\end{tabular}\end{table}

\subsection{Elevation: is the result credible?}
The size of the error has a simple cause. During the 10$^\circ$/s ramp of the doublets the following error of a type-1 loop is $v/k_p$: about 1250\,mdeg with $k_p=8$ and 7\,mdeg with $k_p=1425$. On this bench the gain is about 180 times that of the cascade, and the current does not limit it, so the improvement is an arithmetic consequence of the gain. Whether that gain is usable is another question. Table~\ref{tab:cme} and Figure~\ref{fig:cme} give what the bench can check.
\begin{itemize}
\item \textbf{Noise and delay do not limit it on the bench.} The measured position is exactly constant at rest, so the noise cannot be estimated from the data, and the 26-bit quantisation (\texttt{FB1.BITS}) has no effect. With an assumed white noise of 10\,$\mu$rad the error is unchanged and the peak current goes from 11.5 to 13.2\,A. Two more cycles of delay also change little. The phase voltage stays near 4\,V against a 34\,V bus. These amplitudes are my assumption (A12).
\item \textbf{The margins are what limits it.} On the measured plant the largest $T_{ry}$ in 1--150\,Hz is $+2.3$\,dB, against 0\,dB on the model, and $\min|1+L|$ is 0.46. With $g_{K_t}+10$\,\% the final gains are unstable on a full plant, while the step on the Dahl plant stays bounded. The measured response covers only 1--150\,Hz, so it does not prove stability on the real axis.
\item \textbf{The result is an extrapolation.} The bench has been validated with gains near 8\,1/s. At 1425\,1/s it depends on the model near and above the modes, and on effects that are not in it: cogging and torque ripple, the 8--9\,Hz pair of the measurement, the velocity limits of the drive, the timing of the master.
\item \textbf{Lower gains keep most of the gain.} With $k_p\approx307$ the error is 37\,mdeg, about 10 times below the cascade, with $\min|1+L|=0.62$ on the measured plant, stability for $g_{K_t}\pm10$\,\%, but 87\,mdeg of overshoot and 15\,A. A P on the encoder raised to $k=30$ gives 144\,mdeg and 11\,A, with 304\,mdeg of overshoot, so a large part of the gain is available without an LQG, by raising the gain.
\end{itemize}
The result in elevation should therefore be read as an upper bound from a noise-free bench, to be verified on the axis, and not as a prediction. Only a test on the real axis can settle it.

\begin{table}[H]\centering\footnotesize\setlength{\tabcolsep}{3pt}
\caption{Elevation at growing $q_p$, with the final $q_i$ (which sets a lower bound on $k_p$) and $q_f=3\cdot10^4$ from $q_p=10^5$. Doublets on the Dahl plant: nominal, with 10\,$\mu$rad of noise, with two more cycles of delay. Stable: on the plants sliding and stuck with $g_{K_t}\pm10$\,\%. Last columns: on the measured plant.}\label{tab:cme}
\begin{tabular}{rrrrrrrrcrr}\toprule
& & \multicolumn{3}{c}{nominal} & \multicolumn{2}{c}{noise} & delay & & \multicolumn{2}{c}{measured plant} \\
\cmidrule(lr){3-5}\cmidrule(lr){6-7}\cmidrule(lr){10-11}
$q_p$ & $k_p$ & err [mdeg] & $i$ [A] & over [mdeg] & err [mdeg] & $i$ [A] & err [mdeg] & stable & $\min|1+L|$ & $T_{ry}$ [dB] \\ \midrule
100 & 43 & 358 & 15.9 & 945 & 358 & 15.9 & 373 & yes & 0.62 & 3.1 \\
$1\cdot10^{3}$ & 53 & 132 & 18.2 & 582 & 132 & 18.2 & 143 & yes & 0.37 & 7.2 \\
$1\cdot10^{4}$ & 109 & 78 & 18.7 & 245 & 79 & 18.7 & 83 & yes & 0.47 & 6.6 \\
$1\cdot10^{5}$ & 307 & 37 & 15.3 & 87 & 38 & 15.6 & 38 & yes & 0.62 & 3.0 \\
$1\cdot10^{6}$ & 918 & 23 & 12.0 & 21 & 24 & 13.1 & 23 & yes & 0.52 & 1.6 \\
$2.5\cdot10^{6}$ & 1425 & 20 & 11.5 & 14 & 20 & 13.2 & 19 & no & 0.46 & 2.3
\\ \bottomrule
\end{tabular}\end{table}

\begin{figure}[H]\centering\includegraphics[width=\linewidth]{figures/cm_elevation_checks.pdf}\caption{Elevation: error, current and measured margin against the gain.}\label{fig:cme}\end{figure}

\subsection{Does the LQG at least match the cascade?}\label{sec:equiv}
With no disturbances an LQG is not expected to improve much on the cascade, but it should at least match it. Table~\ref{tab:equiv} tests this with the no-bearing internal model, on the Dahl plant, which is the reference from now on (see below). It searches $q_p$ for the same median rms error as the cascade. It does match it: at the same current the overshoot falls to zero and the settling is shorter. Two facts had hidden this.
\begin{itemize}
\item \textbf{$k_p$ is not comparable with $K_{pp}$.} The term $K_f\hat x_f$ also acts on the states of the velocity loop. With $k_p=K_{pp}$ the LQG has a larger following error, and it takes $k_p\approx1.5\,K_{pp}$ to match it. The row with $K_f=0$ isolates this term: without it the error is lower but the current and overshoot are higher, so it trades one for the other.
\item \textbf{The cascade has no position integral} ($K_{ip}=0$). With $q_i=0.1$ the overshoot returns and the median settling grows to about 6\,s, because the integrator must unwind what it accumulated during the ramp. The book assigns this weight to disturbance rejection (Table~8.2), which these tests do not contain. So on doublets with no disturbance $q_i$ has only a cost, and its benefit must be judged in the analysis of disturbances.
\end{itemize}

\begin{table}[H]\centering\footnotesize\setlength{\tabcolsep}{4pt}
\caption{No-bearing internal model, $q_f=0$: same error as the cascade. Doublets on the Dahl plant. $k$: $K_{pp}$ or $k_p$. Settling to 5\,mdeg, median and largest.}\label{tab:equiv}
\begin{tabular}{lL{6.4cm}rrrrrr}\toprule
axis & controller & $k$ [1/s] & $i$ [A] & err [mdeg] & over [mdeg] & $t_\mathrm{set}$ [s] & max [s] \\ \midrule
azimuth & Cascade & 7.92 & 11.7 & 372 & 43 & 1.27 & 2.63 \\
 & P on the encoder & 7.92 & 11.8 & 373 & 56 & 1.40 & 2.41 \\
 & LQG, $q_p=K_{pp}^2$, $q_i=0.1$ & 7.95 & 9.6 & 447 & 96 & 10.87 & 11.82 \\
 & \textbf{LQG, same error}, $q_i=1\cdot10^{-6}$, $\rho=1$ & 11.64 & 12.0 & 372 & 0 & 0.74 & 1.26 \\
 & \textbf{LQG, same error}, $q_i=1\cdot10^{-6}$, $\rho=1\cdot10^{4}$ & 11.64 & 12.0 & 370 & 0 & 0.52 & 1.02 \\
 & \textbf{LQG, same error}, $q_i=0.1$, $\rho=1$ & 11.68 & 12.0 & 371 & 51 & 5.68 & 11.62 \\
 & same, $K_f=0$ & 11.64 & 15.5 & 292 & 138 & 1.38 & 3.36 \\
elevation & Cascade & 8.02 & 5.2 & 355 & 11 & 1.06 & 4.43 \\
 & P on the encoder & 8.02 & 5.2 & 355 & 16 & 1.02 & 3.09 \\
 & LQG, $q_p=K_{pp}^2$, $q_i=0.1$ & 8.04 & 4.4 & 432 & 45 & 11.27 & 11.82 \\
 & \textbf{LQG, same error}, $q_i=1\cdot10^{-6}$, $\rho=1$ & 12.30 & 5.1 & 353 & 0 & 0.63 & 0.95 \\
 & \textbf{LQG, same error}, $q_i=1\cdot10^{-6}$, $\rho=1\cdot10^{4}$ & 12.04 & 5.1 & 352 & 1 & 0.64 & 0.68 \\
 & \textbf{LQG, same error}, $q_i=0.1$, $\rho=1$ & 12.34 & 5.1 & 352 & 22 & 6.03 & 11.62 \\
 & same, $K_f=0$ & 12.30 & 6.5 & 269 & 64 & 1.46 & 2.87
\\ \bottomrule
\end{tabular}\end{table}

\subsection{Are bearings part of the structural modes of the book?}\label{sec:book_bearings}
No. Section~2.2 derives the structural model from the finite-element model, as its modal model, including the rigid-body mode. Section~8.1.4 defines the flexible part as ``modes of deformation''. Dry friction is treated separately, as a non-linearity (Section~11.2), as a torque input in the time simulation of the velocity loop, and it is not in the LQG model of Chapter~8. Including the bearings in the internal model (D5, D8) was therefore an addition of this work to the method of the book, not part of it. The book does not state explicitly that bearings must be left out: this is a reading of what its models contain. It agrees with the result of Section~\ref{sec:compare}.

\paragraph{Reference for comparisons.} By decision of the user, from now on controllers are compared only with the model with Dahl bearings, the proxy for reality in this simulation study, and not with measured performance. The measurements remain valid for the identification of the model. Table~\ref{tab:cmd} was a requested check of the bench, not a reference for the controller.

\section{Tuning with the no-bearing internal model}\label{sec:nb}
By decision of the user, Steps~1--4 are repeated with the state estimated on a model without bearings (reduced model of the velocity loop: $f\leq150$\,Hz, Section~\ref{sec:step0}), consistent with the models of the book (Section~\ref{sec:book_bearings}), $Q=V$, on the same rules (D12--D15, A1--A11). The simulator stays the Dahl plant. Transfer functions and steps are on the reduced model; current limits and checks on the Dahl plant. All comparisons are with the cascade on the Dahl plant.

\subsection{Step 1}
With $q_i=0.1$, $q_p=0.3$, $q_f=0.1$ the low-frequency behaviour nearly coincides with the book (Table~\ref{tab:nb1}): $k_p$ and $k_i$ match the formula of Section~\ref{sec:weights}, and bandwidth, overshoot and settling are those of the 34-m antenna. The Maxwell model did not have this agreement. By D14 the peaks are not too high, so no weight is raised. Weighting all flexible states up to 100 lowers the bandwidth much less than with Maxwell (Figure~\ref{fig:nb1b}).

\begin{table}[H]\centering\footnotesize\setlength{\tabcolsep}{3pt}
\caption{First guess, no-bearing model. Peaks: 56\,Hz mode (with height), and second mode. Vibration: peak in the 1$^\circ$ step, 30--150\,Hz.}\label{tab:nb1}
\begin{tabular}{lrrrrrrrrr}\toprule
axis & $k_p$ & $k_i$ & band [Hz] & over.\ [\%] & $t_s$ [s] & $T_{ry}$ 56\,Hz [dB] & height [dB] & $T_{ry}$ 2nd [dB] & vib.\ [$\mu$rad] \\ \midrule
azimuth & 0.97 & 0.32 & 0.21 & 17 & 9.2 & $-$62.4 & 5.9 & $-$73.4 & 2.1 \\
elevation & 0.97 & 0.32 & 0.20 & 17 & 9.2 & $-$58.0 & 20.9 & $-$80.1 & 1.5
\\ \bottomrule
\end{tabular}\end{table}

\begin{figure}[H]\centering\includegraphics[width=\linewidth]{figures/nb1_first_guess.pdf}\caption{First guess on the reduced no-bearing model.}\label{fig:nb1a}\end{figure}
\begin{figure}[H]\centering\includegraphics[width=\linewidth]{figures/nb1_process_uniform.pdf}\caption{The same weight on all flexible states, no-bearing model.}\label{fig:nb1b}\end{figure}

\subsection{Steps 2 and 3}
Azimuth is stopped again by the current, now at $q_p=$ \nbQpA, against 86 with Maxwell, and elevation by the plant with stuck bearings, at \nbQpP (A3). In Step~3 the current stops azimuth at $q_i=$ \nbQiA, and the stuck plant stops elevation at \nbQiP. Oscillations and (7.9) are not reached.

\begin{table}[H]\centering\footnotesize\setlength{\tabcolsep}{3pt}
\caption{Step~2, no-bearing model, $q_i=0.1$. $i$: largest current on the doublets, Dahl plant.}\label{tab:nbs2}
\begin{tabular}{lrrrrrrrrrl}\toprule
axis & $q_p$ & $k_p$ & $k_i$ & band [Hz] & $t_s$ [s] & over.\ [\%] & $i$ [A] & $\min|1+L|$ & $T_{ry}$ [dB] & stop \\ \midrule
azimuth & 0.3 & 0.97 & 0.316 & 0.21 & 9.2 & 16.6 & 2.8 & 0.80 & $-$62.4 & -- \\
 & 1 & 1.283 & 0.316 & 0.253 & 9.95 & 11.5 & 3.2 & 0.82 & $-$60.0 & -- \\
 & 3 & 1.912 & 0.316 & 0.351 & 9.99 & 6.4 & 4.0 & 0.83 & $-$56.5 & -- \\
 & 10 & 3.266 & 0.316 & 0.584 & 5.92 & 2.7 & 5.5 & 0.81 & $-$51.9 & -- \\
 & 30 & 5.532 & 0.315 & 0.962 & 0.614 & 1.1 & 7.7 & 0.79 & $-$47.3 & -- \\
 & 100 & 10 & 0.315 & 1.6 & 0.354 & 0.4 & 11.0 & 0.76 & $-$42.1 & -- \\
 & 300 & 17.23 & 0.314 & 2.39 & 0.224 & 1.0 & 14.7 & 0.73 & $-$37.3 & -- \\
 & \textbf{768} & \textbf{27.48} & \textbf{0.313} & \textbf{3.25} & \textbf{0.16} & \textbf{2.0} & \textbf{18.0} & \textbf{0.71} & \textbf{$-$33.2} & \textbf{--} \\
 & 798 & 27.99 & 0.313 & 3.28 & 0.158 & 2.0 & 18.1 & 0.71 & $-$33.0 & current \\
elevation & 0.3 & 0.9689 & 0.316 & 0.204 & 9.21 & 16.5 & 2.2 & 0.80 & $-$58.0 & -- \\
 & 1 & 1.281 & 0.316 & 0.244 & 9.93 & 11.4 & 2.3 & 0.83 & $-$55.6 & -- \\
 & 3 & 1.909 & 0.316 & 0.338 & 9.87 & 6.4 & 2.6 & 0.84 & $-$52.1 & -- \\
 & $1\cdot10^{3}$ & 31.27 & 0.313 & 3.78 & 0.237 & 3.7 & 6.9 & 0.69 & $-$27.5 & -- \\
 & $3\cdot10^{3}$ & 53.86 & 0.311 & 5.01 & 0.187 & 4.4 & 8.1 & 0.69 & $-$22.7 & -- \\
 & $1\cdot10^{4}$ & 97.58 & 0.309 & 6.8 & 0.144 & 4.7 & 9.0 & 0.69 & $-$17.4 & -- \\
 & $3\cdot10^{4}$ & 167.4 & 0.306 & 9.05 & 0.109 & 4.7 & 9.3 & 0.68 & $-$12.6 & -- \\
 & $1\cdot10^{5}$ & 301.5 & 0.301 & 12.4 & 0.076 & 4.7 & 9.2 & 0.67 & $-$7.5 & -- \\
 & $3\cdot10^{5}$ & 513.3 & 0.296 & 16.8 & 0.07 & 5.2 & 9.1 & 0.66 & $-$3.6 & -- \\
 & \textbf{$6.4\cdot10^{5}$} & \textbf{736.4} & \textbf{0.292} & \textbf{20.9} & \textbf{0.087} & \textbf{5.7} & \textbf{9.0} & \textbf{0.64} & \textbf{$-$1.6} & \textbf{--} \\
 & $6.6\cdot10^{5}$ & 749.8 & 0.292 & 21.1 & 0.088 & 5.8 & 9.0 & 0.64 & $-$1.5 & unstable on a full plant
\\ \bottomrule
\end{tabular}\end{table}

\begin{figure}[H]\centering\includegraphics[width=\linewidth]{figures/nb2_step_2.pdf}\caption{Step~2, no-bearing model.}\end{figure}

\begin{table}[H]\centering\footnotesize\setlength{\tabcolsep}{3pt}
\caption{Step~3, no-bearing model.}\label{tab:nbs3}
\begin{tabular}{lrrrrrrrrrl}\toprule
axis & $q_i$ & $k_p$ & $k_i$ & band [Hz] & $t_s$ [s] & over.\ [\%] & $i$ [A] & $\min|1+L|$ & $T_{ry}$ [dB] & stop \\ \midrule
azimuth & 0.1 & 27.48 & 0.313 & 3.25 & 0.16 & 2.0 & 18.0 & 0.71 & $-$33.2 & -- \\
 & \textbf{0.246} & \textbf{27.49} & \textbf{0.491} & \textbf{3.25} & \textbf{0.232} & \textbf{2.0} & \textbf{18.0} & \textbf{0.71} & \textbf{$-$33.2} & \textbf{--} \\
 & 0.255 & 27.49 & 0.5 & 3.25 & 0.233 & 2.0 & 18.0 & 0.71 & $-$33.2 & current \\
elevation & 0.1 & 736.4 & 0.292 & 20.9 & 0.087 & 5.7 & 9.0 & 0.64 & $-$1.6 & -- \\
 & 0.316 & 736.4 & 0.519 & 20.9 & 0.087 & 5.7 & 9.0 & 0.64 & $-$1.6 & -- \\
 & 1 & 736.4 & 0.923 & 20.9 & 0.087 & 5.7 & 9.0 & 0.64 & $-$1.6 & -- \\
 & 3.16 & 736.5 & 1.64 & 20.9 & 0.087 & 5.7 & 9.0 & 0.64 & $-$1.6 & -- \\
 & 10 & 736.5 & 2.92 & 20.9 & 0.087 & 5.7 & 9.0 & 0.64 & $-$1.6 & -- \\
 & 31.6 & 736.5 & 5.19 & 20.9 & 0.087 & 5.7 & 9.0 & 0.64 & $-$1.6 & -- \\
 & 100 & 736.6 & 9.23 & 20.9 & 0.087 & 5.7 & 9.0 & 0.64 & $-$1.6 & -- \\
 & 316 & 736.7 & 16.4 & 20.9 & 0.088 & 5.8 & 9.1 & 0.64 & $-$1.6 & -- \\
 & $1\cdot10^{3}$ & 736.9 & 29.2 & 20.9 & 0.088 & 5.8 & 9.1 & 0.64 & $-$1.6 & -- \\
 & \textbf{$1\cdot10^{3}$} & \textbf{736.9} & \textbf{29.7} & \textbf{20.9} & \textbf{0.088} & \textbf{5.8} & \textbf{9.1} & \textbf{0.64} & \textbf{$-$1.6} & \textbf{--} \\
 & $1.1\cdot10^{3}$ & 736.9 & 30.3 & 20.9 & 0.088 & 5.8 & 9.1 & 0.64 & $-$1.6 & unstable on a full plant
\\ \bottomrule
\end{tabular}\end{table}

\begin{figure}[H]\centering\includegraphics[width=\linewidth]{figures/nb2_step_3.pdf}\caption{Step~3, no-bearing model.}\end{figure}

\subsection{Step 4, iteration and $V$}
In azimuth the peaks stay low and nothing changes. In elevation the peak at 56\,Hz resurfaces at $-1.6$\,dB, so the corresponding weight goes to $10^4$ (Table~\ref{tab:nbs4}). The fine tuning gives one iteration in azimuth, lowering $q_p$ to shorten the settling, after which $q_p$ is raised again to \nbQrA, and none in elevation, where $q_p$ rises to \nbQrP (Table~\ref{tab:nbraise}). $V$ does not change the objective on the model, and by rule A11 it stays $Q=V$. With this model a large $V$ does not hurt: in azimuth at $\rho=100$ the median settling on the doublets falls from 2.6 to 1.95\,s, with the same error (Figure~\ref{fig:nbv}).

\begin{table}[H]\centering\footnotesize
\caption{Step~4, no-bearing model.}\label{tab:nbs4}
\begin{tabular}{lrrrrrrr}\toprule
axis & $q_f$ & $T_{ry}$, 56\,Hz [dB] & $T_{ry}$, second mode [dB] & $\zeta$ & $\min|1+L|$ & band [Hz] & $k_p$ \\ \midrule
azimuth & 0.1 & $-$33.2 & $-$44.2 & 0.023 & 0.71 & 3.25 & 27.49 \\
elevation & 0.1 & $-$1.6 & $-$20.5 & 0.020 & 0.64 & 20.9 & 736.9 \\
elevation & $1\cdot10^{3}$ & $-$2.6 & $-$20.5 & 0.023 & 0.64 & 20.9 & 736.2 \\
elevation & $3\cdot10^{3}$ & $-$4.2 & $-$20.4 & 0.027 & 0.64 & 20.9 & 735.1 \\
elevation & $1\cdot10^{4}$ & $-$7.3 & $-$20.3 & 0.039 & 0.64 & 20.9 & 732.2
\\ \bottomrule
\end{tabular}\end{table}

\begin{table}[H]\centering\footnotesize\setlength{\tabcolsep}{3pt}
\caption{$q_p$ raised after the fine tuning, no-bearing model.}\label{tab:nbraise}
\begin{tabular}{lrrrrrrrrrl}\toprule
axis & $q_p$ & $k_p$ & $k_i$ & band [Hz] & $t_s$ [s] & over.\ [\%] & $i$ [A] & $\min|1+L|$ & $T_{ry}$ [dB] & stop \\ \midrule
azimuth & 538 & 23.04 & 0.492 & 2.91 & 0.179 & 1.7 & 16.7 & 0.72 & $-$34.7 & -- \\
 & \textbf{646} & \textbf{25.21} & \textbf{0.491} & \textbf{3.07} & \textbf{0.168} & \textbf{1.9} & \textbf{17.4} & \textbf{0.71} & \textbf{$-$33.9} & \textbf{--} \\
 & 807 & 28.16 & 0.491 & 3.29 & 0.233 & 2.0 & 18.2 & 0.71 & $-$32.9 & current \\
elevation & $6.4\cdot10^{5}$ & 732.2 & 29.5 & 20.9 & 0.053 & 5.3 & 9.0 & 0.64 & $-$7.3 & -- \\
 & \textbf{$7.6\cdot10^{5}$} & \textbf{798.8} & \textbf{29.4} & \textbf{22} & \textbf{0.06} & \textbf{5.6} & \textbf{9.0} & \textbf{0.64} & \textbf{$-$6.6} & \textbf{--} \\
 & $9.6\cdot10^{5}$ & 888.4 & 29.3 & 23.6 & 0.061 & 5.8 & 9.0 & 0.63 & $-$5.8 & unstable on a full plant, peak
\\ \bottomrule
\end{tabular}\end{table}

\begin{figure}[H]\centering\includegraphics[width=\linewidth]{figures/nb2_effect_of_V.pdf}\caption{Effect of $V$, no-bearing model.}\label{fig:nbv}\end{figure}

\subsection{Result of the tuning}
\begin{table}[H]\centering\footnotesize
\caption{The final vector $q$, no-bearing model, $Q=V$.}\label{tab:nbvector}
\begin{tabular}{lrrL{10.5cm}}\toprule
axis & $q_i$ & $q_p$ & $q_f$ \\ \midrule
azimuth & 0.246 & 646 & 0.43 Hz: 0.1; 2.60 Hz: 0.1; 13.37 Hz: $[0.1,\ 0.1]$; 55.98 Hz: $[0.1,\ 0.1]$; 65.11 Hz: $[0.1,\ 0.1]$ \\
elevation & $1\cdot10^{3}$ & $7.6\cdot10^{5}$ & 1.86 Hz: $[0.1,\ 0.1]$; 49.42 Hz: $[0.1,\ 0.1]$; 56.12 Hz: $[1\cdot10^{4},\ 1\cdot10^{4}]$; 96.38 Hz: $[0.1,\ 0.1]$; 98.25 Hz: $[0.1,\ 0.1]$
\\ \bottomrule
\end{tabular}\end{table}

\begin{table}[H]\centering\footnotesize
\caption{Final tuning against the cascade. Band, $t_s$, overshoot: on the reduced no-bearing model and, for the cascade, on the plant with bearings sliding.}\label{tab:nbfinal}
\begin{tabular}{lrrrr}\toprule
& \multicolumn{2}{c}{azimuth} & \multicolumn{2}{c}{elevation} \\
\cmidrule(lr){2-3}\cmidrule(lr){4-5}
& LQG-PI & cascade & LQG-PI & cascade \\ \midrule
$k_p$ [1/s]; $K_{pp}$ & 25.21 & 7.92 & 798.8 & 8.02 \\
$k_i$ [1/s$^2$] & 0.491 & -- & 29.4 & -- \\
band [Hz] & 3.07 & 2.11 & 22 & 2.69 \\
$t_s$, 1$^\circ$, 2\,\% [s] & 0.168 & 0.423 & 0.06 & 0.622 \\
overshoot [\%] & 1.9 & 4.4 & 5.6 & 4.6 \\
largest current, doublets, Dahl [A] & 17.4 & 11.7 & 9.0 & 5.2 \\
median rms error, doublets, Dahl [mdeg] & 253 & 372 & 50 & 355 \\
median $t_\mathrm{set}$, doublets, Dahl [s] & 2.60 & 1.27 & 0.40 & 1.06 \\
largest overshoot, doublets, Dahl [mdeg] & 38 & 43 & 8 & 11 \\
$\min|1+L|$, internal model & 0.71 & -- & 0.64 & --
\\ \bottomrule
\end{tabular}\end{table}

\begin{figure}[H]\centering\includegraphics[width=\linewidth]{figures/nb2_final_tf_step.pdf}\caption{$T_{ry}$ and step along the tuning, no-bearing model.}\end{figure}
\begin{figure}[H]\centering\includegraphics[width=\linewidth]{figures/nb2_final_doublet.pdf}\caption{3$^\circ$ doublet on the Dahl plant against the cascade, no-bearing model.}\end{figure}

\textbf{The oscillation of azimuth on the Dahl plant has gone,} and the tuning is better than the cascade on both axes, except for the median settling in azimuth, which is longer because of the integrator (2.6 against 1.3\,s) with a current closer to the limit. In elevation $k_p=$ \nbKpP\,1/s remains a very high gain, about half of that with Maxwell, and it depends on A3.

\section{Circle tracking, with disturbances and noise}\label{sec:circle}
\subsection{Setup}
The reference and the metrics are those of the study of the Hilkert observer: radius 10\,m, height 5\,m, centre at 100\,m, 0.2\,rad/s; 5\,s of transient left out, one orbit; pointing error of the line of sight, taken at the node that carries the telescope, and jitter above 1 and 10\,Hz. The plant is the Dahl plant, with the Coriolis, gyroscopic and inertial coupling between the axes and the azimuth inertia depending on elevation (Section~\ref{sec:coupling}), no current limit, and no feedforward (D7). Disturbances and noise are in A13. As well as the tuned controller the test includes three variants, only to read the sensitivity: $k_i\approx0.5k_p$, the typical value of the book; $V=100Q$; and $q_p=10^5$ on both axes with the tuned $q_i$ ($k_p\approx300$\,1/s, higher than tuned in azimuth, lower in elevation).

\subsection{Results}
\paragraph{How the error is computed.} Each axis receives its own reference from the same circle. The pointing error $E$ is the angle between the line of sight, taken at the node that carries the telescope, and the direction to the drone, and the table gives its rms over the orbit. This total contains the tracking error, about 100\,mdeg, and it is the wrong measure for a disturbance: an earlier version of this section concluded from it that the disturbances matter little. The effect of a disturbance is therefore measured as the difference of the line of sight from the same controller with no disturbance, on both axes.

\begin{table}[H]\centering\footnotesize\setlength{\tabcolsep}{3pt}
\caption{Circle on the Dahl plant. C: cascade; L: tuned LQG-PI; V: $V=100Q$; H: $q_p=10^5$.}\label{tab:cc}
\begin{tabular}{lrrrrrrrrrrrrr}\toprule
& \multicolumn{3}{c}{total $E$, rms [mdeg]} & \multicolumn{4}{c}{effect, rms [$\mu$rad]} & \multicolumn{3}{c}{effect, peak [$\mu$rad]} & \multicolumn{3}{c}{jitter $>10$\,Hz [$\mu$rad]} \\
\cmidrule(lr){2-4}\cmidrule(lr){5-8}\cmidrule(lr){9-11}\cmidrule(lr){12-14}
scenario & C & L & H & C & L & V & H & C & L & H & C & L & H \\ \midrule
nominal & 105 & 64.1 & 17.5 & 0 & 0 & 0 & 0 & 0 & 0 & 0 & 0.0285 & 0.0148 & 0.0177 \\
noise 1\,$\mu$rad & 105 & 64.1 & 17.5 & 0.159 & 11 & 14 & 7.18 & 1.56 & 98.7 & 80.6 & 0.0291 & 0.154 & 0.194 \\
noise 10\,$\mu$rad & 105 & 64.1 & 17.4 & 8.09 & 25.1 & 25.2 & 38.6 & 90.6 & 146 & 244 & 0.0664 & 1.64 & 1.9 \\
step 0.5\,N\,m & 103 & 63.1 & 17.6 & 131 & 116 & 85.4 & 44.8 & 1.12e+03 & 1.07e+03 & 272 & 0.608 & 0.575 & 0.908 \\
Gaussian 0.1\,N\,m & 104 & 64.1 & 17.5 & 160 & 69.9 & 57.9 & 53.6 & 642 & 220 & 229 & 8.98 & 8.01 & 10.7 \\
same, mean 0.2\,N\,m & 104 & 64.1 & 17.5 & 161 & 71.1 & 59.3 & 55.1 & 643 & 233 & 231 & 8.98 & 8.01 & 10.7 \\
sine 0.2\,N\,m, 1\,Hz & 108 & 67.5 & 19.8 & 548 & 378 & 285 & 168 & 1.66e+03 & 575 & 401 & 0.736 & 0.123 & 0.244 \\
sine 0.2\,N\,m, 10\,Hz & 104 & 64 & 17.5 & 166 & 49.3 & 45.9 & 66.3 & 747 & 199 & 242 & 13 & 11.9 & 25.2 \\
sine 0.2\,N\,m, 56\,Hz & 104 & 63.4 & 17.2 & 204 & 107 & 80.4 & 52.7 & 706 & 568 & 235 & 9.92 & 9.88 & 9.02 \\
noise, step, Gaussian with mean & 103 & 63.1 & 17.7 & 217 & 138 & 108 & 79.4 & 1.17e+03 & 1.06e+03 & 345 & 9.06 & 8.04 & 10.8
\\ \bottomrule
\end{tabular}\end{table}

\begin{figure}[H]\centering\includegraphics[width=\linewidth]{figures/cc_circle_nominal.pdf}\caption{Circle, no disturbances.}\end{figure}
\begin{figure}[H]\centering\includegraphics[width=\linewidth]{figures/cc_circle_scenarios.pdf}\caption{All scenarios: effect of the disturbance, total error and current.}\end{figure}
\begin{figure}[H]\centering\includegraphics[width=\linewidth]{figures/cc_circle_torque_step.pdf}\caption{Torque step of 0.5\,N\,m at 15\,s.}\end{figure}
\begin{figure}[H]\centering\includegraphics[width=\linewidth]{figures/cc_circle_all_disturbances.pdf}\caption{Noise, torque step and Gaussian torque with a mean of 0.2\,N\,m together.}\end{figure}

\begin{table}[H]\centering\footnotesize\setlength{\tabcolsep}{3pt}
\caption{Feedback and doublets of each case. $|L|$, $T_{wy}$ and band ($T_{ry}$): on the plant with bearings sliding. Doublets on the Dahl plant: rms error, largest current, largest overshoot, median settling.}\label{tab:ccx}
\begin{tabular}{llrrrrrrrr}\toprule
axis & case & $|L|$, 1\,Hz [dB] & $T_{wy}$, 1\,Hz [dB] & $T_{wy}$ peak [dB] & band [Hz] & err [mdeg] & $i$ [A] & over [mdeg] & $t_\mathrm{set}$ [s] \\ \midrule
azimuth & cascade & 2.6 & $-$51.5 & $-$51.5 & 2.11 & 372 & 11.7 & 43 & 1.27 \\
 & LQG-PI tuned & 2.6 & $-$51.8 & $-$51.7 & 3.07 & 253 & 17.4 & 38 & 2.60 \\
 & $k_i\approx0.5k_p$ & 2.9 & $-$51.7 & $-$51.6 & 3.22 & 246 & 18.0 & 417 & 5.35 \\
 & $V=100Q$ & 5.8 & $-$54.6 & $-$54.6 & 3.07 & 252 & 17.4 & 36 & 1.95 \\
 & $q_p=10^5$ & 16.5 & $-$65.1 & $-$65.1 & 9.69 & 71 & 46.2 & 11 & 0.27 \\
elevation & cascade & 3.8 & $-$46.3 & $-$42.9 & 2.69 & 355 & 5.2 & 11 & 1.06 \\
 & LQG-PI tuned & 20.2 & $-$61.9 & $-$59.2 & 22.2 & 50 & 9.0 & 8 & 0.40 \\
 & $k_i\approx0.5k_p$ & 20.2 & $-$61.9 & $-$59.2 & 22.3 & 49 & 9.1 & 44 & 2.16 \\
 & $V=100Q$ & 20.6 & $-$62.3 & $-$59.6 & 22.2 & 50 & 9.0 & 8 & 0.47 \\
 & $q_p=10^5$ & 16.0 & $-$57.8 & $-$55.0 & 12.4 & 78 & 9.2 & 24 & 1.37
\\ \bottomrule
\end{tabular}\end{table}

\begin{figure}[H]\centering\includegraphics[width=\linewidth]{figures/cc_disturbance_rejection.pdf}\caption{Disturbance rejection on the plant with bearings sliding.}\label{fig:twy}\end{figure}

\subsection{Reading}
\begin{itemize}
\item \textbf{Tracking.} The total error is mostly the azimuth tracking. The tuned LQG-PI lowers it from 105 to 64\,mdeg, $k_i\approx0.5k_p$ to 26\,mdeg, and $q_p=10^5$ to 17\,mdeg. The current limit is far on the circle: at most 4.5\,A in azimuth.
\item \textbf{Torque disturbances.} The tuned LQG-PI gains little in azimuth (Table~\ref{tab:ccx}), and much in elevation. With $q_p=10^5$ the effect on the azimuth line of sight falls from 115 to 38\,$\mu$rad on the torque step and from 365 to 87\,$\mu$rad on the sine at 1\,Hz, but it grows at 10\,Hz (total jitter 25 against 12\,$\mu$rad). This variant is not admissible on the doublets, where it asks for 46\,A, so A1 decides this result.
\item \textbf{Gaussian torque and its mean.} The Gaussian torque of 0.1\,N\,m rms moves the line of sight by 160\,$\mu$rad rms with the cascade and 70 with the tuned LQG-PI. The earlier process noise, not Gaussian near the edges, gave 149 and 59. A mean of 0.2\,N\,m changes these values by 0.6 to 1.5\,$\mu$rad. It moves the line of sight by a constant 12.7\,$\mu$rad in azimuth, the static deflection between encoder and telescope, because the velocity integrator of the drive holds it (Section~\ref{sec:thesis}). The largest azimuth current rises by 0.86\,A.
\item \textbf{Noise.} With the tuned controller 1\,$\mu$rad moves the elevation line of sight by 11\,$\mu$rad rms, against 0.1 of the cascade, and 10\,$\mu$rad by 25 against 8. This is larger than the noise itself, probably because it changes the stick-slip of the axis near the reversals; this was verified only for the sine at 56\,Hz. It stays below the effect of the torque disturbances tested.
\end{itemize}

\subsection{Is the simulation right?}
Table~\ref{tab:sines}, on each axis alone, injects a sine at one frequency at a time and reads the amplitude at that frequency, on the difference from the case with no disturbance. In azimuth the simulation agrees with the linear $T_{wy}$ with bearings sliding, within about 10\,\%. In elevation it lies between sliding and stuck, as expected at the low elevation rate, and near 56\,Hz it is close to stuck. There is no error in the injection or in the metric.

\begin{table}[H]\centering\footnotesize\setlength{\tabcolsep}{3pt}
\caption{Sine of 0.2\,N\,m at the bearing port: amplitude [$\mu$rad] at encoder and line of sight on the Dahl plant, against the linear prediction at the encoder with bearings sliding and stuck.}\label{tab:sines}
\begin{tabular}{lrrrrrrrrr}\toprule
& & \multicolumn{4}{c}{cascade} & \multicolumn{4}{c}{tuned LQG-PI} \\
\cmidrule(lr){3-6}\cmidrule(lr){7-10}
axis & $f$ [Hz] & enc. & LOS & sliding & stuck & enc. & LOS & sliding & stuck \\ \midrule
azimuth & 0.3 & 357 & 368 & 361 & 96 & 352 & 362 & 356 & 96.4 \\
 & 1 & 480 & 495 & 531 & 101 & 463 & 477 & 516 & 101 \\
 & 3 & 217 & 227 & 240 & 135 & 231 & 244 & 258 & 128 \\
 & 10 & 20.2 & 28.3 & 19.9 & 24.1 & 19.9 & 27.1 & 19.6 & 23.7 \\
 & 30 & 0.674 & 3.21 & 0.674 & 0.671 & 0.671 & 3.21 & 0.671 & 0.667 \\
 & 56 & 43.5 & 13.8 & 45.2 & 42.9 & 43.3 & 13.8 & 44.9 & 42.6 \\
 & 100 & 1.19 & 0.0938 & 1.19 & 1.2 & 1.19 & 0.0938 & 1.19 & 1.2 \\
elevation & 0.3 & 383 & 383 & 383 & 11.1 & 50.1 & 50.1 & 50.1 & 10.5 \\
 & 1 & 199 & 199 & 964 & 11.1 & 150 & 150 & 160 & 10.6 \\
 & 3 & 63.9 & 63.8 & 880 & 11.3 & 63.4 & 63.4 & 180 & 10.7 \\
 & 10 & 22.3 & 22.1 & 61.5 & 14.1 & 19.9 & 19.7 & 57.4 & 12 \\
 & 30 & 8.27 & 7.2 & 7.55 & 13.4 & 19.2 & 16.8 & 11.1 & 30.6 \\
 & 56 & 11.1 & 0.948 & 30.9 & 10.4 & 12.5 & 1.04 & 31.7 & 11.8 \\
 & 100 & 4.92 & 0.602 & 4.23 & 5.17 & 5.39 & 0.683 & 4.59 & 5.69
\\ \bottomrule
\end{tabular}\end{table}

\paragraph{Why a sine at the resonance has little effect.} In azimuth 0.2\,N\,m at 56\,Hz gives 43\,$\mu$rad at the encoder and only 14\,$\mu$rad at the line of sight, which moves about a third of the encoder in this mode, that is 0.8\,mdeg against 104\,mdeg of tracking. In the total rms it does not show, and it does appear in the jitter, 9.9\,$\mu$rad. The rms effect of this scenario is larger, 105\,$\mu$rad, for another reason: most of it is below 10\,Hz and near the velocity reversals (Table~\ref{tab:dither}). The sine acts as a dither on the friction, the mechanism of Section~11.2 of the book, and it \emph{lowers} the error: in elevation from 11.3 to 7.3\,mdeg with the cascade.

\begin{table}[H]\centering\footnotesize
\caption{Effect of the 56\,Hz sine on the line of sight [$\mu$rad]: below 10\,Hz; below 10\,Hz more than 3\,s from the reversals of the reference; above 30\,Hz. Error rms of that axis [mdeg] without and with the sine.}\label{tab:dither}
\begin{tabular}{llrrrrr}\toprule
axis & case & $<10$\,Hz & far from reversals & $>30$\,Hz & error, nominal & error, sine \\ \midrule
azimuth & cascade & 105 & 16 & 9.85 & 104 & 104 \\
 & LQG-PI tuned & 103 & 14.6 & 9.81 & 64.2 & 63.5 \\
elevation & cascade & 175 & 149 & 0.652 & 11.3 & 7.3 \\
 & LQG-PI tuned & 30.1 & 20.2 & 0.733 & 1.67 & 1.05
\\ \bottomrule
\end{tabular}\end{table}

\subsection{Why the jitter at the mode is not damped}
At the mode, $|L|$ is $-57$\,dB in azimuth and the flexible weight of the mode is 0.1 (Table~\ref{tab:nbvector}): the flexible controller practically does not act, as Step~1 had already shown. The large $q_f=10^4$ is in elevation. $q_i$ acts on the integral, at low frequency, and has no effect at 56\,Hz. Table~\ref{tab:qfaz} raises the weight in azimuth: it does damp, and it costs stability with stuck bearings at $10^5$. In elevation $q_f=10^4$ has no visible effect either, and Section~\ref{sec:twymodes} gives the cause: $q_f$ acts on the regulator, while the peak of $T_{wy}$ is held by the estimator. In Table~\ref{tab:qfaz} the weight moves both, since $V=Q$.

\begin{table}[H]\centering\footnotesize
\caption{Azimuth: raising the weight of the 56\,Hz mode, tuned weights otherwise. Amplitude on the Dahl plant with a sine of 0.2\,N\,m.}\label{tab:qfaz}
\begin{tabular}{lrrrc}\toprule
$q_f$ & enc., 56\,Hz [$\mu$rad] & LOS, 56\,Hz [$\mu$rad] & enc., 1\,Hz [$\mu$rad] & stable \\ \midrule
0.1 & 43.3 & 13.8 & 463 & yes \\
$10^4$ & 40.5 & 13.0 & 420 & yes \\
$3\cdot10^4$ & 34.4 & 11.1 & 244 & yes \\
$10^5$ & 24.2 & 7.9 & 1004 & no
\\ \bottomrule
\end{tabular}\end{table}

\subsection{Why $T_{wy}$ in elevation does not fall at the 56 and 98\,Hz modes}\label{sec:twymodes}
In elevation $T_{wy}$ keeps its peaks at the two modes, although the 56\,Hz mode has $q_f=10^4$. This section is a diagnosis on the elevation axis only: the tuned controller is not changed. Here $Q_f$ and $V_f$ are the weights of the mode in $Q$ and $V$, separated as the book allows (p.~119).

\paragraph{The peak belongs to the estimator.} With the estimation error $\varepsilon=x-\hat x$, on the model and without the delay,
\begin{equation}
x^+=(A-BK)x+BK\varepsilon+B_ww,\qquad\varepsilon^+=(A-K_eC)\varepsilon+B_ww.
\end{equation}
A command enters the plant and the estimator alike and does not excite $\varepsilon$, so $T_{ry}$ only contains the regulator poles, which depend on $Q$. A disturbance enters the plant alone and excites $\varepsilon$, so $T_{wy}$ also contains the estimator poles, which depend on $V$. The tuning read $T_{ry}$, and $q_f$ did lower its peak, from $-1.3$ to $-6.7$\,dB. In $T_{wy}$, however, the least damped pole decides. With $V_f=10^4$ the estimator pole stays almost at the open-loop damping, $\zeta=0.015$ against 0.014, while the regulator pole has 0.039 (Table~\ref{tab:twymodes}). On the plant with bearings sliding, which is the internal model, the closed loop indeed has both, at 56.1\,Hz with $\zeta=0.015$ and at 56.8\,Hz with $\zeta=0.043$.

\paragraph{The same from the loop.} $T_{wy}=P_w/(1+L)$ with $L=F_yP_u$. $F_y$ has zeros at 56.2\,Hz with $\zeta=0.014$, on the poles of the plant, and poles at 57.2\,Hz with $\zeta=0.041$, so it is a notch centred on the mode. The zeros cancel the poles of $P_u$ in $L$, and those of $P_w$ stay in $T_{wy}$. Raising $q_f$ makes the notch deeper and lowers the largest $|L|$ near the mode, from 1.11 to 0.46. So at this level $q_f$ removes feedback at the mode: the command no longer excites it, but a disturbance is not rejected there.

\paragraph{Raising $V_f$.} With $Q_f=10^4$, $V_f=10^6$ takes the estimator pole to $\zeta=0.062$, $T_{wy}$ falls by 6.4\,dB and $T_{ry}$ does not change. Further on it stops at $-85.1$\,dB, the limit of the regulator pole: $20\log_{10}(0.043/0.014)=9.7$\,dB below the open-loop peak of $-76.2$\,dB. With $Q_f=0.1$ the same limit is $-80$\,dB. So both weights are needed.

\begin{table}[H]\centering\footnotesize\setlength{\tabcolsep}{3pt}
\caption{Elevation, 56\,Hz mode. Regulator and estimator pole (frequency [Hz], $\zeta$); peak of $T_{ry}$ between 40 and 70\,Hz; peak of $T_{wy}$ and its frequency with bearings sliding and stuck; largest $|z|$ if unstable with stuck bearings (all are stable when sliding). Without position loop the peak of $T_{wy}$ is $-76.2$\,dB at 56.1\,Hz sliding and at 58.5\,Hz stuck.}\label{tab:twymodes}
\begin{tabular}{llllrrrrrl}\toprule
& & & & & \multicolumn{2}{c}{$T_{wy}$, sliding} & \multicolumn{2}{c}{$T_{wy}$, stuck} & \\
\cmidrule(lr){6-7}\cmidrule(lr){8-9}
$Q_f$ & $V_f$ & regulator & estimator & $T_{ry}$ [dB] & [dB] & [Hz] & [dB] & [Hz] & stuck \\ \midrule
0.1 & 0.1 & 56.1, 0.021 & 56.1, 0.014 & -1.3 & -74.2 & 55.9 & -67.3 & 57.2 & \textbf{1.0007} \\
$10^4$ & $10^4$ & 56.1, 0.039 & 56.1, 0.015 & -6.7 & -76.0 & 56.0 & -60.1 & 57.8 & stable \\
0.1 & $10^6$ & 56.1, 0.021 & 56.3, 0.062 & -1.3 & -79.2 & 56.0 & -75.0 & 58.2 & stable \\
$10^4$ & $10^6$ & 56.1, 0.039 & 56.3, 0.062 & -6.7 & -82.4 & 56.0 & -79.7 & 58.9 & stable \\
$10^4$ & $10^8$ & 56.1, 0.039 & 67.0, 0.498 & -6.7 & -84.8 & 56.4 & -85.1 & 59.3 & \textbf{1.0015} \\
$10^4$ & $10^{12}$ & 56.1, 0.039 & $>400$ & -6.7 & -85.1 & 56.4 & -85.6 & 59.2 & \textbf{1.0074} \\
$10^6$ & $10^6$ & 58.9, 0.306 & 56.3, 0.062 & -15.2 & -87.0 & 55.5 & -88.3 & 58.6 & \textbf{1.0021}
\\ \bottomrule
\end{tabular}\end{table}

\begin{figure}[H]\centering\includegraphics[width=\linewidth]{figures/cc_twy_modes.pdf}\caption{Elevation, around the 56\,Hz mode. Lines: linear model with bearings sliding and stuck. Circles: sines on the Dahl plant, with the axis moving on the circle and at rest.}\label{fig:twymodes}\end{figure}

\paragraph{With stuck bearings the tuned controller amplifies.} With stuck bearings the mode moves to 58.5\,Hz, while the notch stays at 56.1\,Hz. The tuned closed loop then has a pole at 57.8\,Hz with $\zeta=0.001$, and $T_{wy}$ reaches $-60.1$\,dB, 16\,dB above the cascade. The Dahl plant confirms it, Table~\ref{tab:twydahl}: with the axis at rest, a sine of 0.002\,N\,m at 57.8\,Hz gives $-62.5$\,dB against $-78.2$\,dB of the cascade. As the amplitude grows the bearing breaks away and the peak falls to $-79.7$\,dB at 0.2\,N\,m. With the axis moving, small sines follow the sliding model, since the bearing does not reverse. With $V_f=10^6$ the least damped pole has $\zeta=0.016$ and the same sine gives $-82.3$\,dB, 4\,dB better than the cascade. The stuck plant limits how far $V_f$ can go: $V_f=10^8$ is unstable there, $|z|=1.0015$.

\paragraph{Cost of $V_f=10^6$ on the Dahl plant.} None appeared in these tests. On the doublets, error and current stay the same, and the median settling time falls from 0.40 to 0.20\,s. Over a full orbit the error goes from 1.67 to 1.70\,mdeg, and the effect of the noise falls slightly, from 11.0 to 9.4\,$\mu$rad at 1\,$\mu$rad (Table~\ref{tab:twydahl}). These runs concern the elevation axis alone, with one seed and $V_f$ raised at this mode only. They do not replace a full tuning of $V$.

\begin{table}[H]\centering\footnotesize\setlength{\tabcolsep}{3pt}
\caption{Elevation on the Dahl plant, $Q_f=10^4$. Sines: encoder over torque [dB rad/(N\,m)], moving on the circle or at rest. Line of sight [$\mu$rad] at rest, 0.002\,N\,m, 57.8\,Hz. Doublets: rms error [mdeg], largest current [A], median settling [s]. Full orbit: error of the line of sight, rms [mdeg], and effect of noise of 1 and 10\,$\mu$rad on it, rms [$\mu$rad].}\label{tab:twydahl}
\resizebox{\textwidth}{!}{\begin{tabular}{lrrrrrrrrrrrr}\toprule
& \multicolumn{2}{c}{moving, 57\,Hz} & \multicolumn{2}{c}{rest, 57.8\,Hz} & rest, 58.5\,Hz & LOS & \multicolumn{3}{c}{doublets} & \multicolumn{3}{c}{full orbit} \\
\cmidrule(lr){2-3}\cmidrule(lr){4-5}\cmidrule(lr){8-10}\cmidrule(lr){11-13}
& 0.01\,N\,m & 0.2\,N\,m & 0.002\,N\,m & 0.2\,N\,m & 0.002\,N\,m & & err & $i$ & $t_\mathrm{set}$ & err & 1 & 10 \\ \midrule
cascade & -80.3 & -82.5 & -78.2 & -81.6 & -75.9 & 0.038 & 355.1 & 5.2 & 1.06 & 11.27 & 0.2 & 8.6 \\
tuned, $V_f=10^4$ & -81.0 & -79.4 & -62.5 & -79.7 & -79.7 & 0.232 & 50.2 & 9.0 & 0.40 & 1.67 & 11.0 & 25.4 \\
$V_f=10^6$ & -83.2 & -85.2 & -82.3 & -84.4 & -80.0 & 0.024 & 50.2 & 9.0 & 0.20 & 1.70 & 9.4 & 23.2
\\ \bottomrule
\end{tabular}}\end{table}

\paragraph{At the line of sight.} All this is read at the encoder, which in this mode moves much more than the line of sight: with 0.2\,N\,m at 57\,Hz on the moving axis, 15 against 1.6\,$\mu$rad with the cascade, 21 against 2.3 with the tuned controller. The amplification with stuck bearings also reaches the line of sight, but remains small in absolute terms at these amplitudes: 0.23\,$\mu$rad against 0.04 of the cascade.

\paragraph{The modes near 98\,Hz.} They have $q_f=0.1$, so the regulator and estimator poles stay at open-loop damping, $\zeta=0.032$ and 0.031 against 0.031. $F_y$ again has zeros on them, $|L|\le0.08$ and $T_{wy}$ stays at $-91.7$\,dB against $-92.3$\,dB without position loop. Here the weights also have little to work with: the height of the mode in $P_u$ is 26\,dB below that of the 56\,Hz mode. $Q_f=V_f=10^6$ gains only 1.3\,dB, with $\zeta=0.085$ and 0.039. $Q_f=10^6$, $V_f=10^{10}$ is unstable with stuck bearings ($|z|=1.09$) and raises the peak at 56\,Hz to $-57.6$\,dB.

\subsection{Why $T_{wy}$ in azimuth equals the cascade}
$T_{wy}=P_w/(1+L)$, so it depends on the feedback loop $L$ only, while the tracking also depends on the path of the command. Table~\ref{tab:loop} and Figure~\ref{fig:loop} decompose $L$. In azimuth $k_p$ alone would give 12.5\,dB at 1\,Hz. $K_f$ takes about 6\,dB, since the regulator feeds back the states of the velocity loop, with the largest gain on the pole at 2.6\,Hz, which lies inside the bandwidth. The estimator takes about 4\,dB: with $Q=V$, $K_{e,y}=25$\,rad/s, so at this $q_p$ the estimate partly follows the model and not the encoder. What remains is 2.6\,dB, equal to the cascade. The equality is a coincidence of this $q_p$, not a property: at $q_p=10^5$ the loop has 16.5\,dB. In elevation, with $K_{e,y}=573$\,rad/s, the estimator takes nothing and $K_f$ about 15\,dB.

\begin{table}[H]\centering\footnotesize\setlength{\tabcolsep}{4pt}
\caption{$|L|$ [dB] on the plant with bearings sliding, tuned weights. Fast estimator: $V=10^6Q$.}\label{tab:loop}
\begin{tabular}{llrrrrrrr}\toprule
axis & loop & 0.1\,Hz & 0.3\,Hz & 1\,Hz & 3\,Hz & 10\,Hz & 30\,Hz & 56\,Hz \\ \midrule
azimuth & cascade & 22.1 & 12.9 & 2.6 & $-$10.0 & $-$34.4 & $-$44.4 & $-$44.6 \\
 & LQG-PI & 22.7 & 13.2 & 2.6 & $-$7.6 & $-$32.7 & $-$51.6 & $-$57.1 \\
 & $K_f=0$ & 26.1 & 16.7 & 6.6 & $-$3.1 & $-$32.3 & $-$51.6 & $-$57.3 \\
 & fast estimator & 26.2 & 16.7 & 6.1 & $-$4.3 & $-$24.6 & $-$34.3 & $-$34.3 \\
 & fast estimator, $K_f=0$: $k_p$ alone & 31.9 & 22.8 & 12.5 & 0.0 & $-$24.3 & $-$34.3 & $-$34.5 \\
elevation & cascade & 22.1 & 12.8 & 3.8 & $-$8.3 & $-$29.2 & $-$47.6 & $-$39.7 \\
 & LQG-PI & 40.2 & 30.6 & 20.2 & 10.6 & 0.1 & $-$9.8 & $-$8.7 \\
 & $K_f=0$ & 54.5 & 45.0 & 35.2 & 25.9 & 10.5 & $-$7.0 & 4.7 \\
 & fast estimator & 40.6 & 31.1 & 20.6 & 11.1 & 0.5 & $-$9.6 & $-$9.2 \\
 & fast estimator, $K_f=0$: $k_p$ alone & 57.0 & 47.5 & 37.9 & 28.3 & 10.7 & $-$7.2 & 3.5
\\ \bottomrule
\end{tabular}\end{table}

\begin{figure}[H]\centering\includegraphics[width=\linewidth]{figures/cc_loop_decomposition.pdf}\caption{Decomposition of the loop gain.}\label{fig:loop}\end{figure}

\subsection{When the LQG-PI pays}\label{sec:benefits}
The cascade is a position loop of low gain, so comparing with it alone does not separate what comes from the LQG from what comes from a higher $k_p$. Tables~\ref{tab:benefitscmd} and~\ref{tab:benefitsdist} therefore add a proportional loop on the encoder with raised gain, on each axis alone, on the Dahl plant. Its stability limit with bearings sliding is $k\approx134$ in azimuth and 179 in elevation.

\begin{table}[H]\centering\footnotesize\setlength{\tabcolsep}{4pt}
\caption{Feedback and command. Phase margin, band and peak of $T_{ry}$ with bearings sliding. Doublets: rms error [mdeg], largest current [A], largest overshoot [mdeg], median settling [s]. Circle, one orbit: error of the line of sight on that axis, rms [mdeg].}\label{tab:benefitscmd}
\begin{tabular}{llrrrrrrrr}\toprule
axis & case & PM [$^\circ$] & band [Hz] & $T_{ry}$ peak [dB] & err & $i$ & over & $t_\mathrm{set}$ & circle \\ \midrule
azimuth & cascade & 63 & 2.1 & -0.0 & 372 & 11.7 & 43 & 1.27 & 104.3 \\
 & P, $k=$12 & 53 & 3.0 & 1.0 & 284 & 16.0 & 154 & 1.76 & 68.3 \\
 & P, $k=$25 & 35 & 4.9 & 4.4 & 168 & 28.9 & 238 & 0.75 & 32.7 \\
 & P, $k=$60 & 16 & 7.4 & 11.1 & 102 & 66.8 & 278 & 0.67 & 13.6 \\
 & LQG-PI tuned & 67 & 3.1 & 0.0 & 253 & 17.4 & 38 & 2.60 & 64.2 \\
 & LQG-PI, $q_p=10^5$ & 63 & 9.7 & 0.0 & 71 & 46.2 & 11 & 0.27 & 17.3 \\
elevation & cascade & 61 & 2.7 & 0.1 & 355 & 5.2 & 11 & 1.06 & 11.3 \\
 & P, $k=$12 & 44 & 3.5 & 2.6 & 269 & 6.5 & 116 & 0.66 & 7.8 \\
 & P, $k=$60 & 11 & 7.9 & 14.2 & 102 & 19.7 & 301 & 0.41 & 1.8 \\
 & P, $k=$100 & 6 & 10.2 & 20.2 & 69 & 22.8 & 222 & 0.43 & 1.1 \\
 & LQG-PI tuned & 62 & 22.1 & 0.0 & 50 & 9.0 & 8 & 0.40 & 1.7 \\
 & LQG-PI, $q_p=10^5$ & 62 & 12.4 & 0.0 & 78 & 9.2 & 24 & 1.37 & 2.5 \\
 & LQG-PI, $V_f=10^6$ & 61 & 22.1 & 0.0 & 50 & 9.0 & 7 & 0.20 & 1.7
\\ \bottomrule
\end{tabular}\end{table}

\begin{table}[H]\centering\footnotesize\setlength{\tabcolsep}{4pt}
\caption{Disturbances on the line of sight [$\mu$rad], difference from the same controller without disturbance. Over one orbit, rms: noise of 1\,$\mu$rad, torque step of 0.5\,N\,m, process noise, sine of 0.2\,N\,m at 1\,Hz. Over 12\,s, amplitude at the frequency: sines of 0.2 and 0.8\,N\,m.}\label{tab:benefitsdist}
\begin{tabular}{llrrrrrrr}\toprule
axis & case & noise & step & process & 1\,Hz & 0.2\,N\,m, 6.6\,Hz & 0.8\,N\,m, 3\,Hz & 0.8\,N\,m, 6.6\,Hz \\ \midrule
azimuth & cascade & 0.1 & 120.9 & 53.1 & 376.3 & 59 & 807 & 237 \\
 & P, $k=$12 & 0.1 & 87.6 & 44.3 & 275.7 & 63 & 896 & 250 \\
 & P, $k=$25 & 0.1 & 52.1 & 34.4 & 138.4 & 76 & 947 & 298 \\
 & P, $k=$60 & 0.3 & 34.9 & 30.7 & 61.2 & 141 & 465 & 546 \\
 & LQG-PI tuned & 0.1 & 115.7 & 51.9 & 365.4 & 62 & 844 & 246 \\
 & LQG-PI, $q_p=10^5$ & 0.3 & 38.2 & 17.9 & 87.0 & 66 & 352 & 250 \\
elevation & cascade & 0.2 & 51.2 & 147.8 & 396.5 & 35 & 1281 & 380 \\
 & P, $k=$12 & 0.2 & 35.8 & 114.7 & 350.4 & 35 & 1558 & 397 \\
 & P, $k=$60 & 1.1 & 14.8 & 37.4 & 91.5 & 35 & 1224 & 775 \\
 & P, $k=$100 & 2.8 & 8.5 & 26.8 & 58.1 & 35 & 750 & 3425 \\
 & LQG-PI tuned & 11.0 & 11.1 & 30.0 & 100.1 & 33 & 323 & 138 \\
 & LQG-PI, $q_p=10^5$ & 5.5 & 23.4 & 45.6 & 143.5 & 32 & 525 & 215 \\
 & LQG-PI, $V_f=10^6$ & 9.4 & 14.6 & 30.5 & 102.0 & 33 & 331 & 140
\\ \bottomrule
\end{tabular}\end{table}

\paragraph{The benefit: gain at equal margin.} With the proportional loop the margin falls quickly with the gain: at $k=12$ it is already $53^\circ$ in azimuth and $44^\circ$ in elevation, and at a doublet error comparable with the LQG-PI it is $16^\circ$ in azimuth ($k=60$) and $6^\circ$ in elevation ($k=100$). The LQG-PI keeps about the margin of the cascade, a flat $T_{ry}$ and a band of 22\,Hz in elevation and, with $q_p=10^5$, 9.7\,Hz in azimuth. This comes from $K_f$. With $K_f=0$ the same $k_p$ has $20.5^\circ$ in azimuth ($q_p=10^5$) and $18^\circ$ in elevation (tuned). The feedback of the velocity-loop poles, below 10\,Hz, gives most of it: alone it restores $58^\circ$ in azimuth and $62^\circ$ in elevation at $q_p=10^5$. At the tuned gain in elevation, $k_p=799$, it is not enough: alone it gives $27^\circ$, with a peak of $T_{ry}$ of 7\,dB at 55\,Hz, and the feedback of the flexible modes is needed as well to reach $62^\circ$.

\paragraph{Where it shows.}
\begin{itemize}
\item \textbf{Commands with a wide spectrum.} On the doublets in elevation the tuned LQG-PI has 50\,mdeg, 9\,A and 8\,mdeg of overshoot, against 69\,mdeg, 23\,A and 222\,mdeg with $k=100$. In azimuth $q_p=10^5$ has 71\,mdeg, 46\,A and 11\,mdeg, against 102\,mdeg, 67\,A and 278\,mdeg with $k=60$.
\item \textbf{Torques above the friction, at 3--7\,Hz.} With 0.8\,N\,m at 6.6\,Hz in elevation the line of sight moves by 138\,$\mu$rad, against 380 of the cascade and 3425 with $k=100$. In azimuth at 3\,Hz $q_p=10^5$ gives 352\,$\mu$rad against 807 of the cascade and 465 with $k=60$, but at 6.6\,Hz it only equals the cascade.
\item \textbf{Tracking and torque disturbances, once a margin is required.} With $62^\circ$ the LQG-PI gives in elevation 1.7\,mdeg on the circle against 7.8 with $k=12$, which has $44^\circ$, and 30 against 115\,$\mu$rad on the process noise.
\end{itemize}

\paragraph{Where it does not.}
\begin{itemize}
\item \textbf{Smooth tracking and slow disturbances, if the margin is not a constraint.} On the circle and at 1\,Hz only the gain at low frequency counts, and the proportional loop with high gain does as well or better: 1.1 against 1.7\,mdeg, 58 against 100\,$\mu$rad in elevation.
\item \textbf{Torques below the friction.} With 0.2\,N\,m in elevation at 3--8\,Hz all cases give about the same amplitude, within 20\,\% (32 to 35\,$\mu$rad at 6.6\,Hz), since the friction holds the axis and the feedback has little left to do.
\item \textbf{Sensor noise.} With 1\,$\mu$rad the tuned controller is the worst in elevation, 11\,$\mu$rad against 0.2 of the cascade and 2.8 with $k=100$.
\item \textbf{At the modes with $Q=V$} (Section~\ref{sec:twymodes}), and where the plant differs from the internal model: the margin in elevation with stuck bearings is $7^\circ$ at 57.8\,Hz, $22^\circ$ with $V_f=10^6$.
\item \textbf{Azimuth with the tuned weights.} It is close to $k=12$, and about 30\,\% worse on slow torques. The gain needs $q_p$ raised, which costs 46\,A without a current limit, and loses stability above about $60^\circ$ of elevation.
\end{itemize}

\paragraph{Limits of this comparison.} One axis, one seed, no current limit. The alternatives that also add phase or shape the command, such as a lead on the position loop, a prefilter or the feedforward excluded by D7, were not tested, and could reduce the benefit on the commands.

\subsection{Wind gusts}\label{sec:wind}
The book takes the servo error while tracking in wind gusts as the primary measure of performance (Section~8.3) and measures a ratio PI/LQG of 4.9 in azimuth and 6.5 in elevation on a 34-m antenna (Table~8.1); in simulation on the LMT the ratio is about 2 (Section~8.5.4). The argument is $T_{wy}=G(1-T_{ry})$: where the band covers the spectrum of the wind, the error is small, and the bulk of the improvement comes from the band (Section~8.3.3). The scenarios above had no wind: their random torque is flat from 0.1 to 50\,Hz. Table~\ref{tab:wind} therefore applies gusts with the Davenport spectrum (A14), of which 87\,\% of the energy is below 0.1\,Hz and 97\,\% below 1\,Hz.

\begin{table}[H]\centering\footnotesize\setlength{\tabcolsep}{4pt}
\caption{Wind gusts on the Dahl plant while tracking the circle, 120\,s, mean of two seeds. Effect on the line of sight and on the encoder, rms [$\mu$rad], and ratio cascade/case at the line of sight. Largest current [A].}\label{tab:wind}
\begin{tabular}{llrrrrrrr}\toprule
& & & \multicolumn{3}{c}{0.1\,N\,m rms} & \multicolumn{3}{c}{0.5\,N\,m rms} \\
\cmidrule(lr){4-6}\cmidrule(lr){7-9}
axis & case & PM [$^\circ$] & LOS & encoder & ratio & LOS & ratio & $i$ \\ \midrule
azimuth & cascade & 63 & 67.1 & 65.0 & 1.0 & 334.6 & 1.0 & 6.5 \\
 & P, $k=$12 & 53 & 48.9 & 46.6 & 1.4 & 244.2 & 1.4 & 6.5 \\
 & P, $k=$25 & 35 & 27.7 & 25.0 & 2.4 & 139.3 & 2.4 & 6.6 \\
 & P, $k=$60 & 16 & 15.8 & 12.2 & 4.2 & 78.8 & 4.2 & 7.0 \\
 & LQG-PI tuned & 67 & 65.6 & 63.6 & 1.0 & 327.0 & 1.0 & 6.5 \\
 & LQG-PI, $q_p=10^5$ & 63 & 16.6 & 13.7 & 4.0 & 83.3 & 4.0 & 6.9 \\
 & LQG-PI, $k_i=0.5k_p$ & -- & 63.2 & 61.5 & 1.1 & 316.1 & 1.1 & 6.5 \\
elevation & cascade & 61 & 152.4 & 152.4 & 1.0 & 485.8 & 1.0 & 5.2 \\
 & P, $k=$12 & 44 & 118.7 & 118.7 & 1.3 & 387.3 & 1.3 & 5.3 \\
 & P, $k=$60 & 11 & 31.2 & 31.2 & 4.9 & 93.2 & 5.2 & 5.4 \\
 & P, $k=$100 & 6 & 20.2 & 20.3 & 7.5 & 57.4 & 8.5 & 5.4 \\
 & LQG-PI tuned & 62 & 25.3 & 25.3 & 6.0 & 79.3 & 6.1 & 5.6 \\
 & LQG-PI, $q_p=10^5$ & 62 & 43.7 & 43.7 & 3.5 & 127.4 & 3.8 & 5.4 \\
 & LQG-PI, $V_f=10^6$ & 61 & 25.5 & 25.5 & 6.0 & 81.4 & 6.0 & 5.6 \\
 & LQG-PI, $k_i=0.5k_p$ & -- & 26.3 & 26.3 & 5.8 & 81.2 & 6.0 & 5.6
\\ \bottomrule
\end{tabular}\end{table}

\begin{itemize}
\item \textbf{The benefit of the book does appear.} Against the cascade the ratio is 6.0 in elevation with the tuned controller and 4.0 in azimuth with $q_p=10^5$, the same order as Table~8.1 of the book. This is a comparison of simulated ratios with measured ones on another instrument, so it only shows the order of magnitude.
\item \textbf{It does not appear in azimuth with the tuned weights}, ratio 1.0, because $|L|$ at low frequency equals that of the cascade (Table~\ref{tab:loop}). The reason is the tuning, which stopped $q_p$ at 646 on the current of the doublets, while $q_p=10^5$ asks for 46\,A there. In wind the same controller asks for no more than the cascade: 6.9 against 6.5\,A. This is the limit of Section~8.3.5 of the book, the high-gain controller limited by the drive under commands.
\item \textbf{The cause is the gain.} At equal gain the proportional loop does the same in wind, 15.8 against 16.6\,$\mu$rad in azimuth at $k=60$, but with $16^\circ$ of margin against $63^\circ$. With a margin of about $50^\circ$, $k=12$ gives only 1.4 in azimuth and 1.3 in elevation. In the book the band is limited by the fundamental frequency of the structure (Section~8.3.3), 1.7\,Hz on the LMT; here the modes are above 13\,Hz and the limit comes from the velocity loop.
\item \textbf{The integral gain does not matter here.} With $k_i=0.5k_p$, the typical value of the DSN, the ratio stays 1.1 in azimuth and 5.8 in elevation, although the error on the circle falls in azimuth, from 64 to 27\,mdeg.
\end{itemize}

The gusts act at the bearing port, as model 2 of the book; the real wind acts on the structure, and the encoder only partly represents the line of sight, as the book also notes (Section~8.5.4).

\subsection{Cross-coupling between the axes}\label{sec:coupling}
The Lagrange model of the project contains the cross terms, and they were lost in two steps on the way to the Dahl bench: the chains of the two axes are extracted from the model linearised at $\theta=0$, $\dot\psi=0$, where the cross terms in $C$ and $K$ are zero, and each axis is then fitted separately. For academic consistency the simulator of this section now restores the coupling of the telescope: the two axes are simulated together and exchange it at every step of 0.1\,ms.

With the angular velocity of the telescope $\bm\omega=(\dot\psi\sin\theta,\ -\dot\theta,\ \dot\psi\cos\theta)$ and its inertia tensor $I$ about the centre of rotation, the kinetic energy is
\begin{equation}
T=\tfrac12J(\theta)\dot\psi^2+\tfrac12I_{yy}\dot\theta^2+H(\theta)\dot\psi\dot\theta,\quad J=I_{xx}\sin^2\theta+I_{zz}\cos^2\theta+2I_{xz}\sin\theta\cos\theta,\quad H=-(I_{xy}\sin\theta+I_{yz}\cos\theta),
\end{equation}
that is, a mass matrix with $J(\theta)$ and $I_{yy}$ on the diagonal and $H(\theta)$ off it. The Lagrange equations then give three additions to the bench:
\begin{align}
\text{Coriolis and gyroscopic torques:}\quad&\tau_\psi=-J'(\theta)\,\dot\psi\dot\theta-H'(\theta)\,\dot\theta^2,\qquad\tau_\theta=\tfrac12J'(\theta)\,\dot\psi^2,\\
\text{inertial cross terms:}\quad&\tau_\psi=-H(\theta)\,\ddot\theta,\qquad\tau_\theta=-H(\theta)\,\ddot\psi,\\
\text{azimuth inertia of the telescope:}\quad&J(0)\ \longrightarrow\ J(\theta).
\end{align}
All of them act at the nodes that carry the telescope and use the velocities of those nodes, together with the absolute elevation. The dependence on $\theta$ concerns only the telescope, the rest of the azimuth load being fixed. It changes a single row of the state matrix, the acceleration of that node, and the discrete model is computed again every 0.1$^\circ$ of elevation. The inertial terms use the acceleration of the other axis at the previous step; with $H^2/(JI_{yy})\sim10^{-12}$ this delay has no effect, and on a test with $H$ raised up to 0.5\,kg\,m$^2$ against the exact mass matrix the error is about 0.1\,\% of the effect of $H$. As the bench applies a fitted factor to $J(0)=I_{zz}$, 1.135, the same factor is applied to the whole tensor, so that energy and torques stay consistent.

Three things were verified. The torques, $J(\theta)$ and $H(\theta)$ coincide with the symbolic model of the project, within $10^{-15}$\,N\,m and $10^{-16}$\,kg\,m$^2$ (with unit factor). With all these terms off, the coupled simulation coincides with the two separate axes, within $10^{-8}$\,$\mu$rad. All the results of this section are computed with these terms on.

Table~\ref{tab:coupled} separates their effect. On the test circle everything is negligible: below $0.06\,\mu$rad at the peak. Since it is the low elevation and slow rate of the circle that make it so, a much faster and higher circle is also run, with $\dot\psi$ up to 1\,rad/s at about $43^\circ$. There the largest contribution by far is $J(\theta)$, about 350\,$\mu$rad rms and 1.7\,mrad at the peak, on both controllers alike, since the inertia is lower by up to about 20\,\% and the azimuth feedback of the two is similar (Section~\ref{sec:circle}). It also lowers the largest azimuth current from 11 to 8.5\,A. The Coriolis and gyroscopic torques reach 0.19\,N\,m, about half the Coulomb torque of elevation (0.41\,N\,m), and their effect in elevation is 7 times smaller with the LQG-PI. $H$ is negligible on both trajectories, at $10^{-3}\,\mu$rad or below.

\begin{table}[H]\centering\footnotesize\setlength{\tabcolsep}{2.5pt}
\caption{Effect on the line of sight, each term added to the previous one, same controller. Largest torque [N\,m]; rms effect [$\mu$rad] of the torques on azimuth and elevation, of $J(\theta)$ on azimuth, of $H$ on both; peak effect of all terms on azimuth [$\mu$rad]; largest azimuth current [A] with fixed $J$ and with $J(\theta)$. Fast, high circle: centre at 15\,m, radius 10\,m, height 15\,m, 0.5\,rad/s.}\label{tab:coupled}
\resizebox{\textwidth}{!}{\begin{tabular}{lrrlrrrrrrrrrr}\toprule
trajectory & $\dot\psi_\mathrm{max}$ & $\theta$ [$^\circ$] & controller & $\tau_\psi$ & $\tau_\theta$ & az & el & $J$, az & $H$, az & $H$, el & peak & $i$, fixed & $i$ \\ \midrule
test circle & 0.0222 & 3 & cascade & $5.7\cdot10^{-6}$ & $2.3\cdot10^{-5}$ & $1.1\cdot10^{-3}$ & $3.8\cdot10^{-3}$ & $7\cdot10^{-3}$ & $3.5\cdot10^{-6}$ & $4.3\cdot10^{-6}$ & 0.051 & 1.1 & 1.1 \\
 &  &  & LQG-PI tuned & $2.2\cdot10^{-6}$ & $2.3\cdot10^{-5}$ & $4.6\cdot10^{-4}$ & $3.7\cdot10^{-4}$ & $6.7\cdot10^{-3}$ & $9.1\cdot10^{-7}$ & $6.8\cdot10^{-7}$ & 0.049 & 1.1 & 1.1 \\
fast, high circle & 1 & 43 & cascade & 0.0541 & 0.194 & 36 & 75.1 & 353 & $1.7\cdot10^{-4}$ & $6.3\cdot10^{-4}$ & $1.7\cdot10^{3}$ & 11.1 & 8.5 \\
 &  &  & LQG-PI tuned & 0.0604 & 0.196 & 36.3 & 10.2 & 355 & $1.6\cdot10^{-4}$ & $9\cdot10^{-5}$ & $1.7\cdot10^{3}$ & 11.2 & 8.6
\\ \bottomrule
\end{tabular}}\end{table}

\paragraph{Stability with the elevation.} The controllers keep the inertia of $\theta=0$ in their models and gains, so at high elevation the plant differs from them. Table~\ref{tab:stabtheta} checks the loop in azimuth up to $80^\circ$, where the node of the telescope has 64\,\% of its inertia. The cascade and the tuned LQG-PI stay stable. The variant with $q_p=10^5$ loses stability with stuck bearings from $60^\circ$, and also sliding at $80^\circ$, so it would need a gain schedule at high elevation.

\begin{table}[H]\centering\footnotesize\setlength{\tabcolsep}{4pt}
\caption{Azimuth, closed loop with $J(\theta)$ and fixed controller. Largest $|z|$ if unstable. Ratio: inertia of the node of the telescope against $\theta=0$.}\label{tab:stabtheta}
\begin{tabular}{rrllllll}\toprule
& & \multicolumn{2}{c}{cascade} & \multicolumn{2}{c}{tuned LQG-PI} & \multicolumn{2}{c}{$q_p=10^5$} \\
\cmidrule(lr){3-4}\cmidrule(lr){5-6}\cmidrule(lr){7-8}
$\theta$ [$^\circ$] & ratio & sliding & stuck & sliding & stuck & sliding & stuck \\ \midrule
0 & 1.000 & stable & stable & stable & stable & stable & stable \\
3 & 0.999 & stable & stable & stable & stable & stable & stable \\
30 & 0.905 & stable & stable & stable & stable & stable & stable \\
45 & 0.813 & stable & stable & stable & stable & stable & stable \\
60 & 0.723 & stable & stable & stable & stable & stable & \textbf{1.0024} \\
80 & 0.645 & stable & stable & stable & stable & \textbf{1.0016} & \textbf{1.0051}
\\ \bottomrule
\end{tabular}\end{table}

Three limits remain. The modal damping stays that of $\theta=0$. The inertia tensor is the CAD one with a single fitted factor. The tuning, on doublets of one axis at a time at $\theta\approx0$, and the single-axis tests of the previous subsections are unchanged.

The coupled structural model of the project gives the same indication in frequency at the point of the circle: the cross term is 60 to 125\,dB below the direct one (Table~\ref{tab:coupling}).

\textbf{This does not show that the real axes are uncoupled.} These terms only contain the coupling that comes from the kinetic energy of the assumed geometry. The V2 doublets (\path{results_V2_coupling}) show more: moving azimuth deviates the held elevation encoder by far more than the 0.2\,mdeg predicted, with a sign that always follows the torque held before the move, and the currents of the two axes have a coherence of 0.82--0.85 at 56.9\,Hz, with a ratio near 0.07. Neither effect is in the Lagrange model or in the bench, not even with these terms. This is a limit on the validity of the model, used here only to judge it, and not as a reference for the controllers. The deviation of elevation is comparable with the elevation error on the circle, so that error may be optimistic.

\begin{table}[H]\centering\footnotesize
\caption{Coupled structural model at the point of the circle: torque to encoder [dB rad/(N\,m)].}\label{tab:coupling}
\begin{tabular}{rrrrr}\toprule
$f$ [Hz] & az to az & el to el & el to az & az to el \\ \midrule
0.3 & $-$17.7 & $-$9.4 & $-$78.1 & $-$78.1 \\
1 & $-$38.7 & $-$31.0 & $-$110.1 & $-$110.1 \\
10 & $-$88.1 & $-$71.3 & $-$170.0 & $-$170.0 \\
56 & $-$82.4 & $-$122.9 & $-$207.5 & $-$207.6 \\
100 & $-$82.5 & $-$103.7 & $-$225.8 & $-$226.3
\\ \bottomrule
\end{tabular}\end{table}

\paragraph{Limits of this test.} Measured coupling between the axes not modelled, one seed, assumed disturbance levels (A13), no noise on the velocity feedback, no current limit, no feedforward on either scheme. Hilkert's study also showed that without feedforward the tracking error $v/K_{pp}$ dominates. The process noise first used in A13 was generated without discarding the edges of the filter at 0.1\,Hz, so it was not strictly Gaussian near the ends of a run (over 60\,s: excess kurtosis up to 20, peak 0.78\,N\,m). Table~\ref{tab:cc} now uses the corrected signal of Section~\ref{sec:thesis}, with and without the mean; Table~\ref{tab:benefitsdist} still uses the earlier signal.

\section{Pointing performance on a circle and on real trajectories}\label{sec:thesis}
This section compares the cascade tuned on the doublets with the tuned LQG-PI of Gawronski, on a circular trajectory and on two real trajectories recorded with the drone, without and with a disturbance torque. The result of interest is how often a laser beam of 8\,cm radius actually covers the drone.

\subsection{Definition of the pointing error}\label{sec:tp_def}
The pointing error follows Attorrese (2026), Section~5.3, where it is the coupling quantity of the pointer-drone problem. With the fixed position of the pointer $p^{[p]}$ and the \emph{actual} position of the drone $p^{[d]}$,
\begin{equation}
r=p^{[d]}-p^{[p]},\quad\hat r=\frac{r}{\sqrt{\lVert r\rVert^2+\varepsilon}},\quad
\hat n(\psi,\theta)=\begin{bmatrix}\cos\theta\cos\psi\\\cos\theta\sin\psi\\\sin\theta\end{bmatrix},\quad
\bm E=\hat n(\psi,\theta)\times\hat r,
\end{equation}
which are equations (5.21)--(5.23) of that work. Four remarks on how it is used here:
\begin{itemize}
\item Equation (5.23) defines a vector, and the cost of that work penalises its norm. The scalar reported here is $E=\lVert\bm E\rVert=\sin\alpha$, where $\alpha$ is the angle between the line of sight and the drone. At these errors $\sin\alpha$ and $\alpha$ coincide.
\item The regularisation $\varepsilon$ only prevents a singularity near the pointer. At 100\,m it has no effect, and it is set to zero.
\item In (5.22) $\theta$ is measured from the horizontal, although the text calls it polar and equation (5.3) of the same work measures it from the vertical. Here $\theta$ is the elevation, as in (5.22) and in the model.
\item The angles $(\psi,\theta)$ are those of the node that carries the telescope, not of the encoders, so structural deflections are included.
\end{itemize}
The same work uses a limit of $0.2^\circ$, or 3.49\,mrad, set by the company. The beam considered here is more demanding. Since $\lVert r\rVert E$ is exactly the distance of the drone from the line of sight, the \emph{miss distance} is $d=\lVert r\rVert E$, and the \emph{effective pointing} is the fraction of the evaluated time in which the drone lies inside the beam:
\begin{equation}
\eta=\frac{\text{time with }d\le R_b}{\text{evaluated time}},\qquad R_b=8\text{\,cm},
\end{equation}
which at 100\,m corresponds to about 0.8\,mrad. The drone is treated as a point.

\subsection{Simulation setup}\label{sec:tp_setup}
\paragraph{Plant.} The nonlinear model of the mount with Dahl bearings, the proxy of the real system in this study. The two axes are simulated together at 0.1\,ms and exchange the Coriolis, gyroscopic and inertial torques of the telescope, whose azimuth inertia depends on elevation (Section~\ref{sec:coupling}). There is no measurement noise and no current limit.

\paragraph{Controllers.} The present cascade, with $K_{pp}=7.92$ and 8.02\,1/s in azimuth and elevation, and the LQG-PI tuned in Section~\ref{sec:nb}, with $k_p=25.2$ and $k_i=0.49$ in azimuth, $k_p=799$ and $k_i=29.4$ in elevation. The gains of both stop at the current limit of the doublets, which is the application case that sizes the drive (decision D16). Neither has feedforward.

\paragraph{Trajectories.} The circle has radius 10\,m, height 5\,m and centre at 100\,m, at 0.2\,rad/s, so a speed of 2\,m/s. The real trajectories are \texttt{dataset\_7} and \texttt{dataset\_8}, recorded with the motors running, about 30 and 40\,m long at a mean speed of 0.17\,m/s. They are shown in Figure~\ref{fig:tp_traj}.

\begin{figure}[H]\centering\includegraphics[width=0.82\linewidth]{figures/tp_trajectories.pdf}\caption{Real trajectories seen from above, after alignment of the estimates to the ground truth. LUCY lies 100\,m away along $-x$, so the long side of each trajectory is seen as lateral motion.}\label{fig:tp_traj}\end{figure}

\paragraph{Reference and error.} The pointer follows the estimates, and the error is computed against the ground truth:
\begin{itemize}
\item \textbf{Estimators:} EKF2 at 100\,Hz and OpenVINS with \texttt{config\_v3} at 30\,Hz, first run of each.
\item \textbf{Ground truth:} the TagSLAM reference. Its folder \texttt{truths} has a frame attached to a tag, with $z$ horizontal. The same data rigidly rotated with $z$ up (\texttt{truths\_leveled}) are used, as in the evaluation of the estimators, since the elevation axis of the mount needs the vertical.
\item \textbf{Alignment:} each estimate is aligned to the ground truth with a single rigid transformation in position and yaw, computed on the whole window as in the evaluation of the estimators. This reproduces its ATE of 2.6 and 7.8\,cm for \texttt{config\_v3}. The transformation uses the future of the run, so it is the favourable case; Section~\ref{sec:tp_limits} gives the alternative.
\item \textbf{Position of LUCY:} 100\,m from the first ground-truth point of each trajectory, with the drone 5\,m above it, as on the circle, along the $x$ axis of the leveled frame.
\item \textbf{Command:} the azimuth and elevation of the last estimate, held between estimates, without added latency.
\item \textbf{Window:} the interval common to the two estimators, which excludes the initialisation of OpenVINS: 151\,s and 209\,s. The first 5\,s of each run are left out of all metrics.
\end{itemize}

\paragraph{Disturbance.} A torque $\tau_w$ acts at the bearing of each axis, independent between the axes, Gaussian with zero mean, 0.1\,N\,m rms and band 0.1--50\,Hz (Figure~\ref{fig:tp_dist}). On LUCY the source is not known in advance, so the level is a plausible order of magnitude and not a measured value. The peaks, about 0.4\,N\,m, are of the order of the Coulomb friction of the axes, 0.23 and 0.41\,N\,m. As a reference, 0.1\,N\,m is about 2\,N at an arm of 5\,cm, the force of a 5\,m/s wind on about 0.1\,m$^2$ with a drag coefficient of about 1.2. Three realisations are used, the same for every trajectory and controller. Each has zero sample mean by construction, and over 60\,s the excess kurtosis is below 0.1.

\paragraph{Static part of the disturbance.} A second case adds a mean of 0.2\,N\,m, with the same sign on both axes, to the same realisations. It stands for torques that the simulator does not contain and that do not average to zero, such as cable springs and imbalance. The value is plausible as an upper estimate:
\begin{itemize}
\item it is 0.78\,A at the nominal torque constant of 0.258\,N\,m/A, 4\,\% of the current limit of 18\,A;
\item the cable springs identified on the V2 sequences (\path{closedLoopModel/results_v2_identification}), 0.93 and 1.66\,N\,m/rad with neutral angles of $-2.7$ and $-2.2^\circ$, give 0.04--0.06\,N\,m near the measured positions. They reach 0.2\,N\,m only 12 and $7^\circ$ away from the neutral angle;
\item the identification found no imbalance of the telescope, with an offset of the centre of gravity of about 0.3\,mm. With the 151\,N of the telescope this is 0.05\,N\,m, and 0.2\,N\,m would take 1.3\,mm.
\end{itemize}
Errors in friction and in the dynamics do not belong to this term: they change with the velocity and with the frequency, and a constant does not represent them. The static torque already acts before tracking starts, so these runs start from the equilibrium in which the velocity integrator of the drive holds it. A torque that appears suddenly is the step of Section~\ref{sec:circle}.

\begin{figure}[H]\centering\includegraphics[width=0.9\linewidth]{figures/tp_disturbance.pdf}\caption{Disturbance torque on the two axes. Left: two seconds of one realisation. Right: distribution over 60\,s against the Gaussian with the same variance.}\label{fig:tp_dist}\end{figure}

\subsection{Results}\label{sec:tp_results}
Table~\ref{tab:tp_main} and Figure~\ref{fig:tp_eff} collect the results. The percentages of Figure~\ref{fig:tp_eff} are the main result of this section.

\begin{table}[H]\centering\footnotesize\setlength{\tabcolsep}{3pt}
\caption{Pointing error and effective pointing $\eta$. $E$: rms and largest value [mrad]. $d$: rms and 95th percentile of the miss distance [cm]. With disturbance: mean of three realisations, with the range of $\eta$; effect: rms difference of the line of sight from the same run without disturbance [$\mu$rad].}\label{tab:tp_main}
\resizebox{\textwidth}{!}{\begin{tabular}{lllrrrrrrrrr}\toprule
& & & \multicolumn{5}{c}{no disturbance} & \multicolumn{4}{c}{disturbance, $\sigma_\tau=0.1$\,N\,m} \\
\cmidrule(lr){4-8}\cmidrule(lr){9-12}
trajectory & estimator & controller & $E_\mathrm{rms}$ & $E_\mathrm{max}$ & $d_\mathrm{rms}$ & $d_{95}$ & $\eta$ [\%] & $E_\mathrm{rms}$ & $d_\mathrm{rms}$ & $\eta$ [\%] & effect \\ \midrule
circle & -- & cascade & 1.83 & 2.81 & 18.2 & 25.3 & \textbf{13.8} & 1.82 & 18.1 & \textbf{14.2} (13.6--14.5) & 162 \\
 &  & LQG-PI & 1.12 & 1.70 & 11.2 & 15.4 & \textbf{29.1} & 1.12 & 11.2 & \textbf{29.8} (29.3--30.2) & 72 \\
\texttt{dataset\_7} & EKF2 & cascade & 8.20 & 16.23 & 81.7 & 148.4 & \textbf{2.9} & 8.20 & 81.8 & \textbf{2.5} (2.5--2.6) & 284 \\
 &  & LQG-PI & 8.18 & 16.19 & 81.6 & 147.9 & \textbf{2.7} & 8.18 & 81.6 & \textbf{2.4} (2.3--2.5) & 112 \\
\texttt{dataset\_7} & \texttt{config\_v3} & cascade & 0.46 & 1.22 & 4.6 & 7.7 & \textbf{96.3} & 0.45 & 4.4 & \textbf{97.8} (97.4--98.1) & 194 \\
 &  & LQG-PI & 0.34 & 1.09 & 3.4 & 5.8 & \textbf{99.5} & 0.35 & 3.5 & \textbf{99.5} (99.4--99.6) & 91 \\
\texttt{dataset\_8} & EKF2 & cascade & 7.56 & 18.69 & 75.4 & 150.4 & \textbf{0.8} & 7.57 & 75.5 & \textbf{0.8} (0.7--0.9) & 338 \\
 &  & LQG-PI & 7.55 & 18.70 & 75.3 & 150.7 & \textbf{1.0} & 7.55 & 75.3 & \textbf{1.0} (0.9--1.0) & 117 \\
\texttt{dataset\_8} & \texttt{config\_v3} & cascade & 0.87 & 3.48 & 8.7 & 16.3 & \textbf{53.5} & 0.88 & 8.8 & \textbf{55.7} (55.1--56.0) & 226 \\
 &  & LQG-PI & 0.79 & 1.99 & 7.9 & 14.8 & \textbf{61.0} & 0.79 & 7.9 & \textbf{61.5} (61.3--61.7) & 101
\\ \bottomrule
\end{tabular}}\end{table}

\begin{figure}[H]\centering\includegraphics[width=0.92\linewidth]{figures/tp_effective_pointing.pdf}\caption{Effective pointing for a beam of 8\,cm radius. The horizontal line shows the limit given by the estimator: a pointer that follows the estimate with no error.}\label{fig:tp_eff}\end{figure}

\paragraph{Circle.} On the circle the reference is exact, so the whole error comes from the loop (Figure~\ref{fig:tp_circle}). The LQG-PI lowers $E_\mathrm{rms}$ from 1.83 to 1.12\,mrad and raises $\eta$ from 13.8 to 29.1\,\%. Both remain low because the error is almost entirely the velocity lag in azimuth. The azimuth rate of the line of sight reaches 22\,mrad/s, with 14\,mrad/s rms, against 1.8--1.9\,mrad/s rms on the real trajectories. A least-squares fit of the azimuth error on this rate gives $e_\psi \approx \dot\psi/K_v$, with $K_v = 7.8\,\mathrm{s^{-1}}$ for the cascade, practically its $K_{pp} = 7.92$ ($R^2 = 0.996$), and $K_v = 12.9\,\mathrm{s^{-1}}$ for the LQG-PI ($R^2 = 0.974$). The beam has a half-width of about 0.8\,mrad at 100\,m, so the line of sight stays inside only while $|\dot\psi| < 0.8\,\mathrm{mrad} \times K_v$, that is 6.3 and 10.3\,mrad/s. This happens only around the reversals of azimuth, for 20 and 35\,\% of the orbit. There the bearing friction adds its peak: the fit residual is about 250 and 290\,$\mu$rad rms, against 40 and 74\,$\mu$rad elsewhere, so $\eta$ drops to 13.8 and 29.1\,\%. Keeping the lag below 0.8\,mrad over the whole orbit would require $K_v \approx 28\,\mathrm{s^{-1}}$. With a beam radius of 16\,cm $\eta$ would be 43 and 100\,\%; with 35\,cm, the radius equivalent to $0.2^\circ$, it is 100\,\% for both. Both controllers stay within the limit of $0.2^\circ$, with largest values of 2.81 and 1.70\,mrad.

\begin{figure}[H]\centering\includegraphics[width=0.92\linewidth]{figures/tp_circle_miss.pdf}\caption{Circle: miss distance over one orbit, without disturbance and with one realisation of the disturbance.}\label{fig:tp_circle}\end{figure}

\paragraph{Real trajectories, no disturbance.} The estimator sets the limit (Figures~\ref{fig:tp_ds7} and~\ref{fig:tp_ds8}, Table~\ref{tab:tp_loop}).
\begin{itemize}
\item \textbf{EKF2.} Its error in position, 0.81--0.86\,m ATE on these windows, is about ten times the beam radius. $\eta$ stays at 0.8--2.9\,\% with either controller, against 1.1--2.8\,\% of the estimator alone, and the two controllers cannot be told apart. The largest values of $E$, 16--19\,mrad, exceed the limit of $0.2^\circ$.
\item \textbf{\texttt{config\_v3} on \texttt{dataset\_7}.} The estimator alone would give 100\,\%. The cascade reaches 96.3\,\% and the LQG-PI 99.5\,\%.
\item \textbf{\texttt{config\_v3} on \texttt{dataset\_8}.} The ATE of 7.8\,cm, mostly in height, is close to the beam radius, so the limit is 66.3\,\%. The cascade reaches 53.5\,\% and the LQG-PI 61.0\,\%. The largest value of $E$ falls from 3.48\,mrad, at the limit of $0.2^\circ$, to 1.99\,mrad. The height error of the estimate is a drift, not an offset: after the alignment it is about $-15$\,cm in the first 30\,s and $+8$\,cm in the last 40\,s, and the estimate leaves the beam almost only in these two intervals. It also stays between 5 and 8\,cm, near the edge of the beam, for 17\,\% of the time, against never on \texttt{dataset\_7}. There a small loop error is enough to leave the beam, and 80 to 90\,\% of the time lost to the loop falls in these instants. The loop error itself is close to that of \texttt{dataset\_7}, except in elevation with the cascade, 2.7 against 1.8\,cm rms, because the elevation moves more: 0.87 against 0.32\,mrad/s rms.
\end{itemize}

\begin{figure}[H]\centering\includegraphics[width=\linewidth]{figures/tp_dataset_7_miss.pdf}\caption{\texttt{dataset\_7}: miss distance with the two estimators, without and with one realisation of the disturbance. The grey line is the estimator alone. Note the different vertical scales.}\label{fig:tp_ds7}\end{figure}

\begin{figure}[H]\centering\includegraphics[width=\linewidth]{figures/tp_dataset_8_miss.pdf}\caption{\texttt{dataset\_8}: miss distance with the two estimators, without and with one realisation of the disturbance. The grey line is the estimator alone. Note the different vertical scales.}\label{fig:tp_ds8}\end{figure}

\begin{table}[H]\centering\footnotesize\setlength{\tabcolsep}{4pt}
\caption{Sources of error, no disturbance. Estimator alone: ATE on the window [cm], miss distance rms [cm] and $\eta$ [\%] of an ideal pointer. Loop: rms angle between the line of sight and the command [mrad], and its parts in azimuth and elevation. Largest current [A] in azimuth and elevation, without and with disturbance.}\label{tab:tp_loop}
\resizebox{\textwidth}{!}{\begin{tabular}{lllrrrrrrrrrr}\toprule
& & & \multicolumn{3}{c}{estimator alone} & \multicolumn{3}{c}{loop} & \multicolumn{2}{c}{$i_\mathrm{max}$, none} & \multicolumn{2}{c}{$i_\mathrm{max}$, disturbance} \\
\cmidrule(lr){4-6}\cmidrule(lr){7-9}\cmidrule(lr){10-11}\cmidrule(lr){12-13}
trajectory & estimator & controller & ATE & $d_\mathrm{rms}$ & $\eta$ & total & az & el & az & el & az & el \\ \midrule
circle & -- & cascade & -- & 0.0 & 100.0 & 1.83 & 1.82 & 0.20 & 1.1 & 1.6 & 1.7 & 1.5 \\
 &  & LQG-PI &  &  &  & 1.12 & 1.12 & 0.03 & 1.1 & 1.5 & 1.7 & 2.1 \\
\texttt{dataset\_7} & EKF2 & cascade & 85.5 & 81.9 & 2.8 & 0.48 & 0.22 & 0.43 & 1.1 & 1.9 & 1.6 & 2.0 \\
 &  & LQG-PI &  &  &  & 0.20 & 0.17 & 0.11 & 1.2 & 3.1 & 1.6 & 3.6 \\
\texttt{dataset\_7} & \texttt{config\_v3} & cascade & 2.6 & 2.1 & 100.0 & 0.33 & 0.28 & 0.18 & 1.1 & 1.6 & 1.7 & 1.6 \\
 &  & LQG-PI &  &  &  & 0.20 & 0.19 & 0.07 & 1.2 & 3.8 & 1.7 & 4.0 \\
\texttt{dataset\_8} & EKF2 & cascade & 80.7 & 75.5 & 1.1 & 0.66 & 0.26 & 0.61 & 1.2 & 1.9 & 1.6 & 2.1 \\
 &  & LQG-PI &  &  &  & 0.25 & 0.21 & 0.12 & 1.2 & 3.2 & 1.6 & 3.6 \\
\texttt{dataset\_8} & \texttt{config\_v3} & cascade & 7.8 & 7.4 & 66.3 & 0.40 & 0.30 & 0.27 & 1.2 & 1.8 & 1.7 & 1.9 \\
 &  & LQG-PI &  &  &  & 0.24 & 0.23 & 0.06 & 1.3 & 8.4 & 1.7 & 8.4
\\ \bottomrule
\end{tabular}}\end{table}

The loop adds 0.33--0.40\,mrad with the cascade and 0.20--0.24\,mrad with the LQG-PI on \texttt{config\_v3}. The difference comes mostly from elevation: 0.18--0.27 against 0.06--0.07\,mrad. The short peaks of the cascade in the figures, up to 35\,cm, coincide with reversals of the elevation motion. There the current has to reach the friction level, about 1.6\,A, and with the low gain of the cascade this takes an error of 1.5--2.4\,mrad. At the same instants the LQG-PI stays close to the estimator alone. The cost is on the current. The largest elevation current rises from about 1.6--1.9\,A to 3--8\,A, with the largest value at a jump of the estimate on \texttt{dataset\_8}, still far from the limit of 18\,A used on the doublets.

\begin{figure}[H]\centering\includegraphics[width=0.92\linewidth]{figures/tp_miss_cdf.pdf}\caption{Cumulative distribution of the miss distance, no disturbance. The value at the dashed line is the effective pointing.}\label{fig:tp_cdf}\end{figure}

\paragraph{Effect of the disturbance.} The disturbance displaces the line of sight by 157--339\,$\mu$rad with the cascade and 69--118\,$\mu$rad with the LQG-PI, so the LQG-PI rejects it 2.1 to 2.9 times better. Even so, $\eta$ hardly moves, and by at most 2.2 points in any case. At 100\,m, 100\,$\mu$rad is 1\,cm, which adds in quadrature to an error of several centimetres. With the cascade $\eta$ even rises slightly, from 96.3 to 97.8\,\% and from 53.5 to 55.7\,\% on \texttt{config\_v3}: the small torque keeps the bearing moving and reduces the peaks at the reversals, the dither effect already seen in Section~\ref{sec:circle}. The LQG-PI, which has smaller peaks, does not show it.

\paragraph{Effect of the static part.} Table~\ref{tab:tp_mean} repeats the runs with disturbance with the mean of 0.2\,N\,m. Against the same realisation with zero mean, the line of sight moves by a constant: $-12.7\,\mu$rad in azimuth and zero in elevation. The value is the same for every trajectory, controller and realisation, and its largest value equals its mean. The current rises by 0.86\,A in azimuth and 0.74\,A in elevation. $\eta$ changes by at most 0.4 points, and the rms miss distance by less than 0.03\,cm.

The reason is the velocity integrator of the drive, which both schemes have. It holds a constant torque with zero velocity error, and the rest of the loop does not depend on its state. The Dahl friction does not either, since it depends only on the velocity. What remains is the static deflection between the encoder, which the loop brings to the command, and the node of the telescope: 12.7\,$\mu$rad for 0.2\,N\,m in azimuth, that is 1.3\,mm at 100\,m. A torque of the same size appearing suddenly would give the response of the step of Section~\ref{sec:circle}, which peaks at about 1.1\,mrad with 0.5\,N\,m, so at about 0.45\,mrad with 0.2\,N\,m, and then decays.

\begin{table}[H]\centering\footnotesize\setlength{\tabcolsep}{3pt}
\caption{Static part of the disturbance: runs with $\sigma_\tau=0.1$\,N\,m and a mean of 0.2\,N\,m, against the same realisations with zero mean. $\eta$: mean of three realisations, with the range for the case with the mean; $\Delta\eta$ in points. Change of the line of sight from the run with zero mean: mean in azimuth and in elevation, and largest value of the total [$\mu$rad]. Increase of the mean current [A].}\label{tab:tp_mean}
\resizebox{\textwidth}{!}{\begin{tabular}{lllrrrrrrrr}\toprule
& & & \multicolumn{3}{c}{$\eta$ [\%]} & \multicolumn{3}{c}{change of the line of sight} & \multicolumn{2}{c}{current} \\
\cmidrule(lr){4-6}\cmidrule(lr){7-9}\cmidrule(lr){10-11}
trajectory & estimator & controller & zero mean & mean 0.2\,N\,m & $\Delta\eta$ & az & el & largest & az & el \\ \midrule
circle & -- & cascade & 14.2 & \textbf{14.2} (13.7--14.6) & 0.0 & -12.7 & 0.0 & 12.7 & 0.86 & 0.74 \\
 &  & LQG-PI & 29.8 & \textbf{29.4} (29.2--29.8) & -0.4 & -12.7 & 0.0 & 12.7 & 0.86 & 0.74 \\
\texttt{dataset\_7} & EKF2 & cascade & 2.5 & \textbf{2.6} (2.5--2.6) & 0.0 & -12.7 & 0.0 & 12.7 & 0.86 & 0.74 \\
 &  & LQG-PI & 2.4 & \textbf{2.4} (2.3--2.5) & 0.0 & -12.7 & 0.0 & 12.7 & 0.86 & 0.74 \\
\texttt{dataset\_7} & \texttt{config\_v3} & cascade & 97.8 & \textbf{98.1} (97.7--98.3) & +0.3 & -12.7 & 0.0 & 12.7 & 0.86 & 0.74 \\
 &  & LQG-PI & 99.5 & \textbf{99.5} (99.4--99.6) & 0.0 & -12.7 & 0.0 & 12.7 & 0.86 & 0.74 \\
\texttt{dataset\_8} & EKF2 & cascade & 0.8 & \textbf{0.8} (0.7--0.9) & 0.0 & -12.7 & 0.0 & 12.7 & 0.86 & 0.74 \\
 &  & LQG-PI & 1.0 & \textbf{1.0} (1.0--1.0) & 0.0 & -12.7 & 0.0 & 12.7 & 0.86 & 0.74 \\
\texttt{dataset\_8} & \texttt{config\_v3} & cascade & 55.7 & \textbf{55.7} (55.2--55.9) & 0.0 & -12.7 & 0.0 & 12.7 & 0.86 & 0.74 \\
 &  & LQG-PI & 61.5 & \textbf{61.5} (61.3--61.7) & 0.0 & -12.7 & 0.0 & 12.7 & 0.86 & 0.74
\\ \bottomrule
\end{tabular}}\end{table}

\paragraph{Higher speed.} The real trajectories are slow, 0.19\,m/s on average. To see whether faster motion separates the two controllers more, the same paths with \texttt{config\_v3} are run 2, 3, 5 and 10 times faster, up to 1.9\,m/s, the speed of the circle (Table~\ref{tab:tp_speed}, Figure~\ref{fig:tp_speed}). Truth and estimate are compressed in time. The estimator keeps its rate of 30\,Hz, so each estimate is held over a longer distance, and the error of the estimate is that of the original run. The transient left out scales with the speed; leaving out 2\,s more changes $\eta$ by at most 1.2 points.
\begin{itemize}
\item \textbf{The gap between the controllers grows.} At twice the speed $\eta$ on \texttt{dataset\_7} falls from 96.3 to 58.7\,\% with the cascade and from 99.5 to 83.9\,\% with the LQG-PI; on \texttt{dataset\_8} from 53.5 to 27.7 and from 61.0 to 47.0\,\%. The gap is largest at 2 to 3 times the speed, 17 to 25 points, against 3 to 8 at the recorded speed. Beyond that both controllers fall, and at 1.9\,m/s they are of the same order as on the circle.
\item \textbf{The loop error grows almost in proportion to the speed} and is mostly in azimuth: it is the velocity lag of the circle. The ratio LQG-PI/cascade stays between 0.57 and 0.65, so the gap in $\eta$ is largest where the error of the cascade has just crossed the beam half-angle and that of the LQG-PI has not.
\item \textbf{The cost is the current in elevation.} Each new estimate is a step of the command, up to 0.9\,mrad at five times the speed, and the stiff elevation of the LQG-PI answers with a current peak 2\,ms later. The peaks exceed the 18\,A of the drive on \texttt{dataset\_8} from three times the speed and on \texttt{dataset\_7} at ten times, for 0.01 to 0.06\,\% of the time, and 0.6\,\% on \texttt{dataset\_8} at ten times. The model has no current limit, so these points, hollow in the figure, are not reachable as computed. The cascade stays below 2.3\,A in elevation and 4.1\,A in azimuth.
\item \textbf{The estimator alone falls at 1.9\,m/s}, to 87.8 and 51.4\,\%, because an estimate held for 33\,ms lags the drone by up to 6\,cm.
\end{itemize}

\begin{figure}[H]\centering\includegraphics[width=0.92\linewidth]{figures/tp_speed.pdf}\caption{Same paths with \texttt{config\_v3} at higher speed, no disturbance: effective pointing and rms angle between the line of sight and the command. Hollow markers: the LQG-PI asks for more than the 18\,A of the drive, which the model does not limit.}\label{fig:tp_speed}\end{figure}

\begin{table}[H]\centering\footnotesize\setlength{\tabcolsep}{4pt}
\caption{Same paths with \texttt{config\_v3} at higher speed, no disturbance. Factor on the speed and mean speed [m/s]; $\eta$ [\%] of the estimator alone, the cascade and the LQG-PI, and their difference [points]; rms loop error [mrad]; largest elevation current [A], in bold when above 18\,A.}\label{tab:tp_speed}
\begin{tabular}{lrrrrrrrrrr}\toprule
& & & \multicolumn{4}{c}{$\eta$ [\%]} & \multicolumn{2}{c}{loop [mrad]} & \multicolumn{2}{c}{$i_\mathrm{max}$, el [A]} \\
\cmidrule(lr){4-7}\cmidrule(lr){8-9}\cmidrule(lr){10-11}
trajectory & factor & speed & estimator & cascade & LQG-PI & $\Delta\eta$ & cascade & LQG-PI & cascade & LQG-PI \\ \midrule
\texttt{dataset\_7} & 1 & 0.19 & 100.0 & 96.3 & 99.5 & +3.2 & 0.33 & 0.20 & 1.6 & 3.8 \\
 & 2 & 0.39 & 100.0 & 58.7 & 83.9 & +25.1 & 0.58 & 0.37 & 1.7 & 5.7 \\
 & 3 & 0.58 & 100.0 & 43.5 & 62.3 & +18.8 & 0.84 & 0.53 & 1.7 & 8.2 \\
 & 5 & 0.97 & 99.3 & 32.8 & 44.0 & +11.2 & 1.34 & 0.86 & 1.8 & 11.8 \\
 & 10 & 1.94 & 87.8 & 24.0 & 32.6 & +8.6 & 2.55 & 1.67 & 1.8 & \textbf{19.2} \\
\texttt{dataset\_8} & 1 & 0.19 & 66.3 & 53.5 & 61.0 & +7.5 & 0.40 & 0.24 & 1.8 & 8.4 \\
 & 2 & 0.38 & 65.2 & 27.7 & 47.0 & +19.3 & 0.69 & 0.39 & 2.0 & 14.5 \\
 & 3 & 0.57 & 63.9 & 15.1 & 32.3 & +17.2 & 0.93 & 0.53 & 2.0 & \textbf{21.1} \\
 & 5 & 0.95 & 61.3 & 12.1 & 21.4 & +9.3 & 1.38 & 0.80 & 2.2 & \textbf{33.4} \\
 & 10 & 1.91 & 51.4 & 6.4 & 19.4 & +13.0 & 2.33 & 1.47 & 2.3 & \textbf{57.7}
\\ \bottomrule
\end{tabular}\end{table}

\subsection{Discussion}\label{sec:tp_discussion}
\begin{itemize}
\item \textbf{The sources of error have a clear order: estimator, then velocity lag, then disturbance.} With EKF2 no controller can help. With \texttt{config\_v3} the control matters, because estimator and beam are of the same size.
\item \textbf{The LQG-PI brings the pointer close to the limit of the estimator.} It leaves 0.5\,\% of the time out of the beam against 3.7\,\% on \texttt{dataset\_7}, and covers 61.0 of the 66.3\,\% available on \texttt{dataset\_8}, against 53.5.
\item \textbf{Its better disturbance rejection does not change the coverage} at this level of torque. It would matter with a much better estimator or a much narrower beam.
\item \textbf{A static torque does not change the coverage with either controller.} The velocity integrator of the drive holds it, and only a deflection of 1.3\,mm at the drone remains. A constant offset is therefore not the way unmodelled torques would show up; position- and velocity-dependent terms, and torques that appear suddenly, are.
\item \textbf{The result on the circle shows the other limit.} At 2\,m/s the velocity lag dominates, and even the LQG-PI stays out of the beam for 70\,\% of the time, since the gains stop at the doublets and there is no feedforward.
\item \textbf{Speed decides how much the controller matters.} On the same paths the gap grows from 3--8 points at the recorded speed to 17--25 points at 0.4--0.6\,m/s, and then both controllers fall as the velocity lag exceeds the beam. At these speeds the stiff elevation of the LQG-PI would also need a smoother command than estimates held at 30\,Hz, or it exceeds the current of the drive.
\end{itemize}

\subsection{Limits}\label{sec:tp_limits}
\begin{itemize}
\item \textbf{Alignment of the estimates.} It uses the whole run, as in the evaluation of the estimators. Computing it from the first 5\,s only, as a real system would have to, changes the limit given by the estimator alone: from 66.4 to 14.0\,\% with \texttt{config\_v3} on \texttt{dataset\_8}, where the ATE grows to 17\,cm, while \texttt{dataset\_7} stays at 100\,\%. EKF2 stays low, 3.4 and 4.1\,\%. The registration between estimator and pointer is therefore as important as the estimator itself. Closed-loop runs with this alignment were not made.
\item \textbf{Latency.} Processing and communication delays are not modelled. At 0.17\,m/s, 50\,ms are about 1\,cm, but they would matter in faster motion.
\item \textbf{Geometry.} One position of LUCY. Placing it along $y$ leaves \texttt{config\_v3} unchanged (100 and 66.2\,\%) and gives 6.1 and 1.3\,\% with EKF2, for the estimator alone.
\item \textbf{Faster paths.} They keep the error of the estimate of the original run, while an estimator at higher speed would err differently, and they ignore the current limit of the drive.
\item \textbf{Trajectories.} Slow and at low elevation, about $3^\circ$. They do not exercise the azimuth inertia at high elevation, nor flight conditions.
\item \textbf{Ground truth.} TagSLAM has an uncertainty of about 1\,cm, which is not small against the ATE of 2.6\,cm.
\item \textbf{Model.} No measurement noise, no current limit, Dahl plant as the proxy of the real system. The coupling measured between the axes is not modelled (Section~\ref{sec:coupling}). The disturbance level is assumed, and three realisations give only an indication of its spread. The cable springs identified on V2 are not in the simulator, so the dependence of their torque on the position is not tested. The margin of the tuned LQG-PI with stuck bearings at 57.8\,Hz is low (Section~\ref{sec:twymodes}), and these runs do not exercise it specifically.
\item \textbf{Effective pointing.} It counts time with the centre of the drone inside the beam, and does not consider beam divergence, atmospheric effects or the size of the receiver.
\end{itemize}



\section{Open points}
\begin{itemize}
\item Choice of $q_i$: the doublets favour a low value, the circle a high one (Section~\ref{sec:circle}).
\item Elevation gain: $k_p\approx800$\,1/s depends on A3; its real benefit needs a test on the axis.
\item Real configuration of the PLC and of the velocity loop: $K_{pv}$, $K_{iv}$ are effective gains, and the plant of Step~0 inherits them.
\item Velocity-loop controller of the DSN: Section~8.5 of the book says proportional, Section~8.5.2 says PI.
\item Feedforward in the operating DSN: Chapter~8 does not state it; reference~[2] of Chapter~7 was not read.
\end{itemize}

\section{Points to study}
\begin{itemize}
\item \textbf{Harmonic linearisation in detail} (to study): what it means to replace a nonlinearity with its first harmonic, why the result depends on amplitude, and when it is valid. Reference texts: Gelb and Vander Velde, \emph{Multiple-Input Describing Functions and Nonlinear System Design} (1968); Slotine and Li, \emph{Applied Nonlinear Control} (1991), Chapter~5.
\end{itemize}

\paragraph{Reproducibility.} From \path{sim/gimbal_multidof}:
\begin{itemize}
\item \path{analisi_lqg_punto_lavoro.py}: results in \path{results_lqg_tuning};
\item \path{analisi_maxwell_tempo.py}: results in \path{results_lqg_tuning/time_domain_maxwell};
\item \path{analisi_lqg_passo1_flessibili.py}: results in \path{results_lqg_tuning/step1_flexible};
\item \path{analisi_tesi_puntamento.py}, then again with \texttt{--figure}: results in \path{results_lqg_tuning/no_bearings/thesis_pointing};
\item \path{report_lqg_tuning/build_lqg_tuning_report.py}: this report.
\end{itemize}

\end{document}
